% Consolidation des acquis de la méthodologie de la dissertation philosophique — Philosophie Tle Tech — Cours signé OnBuch+
% Programme officiel MINESEC. Compiler avec : tectonic N14-consolidation-acquis-methodologie-dissertation-philosop.tex
\newcommand{\DOCMATIERE}{Philosophie}
\newcommand{\DOCNIVEAU}{Tle Tech}
\newcommand{\DOCMODULE}{Philosophie – classe de Terminale technique}
\newcommand{\DOCLECON}{N14}
\newcommand{\DOCTITRE}{Consolidation des acquis de la méthodologie de la dissertation philosophique}
% OnBuch+ — préambule « cours signés » ENRICHI (compilé avec tectonic / XeLaTeX)
% Système visuel « Pop » d'OnBuch+ : fond crème (jamais blanc), bordures encre
% épaisses, ombres « dures » décalées, titres Archivo Black en capitales.
% Version MATHÉMATIQUES : graphes (pgfplots/tikz), boîtes pédagogiques
% (définition, propriété, méthode, attention, le savais-tu, exercice résolu,
% exercice, TP, bilan), page de garde et signature OnBuch+ ancrée en bas.
%
% Macros attendues dans main.tex AVANT \input{preamble} :
%   \DOCMATIERE  \DOCNIVEAU  \DOCMODULE  \DOCLECON (numéro)  \DOCTITRE
\documentclass[11pt]{article}

\usepackage{fontspec}
\usepackage[french]{babel}
\usepackage[a4paper,margin=2cm,top=3.4cm,bottom=2.8cm,headheight=1.4cm,footskip=1.3cm]{geometry}
\usepackage{amsmath,amssymb,amsfonts}
\usepackage{mathtools} % \xmapsto, \coloneqq... souvent utilisés par les modèles
\usepackage{cancel} % simplifications barrées (bilans de forces)
% \abs{z} (module, valeur absolue) et \norm{u} (norme), utilisés par les modèles
\providecommand{\abs}{}\let\abs\relax\DeclarePairedDelimiter{\abs}{\lvert}{\rvert}
\providecommand{\norm}{}\let\norm\relax\DeclarePairedDelimiter{\norm}{\lVert}{\rVert}
\usepackage[table]{xcolor} % [table] : fournit aussi \rowcolors (alternance de lignes)
\usepackage{pagecolor}
\usepackage{tikz}
\usetikzlibrary{arrows.meta,positioning,calc,shapes.geometric,shapes.misc,shapes.arrows,decorations.pathmorphing,decorations.pathreplacing,decorations.markings,patterns,shadows,mindmap,trees,fit,backgrounds,matrix,intersections,angles,quotes}
\usepackage{pgfplots}
\usepackage[version=4]{mhchem} % équations nucléaires \ce{^{235}_{92}U}
\usepackage[european,siunitx]{circuitikz} % schémas électriques
\pgfplotsset{compat=1.18}
\usepgfplotslibrary{fillbetween,groupplots} % aires entre courbes (calcul intégral)
\usepackage{enumitem}
\usepackage{fancyhdr}
% [most] charge toutes les bibliothèques tcolorbox (listings, minted, raster,
% documentation...) et gonfle inutilement la mémoire TeX sur un document long
% (~300 boîtes) : on ne charge que ce qui est réellement utilisé.
\usepackage[skins,breakable]{tcolorbox}
\usepackage{graphicx}
\usepackage{colortbl}
\usepackage{booktabs}
\usepackage{tabularx}
\usepackage{multirow}
\usepackage{diagbox} % case d'en-tête diagonale des tableaux à double entrée
% Colonnes extensibles centrée / gauche / droite (C, L, R), souvent utilisées par les modèles
\newcolumntype{C}{>{\centering\arraybackslash}X}
\newcolumntype{L}{>{\raggedright\arraybackslash}X}
\newcolumntype{R}{>{\raggedleft\arraybackslash}X}
\usepackage{array}
\usepackage{multicol}
\usepackage{siunitx}
% parse-numbers=false : accepte aussi \SI{2,5 \times 10^{-3}}{...}
\sisetup{locale=FR,output-decimal-marker={,},group-minimum-digits=4,parse-numbers=false}
\DeclareSIUnit\ohm{\ensuremath{\symup{\Omega}}} % U+2126 absent de la police
\DeclareSIUnit\division{div} % réglages d'oscilloscope (ms/div, V/div)
\DeclareSIUnit\tr{tr} % tours (vitesse de rotation, ex. \SI{1500}{\tr\per\minute})
\DeclareSIUnit\molar{M} % concentration molaire (ex. \SI{3}{\molar}, \SI{1}{\micro\molar})
\DeclareSIUnit\an{an} % année (le modèle confond parfois avec \year, un compteur TeX) — ex. \SI{5}{\kilo\meter\per\an}
\DeclareSIUnit\FCFA{FCFA} % franc CFA (ex. \SI{500}{\FCFA\per\kilo\gram})
\DeclareSIUnit\degreeBrix{\textdegree Brix} % teneur en sucre d'un sirop/jus (réfractométrie)
\DeclareSIUnit\month{mois} % le modèle utilise parfois \month, un compteur TeX, comme unité de durée
\DeclareSIUnit\microliter{\micro\liter} % le modèle écrit parfois \microliter en un seul token
\DeclareSIUnit\million{millions} % dénombrement (ex. \SI{15}{\million\per\mL} spermatozoïdes)
\usepackage{qrcode}
% \degree : commande fréquemment utilisée par les modèles pour un angle ou une
% température en dehors de \SI{}{} (non fournie par les paquets chargés).
\providecommand{\degree}{^{\circ}}
% \qed (fin de démonstration), utilisé par les modèles sans amsthm
\providecommand{\qed}{\hfill$\square$}
% \barre : double barre des valeurs interdites dans un tableau de variations (tkz-tab)
\providecommand{\barre}{\ensuremath{\|}}
% environnement proof (amsthm non chargé) utilisé par les modèles
\ifdefined\proof\else\newenvironment{proof}[1][Démonstration]{\par\noindent\textit{#1.}\ }{\qed\par}\fi
\usepackage{lastpage}
\usepackage{hyperref}
\hypersetup{hidelinks}

% ============================================================
% Palette « Pop » OnBuch+ — mêmes hex que la classe P (plus_ui.dart)
% ============================================================
\definecolor{poporange}{HTML}{F59321}
\definecolor{popdark}{HTML}{E07A0C}
\definecolor{popcream}{HTML}{FFF6E8}
\definecolor{popink}{HTML}{1C1714}
\definecolor{poppurple}{HTML}{7A5AE0}
\definecolor{popgreen}{HTML}{2E9E6A}
\definecolor{poppink}{HTML}{E0457B}
\definecolor{popblue}{HTML}{2F7DE1}
\definecolor{popgold}{HTML}{C98A12}
\definecolor{popmuted}{HTML}{8A7F76}
\definecolor{popline}{HTML}{EADFD2}
\definecolor{popgray}{HTML}{8A7F76}
\colorlet{popgrey}{popgray}
% Alias tolérants (le modèle nomme parfois les couleurs différemment)
\colorlet{popwhite}{white}
\colorlet{popblack}{popink}
\colorlet{popred}{poppink}
\colorlet{popyellow}{popgold}
% Teintes claires (fonds de tableaux/encadrés) — le modèle nomme parfois
% les couleurs avec un suffixe L (ex. popblueL) sans les avoir définies.
\colorlet{poporangeL}{poporange!20}
\colorlet{popdarkL}{popdark!20}
\colorlet{poppurpleL}{poppurple!20}
\colorlet{popgreenL}{popgreen!20}
\colorlet{poppinkL}{poppink!20}
\colorlet{popblueL}{popblue!20}
\colorlet{popgoldL}{popgold!20}
\colorlet{popredL}{popred!20}
\colorlet{popgrayL}{popgray!20}
% Teintes claires pour fonds de tableaux / zones
\colorlet{popblueL}{popblue!10!white}
\colorlet{poporangeL}{poporange!14!white}
\colorlet{popgreenL}{popgreen!12!white}
\colorlet{poppurpleL}{poppurple!12!white}
\colorlet{poppinkL}{poppink!12!white}
\colorlet{popgoldL}{popgold!14!white}

\pagecolor{popcream}

% ============================================================
% Typographie
% ============================================================
\defaultfontfeatures{Ligatures=TeX}
\setmainfont{PlusJakartaSans-Medium.ttf}[Path=../../../fonts/,BoldFont=PlusJakartaSans-Bold.ttf]
% Maths en sans-serif (Fira Math) avec chiffres et lettres droites de Plus
% Jakarta, pour que les formules se fondent dans le texte.
\usepackage[math-style=ISO,bold-style=ISO]{unicode-math}
\setmathfont{FiraMath-Regular.otf}[Path=../../../fonts/]
\setmathfont{PlusJakartaSans-Medium.ttf}[Path=../../../fonts/,range={up/num,up/{Latin,latin},\mathup}]
\usepackage{icomma} % virgule décimale sans espace en mode math
% Caractères mathématiques Unicode absents des polices de texte (Plus Jakarta,
% Archivo Black) : rendus par leur équivalent mathématique, en texte comme en
% formule — sinon ils s'affichent en carré vide (ex. « Arithmétique dans ℤ »).
\usepackage{newunicodechar}
\newunicodechar{ℝ}{\ensuremath{\mathbb{R}}}
\newunicodechar{ℤ}{\ensuremath{\mathbb{Z}}}
\newunicodechar{ℂ}{\ensuremath{\mathbb{C}}}
\newunicodechar{ℕ}{\ensuremath{\mathbb{N}}}
\newunicodechar{ℚ}{\ensuremath{\mathbb{Q}}}
\newunicodechar{𝔻}{\ensuremath{\mathbb{D}}}
\newunicodechar{α}{\ensuremath{\alpha}}
\newunicodechar{β}{\ensuremath{\beta}}
\newunicodechar{γ}{\ensuremath{\gamma}}
\newunicodechar{δ}{\ensuremath{\delta}}
\newunicodechar{ε}{\ensuremath{\varepsilon}}
\newunicodechar{θ}{\ensuremath{\theta}}
\newunicodechar{λ}{\ensuremath{\lambda}}
\newunicodechar{μ}{\ensuremath{\mu}}
\newunicodechar{σ}{\ensuremath{\sigma}}
\newunicodechar{τ}{\ensuremath{\tau}}
\newunicodechar{φ}{\ensuremath{\varphi}}
\newunicodechar{ψ}{\ensuremath{\psi}}
\newunicodechar{ω}{\ensuremath{\omega}}
\newunicodechar{Σ}{\ensuremath{\Sigma}}
% majuscules produites par \MakeUppercase dans les titres (α-aminés -> Α)
\newunicodechar{Α}{\ensuremath{\alpha}}
\newunicodechar{Β}{\ensuremath{\beta}}
\newunicodechar{Γ}{\ensuremath{\Gamma}}
\newunicodechar{Δ}{\ensuremath{\Delta}}
\newunicodechar{Ε}{\ensuremath{\varepsilon}}
\newunicodechar{Θ}{\ensuremath{\Theta}}
\newunicodechar{Λ}{\ensuremath{\Lambda}}
\newunicodechar{Μ}{\ensuremath{\mu}}
\newunicodechar{Τ}{\ensuremath{\tau}}
\newunicodechar{Φ}{\ensuremath{\Phi}}
\newunicodechar{Ψ}{\ensuremath{\Psi}}
\newunicodechar{Ω}{\ensuremath{\Omega}}
\newunicodechar{↦}{\ensuremath{\mapsto}}
\newunicodechar{∈}{\ensuremath{\in}}
\newunicodechar{∉}{\ensuremath{\notin}}
\newunicodechar{∧}{\ensuremath{\wedge}}
\newunicodechar{⇔}{\ensuremath{\Leftrightarrow}}
\newunicodechar{⟺}{\ensuremath{\Longleftrightarrow}}
\newunicodechar{⇒}{\ensuremath{\Rightarrow}}
\newunicodechar{⟹}{\ensuremath{\Longrightarrow}}
\newunicodechar{∘}{\ensuremath{\circ}}
\newunicodechar{∪}{\ensuremath{\cup}}
\newunicodechar{∩}{\ensuremath{\cap}}
\newunicodechar{≡}{\ensuremath{\equiv}}
\newunicodechar{⊂}{\ensuremath{\subset}}
\newunicodechar{⊥}{\ensuremath{\perp}}
\newunicodechar{∃}{\ensuremath{\exists}}
\newunicodechar{∀}{\ensuremath{\forall}}
\newunicodechar{∛}{\ensuremath{\sqrt[3]{\,}}}
\newunicodechar{∜}{\ensuremath{\sqrt[4]{\,}}}
\newunicodechar{⃗}{\ensuremath{{}^{\to}}}
\newunicodechar{ⁿ}{\ifmmode^{n}\else\textsuperscript{\ensuremath{n}}\fi}
\newunicodechar{ˣ}{\ifmmode^{x}\else\textsuperscript{\ensuremath{x}}\fi}
\newunicodechar{ᵇ}{\ifmmode^{b}\else\textsuperscript{\ensuremath{b}}\fi}
\newunicodechar{ʳ}{\ifmmode^{r}\else\textsuperscript{\ensuremath{r}}\fi}
\newunicodechar{ᵘ}{\ifmmode^{u}\else\textsuperscript{\ensuremath{u}}\fi}
\newunicodechar{ʸ}{\ifmmode^{y}\else\textsuperscript{\ensuremath{y}}\fi}
\newunicodechar{ᵃ}{\ifmmode^{a}\else\textsuperscript{\ensuremath{a}}\fi}
\newunicodechar{ᵉ}{\ifmmode^{e}\else\textsuperscript{\ensuremath{e}}\fi}
\newunicodechar{ᵏ}{\ifmmode^{k}\else\textsuperscript{\ensuremath{k}}\fi}
\newunicodechar{ᵖ}{\ifmmode^{p}\else\textsuperscript{\ensuremath{p}}\fi}
\newunicodechar{ᴺ}{\ifmmode^{N}\else\textsuperscript{\ensuremath{N}}\fi}
\newunicodechar{ᶜ}{\ifmmode^{c}\else\textsuperscript{\ensuremath{c}}\fi}
\newunicodechar{ʰ}{\ifmmode^{h}\else\textsuperscript{\ensuremath{h}}\fi}
\newunicodechar{ᵠ}{\ifmmode^{q}\else\textsuperscript{\ensuremath{q}}\fi}
\newunicodechar{ₙ}{\ifmmode_{n}\else\textsubscript{\ensuremath{n}}\fi}
\newunicodechar{ₐ}{\ifmmode_{a}\else\textsubscript{\ensuremath{a}}\fi}
\newunicodechar{ᵢ}{\ifmmode_{i}\else\textsubscript{\ensuremath{i}}\fi}
\newunicodechar{ᵣ}{\ifmmode_{r}\else\textsubscript{\ensuremath{r}}\fi}
\newunicodechar{ₖ}{\ifmmode_{k}\else\textsubscript{\ensuremath{k}}\fi}
\newunicodechar{ᵦ}{\ifmmode_{\beta}\else\textsubscript{\ensuremath{\beta}}\fi}
\newunicodechar{ₓ}{\ifmmode_{x}\else\textsubscript{\ensuremath{x}}\fi}
\newfontfamily\popdisplay{ArchivoBlack-Regular.ttf}[Path=../../../fonts/]
\newfontfamily\poplogo{SpaceGrotesk-Bold.ttf}[Path=../../../fonts/]
\newfontfamily\popsemi{PlusJakartaSans-SemiBold.ttf}[Path=../../../fonts/]
\newfontfamily\popxbold{PlusJakartaSans-ExtraBold.ttf}[Path=../../../fonts/]
\newfontfamily\popxboldls{PlusJakartaSans-ExtraBold.ttf}[Path=../../../fonts/,LetterSpace=3.0]

\setlist[enumerate]{leftmargin=1.7em,itemsep=5pt,topsep=4pt}
\setlist[itemize]{leftmargin=1.7em,itemsep=4pt,topsep=4pt,label={\color{poporange}\textbullet}}
\setlength{\parindent}{0pt}
\setlength{\parskip}{5pt}
\renewcommand{\arraystretch}{1.25}

% pgfplots : style Pop par défaut
\pgfplotsset{
  every axis/.append style={axis line style={popink,line width=1pt},tick style={popink},
    grid style={popline,line width=0.5pt},label style={font=\small},tick label style={font=\footnotesize}},
  popcurve/.style={poporange,line width=1.8pt,no markers,smooth},
}
% Même style disponible dans un \draw TikZ simple (hors pgfplots)
\tikzset{popcurve/.style={poporange,line width=1.8pt}}
\tikzset{popgrid/.style={popline,very thin}}
\colorlet{popcurve}{poporange} % le nom du style est parfois utilisé comme couleur

% ============================================================
% Pill / squiggle
% ============================================================
\newcommand{\poppill}[3]{%
  \tikz[baseline=-0.55ex]\node[draw=popink,line width=1.2pt,rounded corners=2.6pt,%
    fill=#2,inner xsep=6.5pt,inner ysep=3pt]{%
    \popxboldls\fontsize{7}{7}\selectfont\color{#3}\MakeUppercase{#1}};%
}
\newcommand{\popsquiggle}[1]{%
  \tikz[baseline=-0.4ex]{
    \draw[#1,line width=2.6pt,line cap=round]
      plot [smooth] coordinates {(0,0.09) (0.34,0.01) (0.68,0.21) (1.02,0.01) (1.36,0.10)};
  }%
}
% Mot-clé surligné (marqueur orange sous le texte)
\newcommand{\cle}[1]{{\popxbold\color{popdark}#1}}

% ============================================================
% En-tête / pied
% ============================================================
\pagestyle{fancy}
\fancyhf{}
\renewcommand{\headrulewidth}{0pt}
\renewcommand{\footrulewidth}{0pt}
\lhead{\raisebox{-0.15ex}{\poplogo\fontsize{13}{13}\selectfont\color{popink}OnBuch\color{poporange}+}}
\rhead{\poppill{\DOCMATIERE\ \textbullet\ \DOCNIVEAU\ \textbullet\ Leçon \DOCLECON}{white}{popink}}
\lfoot{\popxbold\fontsize{7.6}{7.6}\selectfont\color{popmuted}\MakeUppercase{Cours signé OnBuch+}\ \textbullet\ {\popsemi\DOCTITRE}}
\rfoot{\tikz[baseline=-0.6ex]\node[rounded corners=8.5pt,fill=popink,minimum height=17pt,minimum width=17pt,inner xsep=5pt,inner sep=0]{\popxbold\fontsize{8}{8}\selectfont\color{popcream}\thepage\,/\,\pageref*{LastPage}};}

% ============================================================
% Titres
% ============================================================
\newcommand{\chaptitre}[2]{%
  \begin{tcolorbox}[enhanced,colback=poporange,colframe=popink,boxrule=2.6pt,arc=11pt,
      drop shadow=popink,left=15pt,right=15pt,top=12pt,bottom=11pt,before skip=3pt,after skip=18pt]
    \poppill{#2}{popink}{popcream}\par\vspace{9pt}
    {\raggedright\hyphenpenalty=10000\exhyphenpenalty=10000\popdisplay\fontsize{22}{24}\selectfont\color{popcream}\MakeUppercase{#1}\par}
  \end{tcolorbox}
}
% Section numérotée : pastille orange avec le numéro + titre Archivo
\newcounter{popsec}
\newcommand{\coursec}[1]{%
  \stepcounter{popsec}\setcounter{popsub}{0}%
  \par\vspace{16pt}\noindent
  \begin{minipage}{\linewidth}
  \tikz[baseline=-0.7ex]\node[draw=popink,line width=1.6pt,fill=poporange,rounded corners=4pt,minimum size=22pt,inner sep=0]{\popdisplay\fontsize{12}{12}\selectfont\color{popcream}\thepopsec};%
  \hspace{8pt}{\popdisplay\fontsize{15}{17}\selectfont\color{popink}\MakeUppercase{#1}}\par
  \vspace{4pt}\noindent{\color{poporange}\rule{36pt}{3.6pt}}
  \end{minipage}\par\nopagebreak\vspace{8pt}
}
\newcounter{popsub}
\newcommand{\courssub}[1]{%
  \stepcounter{popsub}%
  \par\vspace{9pt}\noindent{\popxbold\fontsize{11.5}{13}\selectfont\color{popink}{\color{poporange}\thepopsec.\thepopsub}\hspace{6pt}#1}\par\nopagebreak\vspace{4pt}
}

% ============================================================
% Boîtes pédagogiques — même recette : fond blanc, bordure épaisse colorée,
% ombre dure encre, badge plein. Titre optionnel [#1].
% ============================================================
\tcbset{popbase/.style={breakable,enhanced,colback=white,boxrule=2.3pt,arc=9pt,
  shadow={3pt}{-3pt}{0pt}{popink},left=14pt,right=14pt,top=11pt,bottom=11pt,before skip=11pt,after skip=15pt,
  fontupper=\popsemi\fontsize{10.6}{14.5}\selectfont\color{popink},
  fontlower=\popsemi\fontsize{10.6}{14.5}\selectfont\color{popink}}}
% Coche dessinée (le glyphe ✓ n'existe pas dans les polices du document)
\newcommand{\popcheck}[1]{\tikz[baseline=-0.5ex]\draw[#1,line width=1.6pt,line cap=round,line join=round](-0.13,0.0)--(-0.04,-0.09)--(0.13,0.1);}
\providecommand{\coche}{\popcheck{popgreen}}
\providecommand{\resultat}[1]{\par\smallskip\noindent\fcolorbox{popgreen}{popgreenL}{\parbox{\dimexpr\linewidth-2\fboxsep-2\fboxrule\relax}{\textbf{Résultat.} #1}}\par\smallskip}
\newcommand{\popboxhead}[3]{\poppill{#1}{#2}{white}\ifx\relax#3\relax\else\hspace{6pt}{\popxbold\fontsize{10.6}{12}\selectfont\color{popink}#3}\fi\par\vspace{7pt}}

\newtcolorbox{aretenir}[1][]{popbase,colframe=popgreen,before upper={\popboxhead{À retenir}{popgreen}{#1}}}
\newtcolorbox{exemplebox}[1][]{popbase,colframe=poporange,before upper={\popboxhead{Exemple}{poporange}{#1}}}
\newtcolorbox{definition}[1][]{popbase,colframe=popblue,before upper={\popboxhead{Définition}{popblue}{#1}}}
\newtcolorbox{propriete}[1][]{popbase,colframe=poppurple,before upper={\popboxhead{Propriété}{poppurple}{#1}}}
\newtcolorbox{methode}[1][]{popbase,colframe=popgold,before upper={\popboxhead{Méthode}{popgold}{#1}}}
\newtcolorbox{attention}[1][]{popbase,colframe=poppink,before upper={\popboxhead{Attention}{poppink}{#1}}}
\newtcolorbox{experience}[1][]{popbase,colframe=popblue,colback=popblueL,before upper={\popboxhead{Expérience}{popblue}{#1}}}
% « Le savais-tu ? » : carte violette pleine (contexte, culture, Cameroun)
\newtcolorbox{savaistu}[1][]{popbase,colback=poppurple,colframe=popink,
  fontupper=\popsemi\fontsize{10.4}{14.2}\selectfont\color{white},
  before upper={\poppill{Le savais-tu\,?}{white}{poppurple}\ifx\relax#1\relax\else\hspace{6pt}{\popxbold\color{white}#1}\fi\par\vspace{7pt}}}
% Exercice résolu : énoncé en haut, \tcblower puis la solution sur fond crème
\newcounter{popexo}\newcounter{popexr}
\newtcolorbox{exoresolu}[1][]{popbase,colframe=poporange,colbacklower=popcream,
  segmentation style={popink,line width=1.2pt,dashed},
  before upper={\stepcounter{popexr}\popboxhead{Exercice résolu \thepopexr}{poporange}{#1}},
  before lower={\poppill{Solution}{popink}{popcream}\par\vspace{6pt}}}
% Exercice (non résolu en place) — la correction est dans la partie Corrigés
\newtcolorbox{exercice}[1][]{popbase,colframe=popink,
  before upper={\stepcounter{popexo}\popboxhead{Exercice \thepopexo}{popink}{#1}}}
% Corrigé
\newtcolorbox{corrige}[1][]{popbase,colframe=popgreen,colback=popgreenL,
  before upper={\popboxhead{Corrigé}{popgreen}{#1}}}
% Objectifs / prérequis / bilan
\newtcolorbox{objectifs}{popbase,colframe=popink,colback=poporangeL,
  before upper={\popboxhead{Objectifs de la leçon}{popink}{}}}
\newtcolorbox{prerequis}{popbase,colframe=popink,colback=popblueL,
  before upper={\popboxhead{Prérequis}{popblue}{}}}
\newtcolorbox{bilan}{popbase,colframe=popink,colback=white,boxrule=2.8pt,
  before upper={\popboxhead{Fiche bilan}{poporange}{L'essentiel à connaître pour l'examen}}}
% Figure encadrée avec légende
\newtcolorbox{popfigure}[1][]{enhanced,breakable=false,colback=white,colframe=popline,boxrule=1.4pt,arc=8pt,
  left=10pt,right=10pt,top=10pt,bottom=8pt,before skip=10pt,after skip=12pt,halign=center,
  lower separated=false,halign lower=center,
  fontlower=\popsemi\fontsize{9}{11.5}\selectfont\color{popmuted},
  #1}
\newcommand{\legende}[1]{\par\vspace{6pt}{\popsemi\fontsize{9}{11.5}\selectfont\color{popmuted}#1}}

% ============================================================
% Signature OnBuch+
% ============================================================
\newcommand{\poplogoinline}[1]{{\poplogo\fontsize{#1}{#1}\selectfont\color{popink}OnBuch\color{poporange}+}}
\newcommand{\signatureonbuch}{%
  \begin{tcolorbox}[enhanced,colback=popink,colframe=popink,boxrule=2.2pt,arc=10pt,
      drop shadow=poporange,left=14pt,right=14pt,top=10pt,bottom=10pt,before skip=0pt,after skip=0pt]
    \tikz[baseline=-0.7ex]\node[circle,fill=popgreen,draw=popcream,line width=1.3pt,minimum size=22pt,inner sep=0]{\popcheck{white}};%
    \hspace{9pt}\begin{minipage}[c]{0.62\linewidth}
      {\popxbold\fontsize{12}{13}\selectfont\color{popcream}Signé\ \textbullet\ {\poplogo OnBuch\color{poporange}+}}\par\vspace{2pt}
      {\raggedright\popsemi\fontsize{8}{10.5}\selectfont\color{popline}\MakeUppercase{Équipe pédagogique OnBuch+}\\\MakeUppercase{Conforme au programme officiel MINESEC}\par}
    \end{minipage}\hfill
    {\poplogo\fontsize{22}{22}\selectfont\color{popcream}OnBuch\color{poporange}+}
  \end{tcolorbox}%
}

% ============================================================
% Page de garde
% #1 titre · #2 matière · #3 niveau · #4 module · #5 n° leçon · #6 durée ·
% #7 description / objectifs (texte ou itemize)
% ============================================================
\newcommand{\pagegarde}[7]{%
  \thispagestyle{empty}
  \newgeometry{a4paper,margin=1.7cm,top=1.6cm,bottom=1.4cm}%
  \begin{tikzpicture}[remember picture,overlay]
    \fill[poporange!18] ([xshift=-3.2cm,yshift=-3.6cm]current page.north east) circle (4.6cm);
    \fill[poppurple!14] ([xshift=2.4cm,yshift=7.6cm]current page.south west) circle (3.8cm);
  \end{tikzpicture}%
  {\poplogo\fontsize{30}{30}\selectfont\color{popink}OnBuch\color{poporange}+}\hfill
  \raisebox{0.6ex}{\poppill{Cours signé \textbullet\ Premium}{popink}{popcream}}\par
  \vspace{1pt}\popsquiggle{poppurple}\par
  \vspace{8pt}
  \begin{tcolorbox}[enhanced,colback=poporange,colframe=popink,boxrule=2.8pt,arc=13pt,
      drop shadow=popink,left=18pt,right=18pt,top=12pt,bottom=13pt,after skip=10pt]
    \poppill{#2 \textbullet\ #3}{popink}{popcream}\hspace{5pt}\poppill{#4}{white}{popink}\par\vspace{8pt}
    {\popdisplay\fontsize{14}{14}\selectfont\color{popink}LEÇON #5}\par\vspace{4pt}
    {\raggedright\hyphenpenalty=10000\exhyphenpenalty=10000\popdisplay\fontsize{25}{27}\selectfont\color{popcream}\MakeUppercase{#1}\par}
    \vspace{8pt}
    {\popxbold\fontsize{9.5}{9.5}\selectfont\color{popink}\MakeUppercase{Durée conseillée : #6}}
  \end{tcolorbox}
  \begin{tcolorbox}[enhanced,colback=white,colframe=popink,boxrule=2.2pt,arc=10pt,
      drop shadow=popink,left=16pt,right=16pt,top=10pt,bottom=10pt,after skip=8pt]
    \poppill{Dans cette leçon}{poppurple}{white}\par\vspace{5pt}
    \popsemi\fontsize{9.3}{12.6}\selectfont\color{popink}#7
  \end{tcolorbox}
  \noindent\begin{minipage}[t]{0.487\linewidth}
    \begin{tcolorbox}[enhanced,colback=popgreen,colframe=popink,boxrule=2.2pt,arc=10pt,
        drop shadow=popink,left=11pt,right=11pt,top=10pt,bottom=10pt,height=3.1cm,valign=center]
      \tcbox[colback=white,colframe=popink,boxrule=1.3pt,arc=5pt,boxsep=0pt,nobeforeafter,
        left=3pt,right=3pt,top=3pt,bottom=3pt,valign=center]{\qrcode[height=1.9cm]{https://play.google.com/store/apps/details?id=com.luvvix.onbuch&hl=fr}}%
      \hspace{8pt}\begin{minipage}[c]{0.5\linewidth}
        {\popdisplay\fontsize{11}{12.5}\selectfont\color{white}\MakeUppercase{Télécharge OnBuch}}\par\vspace{4pt}
        {\raggedright\popsemi\fontsize{8.4}{11}\selectfont\color{white}Gratuit \textbullet\ Google Play\\Tuteur IA, annales, concours, résultats\par}
      \end{minipage}
    \end{tcolorbox}
  \end{minipage}\hfill
  \begin{minipage}[t]{0.487\linewidth}
    \begin{tcolorbox}[enhanced,colback=poppurple,colframe=popink,boxrule=2.2pt,arc=10pt,
        drop shadow=popink,left=11pt,right=11pt,top=10pt,bottom=10pt,height=3.1cm,valign=center]
      \tikz[baseline=-0.7ex]\node[circle,fill=white,draw=popink,line width=1.3pt,minimum size=1.6cm,inner sep=0]{\popdisplay\fontsize{24}{24}\selectfont\color{poppurple}+};%
      \hspace{8pt}\begin{minipage}[c]{0.6\linewidth}
        {\popdisplay\fontsize{11}{12.5}\selectfont\color{white}\MakeUppercase{Passe à OnBuch+}}\par\vspace{4pt}
        {\raggedright\popsemi\fontsize{8.4}{11}\selectfont\color{white}Tous les cours signés, Tuteur IA illimité, fiches PDF — onglet OnBuch+ de l'app\par}
      \end{minipage}
    \end{tcolorbox}
  \end{minipage}
  \vspace{8pt}
  \popcontenu
  \vfill
  \signatureonbuch
  \restoregeometry
  \newpage
}

% Rangée « Ce cours contient » (page de garde)
\newcommand{\poptile}[3]{% #1 couleur #2 gros texte #3 légende
  \begin{tcolorbox}[enhanced,colback=#1,colframe=popink,boxrule=2pt,arc=9pt,shadow={2.5pt}{-2.5pt}{0pt}{popink},
      left=6pt,right=6pt,top=9pt,bottom=9pt,halign=center,valign=center,height=1.75cm,width=\linewidth,nobeforeafter]
    {\popdisplay\fontsize{12.5}{13}\selectfont\color{white}\MakeUppercase{#2}}\par\vspace{3pt}
    {\centering\popsemi\fontsize{7.8}{9.5}\selectfont\color{white}#3\par}
  \end{tcolorbox}}
\newcommand{\popcontenu}{%
  \par\noindent{\popxboldls\fontsize{7.5}{7.5}\selectfont\color{popmuted}CE COURS SIGNÉ CONTIENT}\par\vspace{7pt}
  \noindent\begin{minipage}[t]{0.235\linewidth}\poptile{popblue}{Cours}{complet, illustré, pas à pas}\end{minipage}\hfill
  \begin{minipage}[t]{0.235\linewidth}\poptile{poporange}{Méthodes}{et exercices résolus}\end{minipage}\hfill
  \begin{minipage}[t]{0.235\linewidth}\poptile{poppink}{Exercices}{type examen + corrigés}\end{minipage}\hfill
  \begin{minipage}[t]{0.235\linewidth}\poptile{popgreen}{Fiche bilan}{l'essentiel à retenir}\end{minipage}\par}

% Sommaire
\newtcolorbox{sommaire}{enhanced,colback=white,colframe=popink,boxrule=2.2pt,arc=10pt,
  drop shadow=popink,left=18pt,right=18pt,top=13pt,bottom=13pt,before skip=6pt,after skip=20pt,
  before upper={\poppill{Sommaire}{poppurple}{white}\par\vspace{9pt}\popsemi\fontsize{10.6}{14.5}\selectfont\color{popink}}}

% Signature de fin de cours — poussée tout en bas de la dernière page
\newcommand{\findecours}{%
  \par\vfill\nopagebreak
  \begin{center}\popsquiggle{poporange}\end{center}\vspace{4pt}
  \signatureonbuch
}
\AtBeginDocument{\def\llbracket{\mathopen{[\mkern-3.5mu[}}\def\rrbracket{\mathclose{]\mkern-3.5mu]}}} % intervalles d'entiers (glyphe ⟦ absent de la police)
% Glyphes absents de Fira Math : redéfinis à partir de glyphes disponibles
\AtBeginDocument{%
  \def\setminus{\mathbin{\backslash}}\let\smallsetminus\setminus
  \def\vdots{\mathord{\vbox{\baselineskip3.2pt\lineskiplimit0pt\kern4pt\hbox{.}\hbox{.}\hbox{.}}}}%
  \def\ddots{\mathinner{\mkern1mu\raise6.5pt\hbox{.}\mkern2mu\raise3.6pt\hbox{.}\mkern2mu\raise0.7pt\hbox{.}\mkern1mu}}%
  \def\triangle{\mathord{\scalebox{0.9}{$\symup{\Delta}$}}}%
  \def\nabla{\mathord{\raisebox{\depth}{\scalebox{1}[-1]{$\symup{\Delta}$}}}}%
  \def\checkmark{\popcheck{popdark}}%
  \def\star{\mathbin{\ast}}\def\bigstar{\mathord{\ast}}%
  \def\diamond{\mathbin{\scalebox{0.6}{$\Diamond$}}}%
  \def\bigcup{\mathop{\mathchoice{\raisebox{-0.3em}{\scalebox{1.7}{$\cup$}}}{\scalebox{1.25}{$\cup$}}{\cup}{\cup}}\displaylimits}%
  \def\bigcap{\mathop{\mathchoice{\raisebox{-0.3em}{\scalebox{1.7}{$\cap$}}}{\scalebox{1.25}{$\cap$}}{\cap}{\cap}}\displaylimits}%
  \def\lesseqqgtr{\mathrel{\lessgtr}}\def\bigoplus{\mathop{\oplus}\displaylimits}%
}
\providecommand{\ssi}{si et seulement si} % abréviation fréquente
\AtBeginDocument{\providecommand{\wideparen}{\overparen}} % arc de cercle

%% FIGFIX-BEGIN v3 — correctifs globaux des figures (docs/onbuch-generation/tools/figfix)
\usepackage{adjustbox}\usepackage{etoolbox}
% toute figure TikZ/pgfplots plus large que la ligne est réduite à la largeur disponible
\BeforeBeginEnvironment{tikzpicture}{\begin{adjustbox}{max width=\linewidth}}
\AfterEndEnvironment{tikzpicture}{\end{adjustbox}}
% légendes pgfplots lisibles par-dessus les courbes
\pgfplotsset{every axis/.append style={legend style={fill=white,fill opacity=0.92,text opacity=1,draw=black!15,inner sep=3pt}}}
% étiquettes d'arêtes (syntaxe edge["..."]) avec fond blanc
\tikzset{every edge quotes/.append style={fill=white,inner sep=1.2pt}}
% bulles de cartes mentales : texte à la ligne dans la bulle, bulle un peu plus grande
\tikzset{every concept/.append style={text width=2.0cm,align=center,minimum size=2.3cm,inner sep=2pt}}
%% FIGFIX-END
\begin{document}

\pagegarde{Consolidation des acquis de la méthodologie de la dissertation philosophique}{Philosophie}{Tle Tech}{Philosophie – classe de Terminale technique}{N14}{2 séances (fiche de progression harmonisée)}{Cette leçon de consolidation rassemble et systématise l'ensemble des exigences méthodologiques de la dissertation philosophique : introduction, développement dialectique (thèse, antithèse, synthèse), conclusion et usage des citations. À travers des travaux pratiques guidés et des sujets types du Baccalauréat technique camerounais, l'élève s'entraîne à rédiger et justifier une démarche complète applicable à des situations concrètes de la vie courante.
\vspace{4pt}
\begin{multicols}{2}\small
\begin{itemize}
\item À la fin de cette leçon, je sais rappeler les principes et les exigences méthodologiques de la dissertation philosophique
\item rédiger une introduction complète (accroche, analyse du sujet, problématique, annonce du plan)
\item construire un développement dialectique structuré en thèse, antithèse et synthèse
\item formuler une conclusion qui répond à la problématique et ouvre sur une question nouvelle
\item choisir, introduire et commenter correctement une citation philosophique
\item analyser un sujet de dissertation pour dégager ses termes clés et son problème
\end{itemize}\end{multicols}}

\begin{sommaire}
\begin{multicols}{2}
\begin{enumerate}
\item Rappel des principes et exigences de la dissertation philosophique
\item TP 1 : L'introduction
\item TP 2 : Le développement — thèse, antithèse, synthèse
\item TP 3 : La conclusion
\item TP 4 : Choix et utilisation des citations
\item Application à la vie courante et entraînement à l'examen
\item Activité d'intégration
\item Méthodes et exercices résolus
\item Exercices
\item Corrigés des exercices
\item Fiche bilan
\end{enumerate}
\end{multicols}
\end{sommaire}

\begin{objectifs}
\begin{itemize}
\item À la fin de cette leçon, je sais rappeler les principes et les exigences méthodologiques de la dissertation philosophique
\item rédiger une introduction complète (accroche, analyse du sujet, problématique, annonce du plan)
\item construire un développement dialectique structuré en thèse, antithèse et synthèse
\item formuler une conclusion qui répond à la problématique et ouvre sur une question nouvelle
\item choisir, introduire et commenter correctement une citation philosophique
\item analyser un sujet de dissertation pour dégager ses termes clés et son problème
\item appliquer la méthode de la dissertation à des situations concrètes de la vie courante camerounaise
\item rédiger et justifier une démarche argumentative complète en temps limité, comme à l'examen
\end{itemize}
\end{objectifs}

\begin{prerequis}
\begin{itemize}
\item Notion d'argumentation et de raisonnement (Première)
\item Les grandes étapes de la dissertation vues en cours d'année de Terminale
\item Connaissance des auteurs et doctrines étudiés au programme de Terminale
\item Techniques de lecture et d'analyse d'un texte philosophique
\end{itemize}
\end{prerequis}

\newpage

\coursec{Rappel des principes et exigences de la dissertation philosophique}

Au lycée de Nkongsamba, Amina, élève de Terminale technique, obtient 6/20 à sa dissertation philosophique alors qu'elle maîtrisait parfaitement le contenu du cours : elle connaissait Platon, Descartes et Kant sur le bout des doigts. Pourtant, son introduction manquait de problématique et son développement alignait des citations sans véritable argumentation. Son professeur lui fait alors remarquer qu'au Baccalauréat technique, la dissertation représente souvent l'épreuve décisive de philosophie, celle qui fait la différence entre l'admis et le recalé. Comme Amina, beaucoup d'élèves camerounais échouent non par ignorance des doctrines, mais par méconnaissance de la \cle{méthode}. Comment transformer ses connaissances philosophiques en une copie structurée, argumentée et convaincante ? C'est toute la problématique de cette leçon de consolidation, qui reprend pas à pas les principes et les exigences de la dissertation philosophique.

\courssub{Qu'est-ce qu'une dissertation philosophique ?}

\begin{definition}[La dissertation philosophique]
La \cle{dissertation philosophique} est une réflexion écrite, \textbf{personnelle} et \textbf{méthodique}, qui traite un \textbf{problème} en mobilisant des \textbf{arguments} (raisons logiques) et des \textbf{références} (auteurs, doctrines, exemples). Elle ne se contente pas d'exposer des connaissances : elle discute une question, confronte des thèses et aboutit à une réponse justifiée.
\end{definition}

Commençons par l'intuition. Quand un ami te demande : « Crois-tu que l'argent fait le bonheur ? », tu ne récites pas un dictionnaire : tu réfléchis, tu donnes ton avis, tu le justifies, tu envisages l'avis contraire, puis tu conclus. La dissertation philosophique fait exactement cela, mais de manière écrite, rigoureuse et organisée. C'est une \emph{conversation raisonnée} avec le sujet.

Trois traits distinguent donc la dissertation de tout autre exercice scolaire :
\begin{itemize}
\item Elle est \textbf{personnelle} : l'élève pense par lui-même. Les auteurs cités ne remplacent pas sa réflexion, ils la nourrissent.
\item Elle est \textbf{structurée} : introduction, développement en parties ordonnées, conclusion. Rien n'est laissé au hasard.
\item Elle est \textbf{argumentée} : chaque idée est soutenue par une raison, illustrée par un exemple, éventuellement appuyée sur une référence philosophique.
\end{itemize}

\begin{attention}[Le hors-sujet, faute capitale]
L'erreur la plus fréquente et la plus sanctionnée consiste à \textbf{réciter le cours sans répondre à la question posée}. Un élève qui, devant le sujet « La tradition est-elle un obstacle au progrès ? », récite tout son cours sur la tradition (définitions, auteurs, exemples) sans jamais se demander si elle \emph{empêche} ou non le progrès, fait du hors-sujet. Le correcteur du Baccalauréat technique sanctionne sévèrement le hors-sujet, \textbf{même s'il est bien écrit}. Une copie brillante mais à côté de la question ne peut dépasser la moyenne. Règle d'or : chaque paragraphe doit contenir au moins un mot du sujet ou un synonyme direct.
\end{attention}

\courssub{Les trois exigences fondamentales de la copie}

Toute dissertation de qualité repose sur trois exigences que le correcteur évalue explicitement. On peut les retenir grâce aux trois « C » de la copie : \textbf{C}onformité (exactitude), \textbf{C}ohérence, \textbf{C}larité.

\textbf{1. L'exactitude : la fidélité au sujet.} La copie doit traiter \emph{exactement} la question posée, ni plus, ni moins. Cela suppose de respecter les termes du sujet sans les déformer, de ne pas transformer une question sur la \emph{liberté} en question sur le \emph{devoir}, et de ne pas élargir le sujet à tout le chapitre. L'exactitude se vérifie à chaque étape : ma problématique reprend-elle les termes du sujet ? Ma conclusion répond-elle à la question initiale ?

\textbf{2. La cohérence : la logique interne.} La cohérence exige que les parties de la copie s'enchaînent logiquement et que l'élève ne se contredise pas. Si la première partie affirme que la tradition est un obstacle et la conclusion qu'elle n'a jamais été un obstacle, la copie est incohérente. La cohérence suppose aussi des transitions claires entre les parties : chaque partie doit naître de la précédente par un lien logique explicite (« Cependant... », « Mais cette réponse est insuffisante car... »).

\textbf{3. La clarté : l'expression et l'orthographe.} Une pensée profonde mal exprimée est une pensée perdue. La clarté concerne la langue (orthographe, grammaire, conjugaison), le style (phrases construites, vocabulaire précis, pas de langage SMS) et la présentation (écriture lisible, alinéas, parties visibles). Au Baccalauréat technique, les fautes graves et répétées font perdre des points précieux.

\begin{aretenir}[Les trois exigences en une phrase]
Une bonne dissertation philosophique est une copie \textbf{exacte} (elle répond à la question), \textbf{cohérente} (elle se tient logiquement) et \textbf{claire} (elle se lit sans effort). Toute la méthode qui suit sert ces trois exigences.
\end{aretenir}

\courssub{L'analyse du sujet : la première victoire de l'élève}

Tout se joue dans les trente premières minutes. Un sujet de dissertation n'est pas une invitation à parler d'un thème : c'est une \emph{question précise} qu'il faut d'abord décoder. L'analyse du sujet consiste à passer du \textbf{sujet brut} (la phrase donnée par l'examinateur) à la \textbf{problématique} (le problème que cette phrase soulève).

\begin{methode}[Analyser un sujet en quatre étapes]
\begin{enumerate}
\item \textbf{Souligner les termes clés} : repérer les notions philosophiques et les mots qui portent la question (verbes, adjectifs, oppositions).
\item \textbf{Définir chaque terme} : donner au brouillon une définition précise de chaque notion, en distinguant le sens courant du sens philosophique.
\item \textbf{Repérer la tension ou la question implicite} : identifier l'opposition, le doute ou le paradoxe que le sujet met en scène (c'est là que naît le problème).
\item \textbf{Formuler la problématique} : rédiger une ou deux questions qui expriment le problème et qui guideront tout le développement.
\end{enumerate}
\end{methode}

\begin{exemplebox}[Analyse du sujet : « La tradition est-elle un obstacle au progrès ? »]
\textbf{Étape 1 — Termes clés :} \emph{tradition}, \emph{obstacle}, \emph{progrès}.

\textbf{Étape 2 — Définitions :} la \emph{tradition} est l'ensemble des coutumes, croyances et pratiques transmises de génération en génération (au Cameroun : rites, langues, autorité coutumière) ; un \emph{obstacle} est ce qui empêche ou freine un mouvement ; le \emph{progrès} est l'amélioration des conditions de vie et des connaissances (technique, scientifique, sociale).

\textbf{Étape 3 — Tension :} le sujet oppose la \textbf{fidélité au passé} (conserver) à l'\textbf{innovation} (transformer). La tradition protège-t-elle une identité précieuse ou enchaîne-t-elle la société à des pratiques dépassées ?

\textbf{Étape 4 — Problématique :} « La transmission du passé est-elle incompatible avec le mouvement vers l'avenir, ou la tradition peut-elle au contraire orienter et féconder le progrès ? »
\end{exemplebox}

\begin{popfigure}
\begin{center}
\begin{tikzpicture}[font=\small]
\fill[popblue!25] (0,4.6) -- (9,4.6) -- (7.4,3.4) -- (1.6,3.4) -- cycle;
\fill[poporange!30] (1.6,3.4) -- (7.4,3.4) -- (6.2,2.2) -- (2.8,2.2) -- cycle;
\fill[popgreen!30] (2.8,2.2) -- (6.2,2.2) -- (5.4,1.0) -- (3.6,1.0) -- cycle;
\fill[poppink!40] (3.6,1.0) -- (5.4,1.0) -- (4.9,0) -- (4.1,0) -- cycle;
\node[align=center] at (4.5,4.0) {\textbf{1. Sujet brut} : termes clés soulignés};
\node[align=center] at (4.5,2.8) {\textbf{2. Définitions} de chaque notion};
\node[align=center] at (4.5,1.6) {\textbf{3. Tension} repérée};
\node[align=center] at (4.5,0.5) {\textbf{4. Problématique}};
\draw[->,thick,popdark] (4.5,-0.3) -- (4.5,-1.0);
\node[align=center,popdark] at (4.5,-1.6) {\textbf{Le problème à discuter}\\ guide tout le devoir};
\end{tikzpicture}
\end{center}
\legende{Schéma en entonnoir : le passage du sujet brut à la problématique. Chaque étape rétrécit et précise la réflexion : on part de la phrase de l'examinateur pour aboutir au problème unique qui commandera toute la dissertation.}
\end{popfigure}

\courssub{Les types de sujets au Baccalauréat technique camerounais}

À l'examen, les sujets de dissertation se présentent sous trois formes principales. Savoir les reconnaître permet d'adapter immédiatement sa stratégie.

\begin{center}
\begin{tabularx}{\linewidth}{|l|X|X|}
\hline
\rowcolor{popblueL} \textbf{Type de sujet} & \textbf{Exemple} & \textbf{Stratégie attendue} \\
\hline
\textbf{Question directe} & « La conscience fait-elle la grandeur de l'homme ? » & Répondre par oui, par non, puis nuancer : construire thèse, antithèse et synthèse. \\
\hline
\textbf{Affirmation à discuter} & « Toute vérité est relative. Discutez. » & Examiner la thèse proposée : la justifier, la contester, puis l'évaluer. \\
\hline
\textbf{Comparaison de deux thèses} & « La paix vient-elle de la force ou du droit ? » & Confronter les deux options l'une à l'autre, puis dépasser l'alternative. \\
\hline
\end{tabularx}
\end{center}

Quel que soit le type, l'exigence reste identique : le sujet pose un \emph{problème} et attend une \emph{discussion}, jamais une simple récitation.

\courssub{Dissertation ou exposé : ne pas confondre les exercices}

Beaucoup d'élèves, comme Amina, rédigent en réalité un \textbf{exposé} alors que l'examinateur attend une \textbf{dissertation}. La différence est décisive : l'exposé \emph{informe}, la dissertation \emph{problématise}.

\begin{popfigure}
\begin{center}
\begin{tikzpicture}[font=\small]
\node[draw,thick,popblue,fill=popblue!15,rounded corners,minimum width=6.4cm,minimum height=4.6cm,align=left,anchor=north west] at (0,0)
{\textbf{LA DISSERTATION}\\[2pt]
\textbf{Objectif :} discuter un problème\\
\textbf{Structure :} thèse -- antithèse --\\
\hspace{0.6cm}synthèse, parties liées par des\\
\hspace{0.6cm}transitions\\
\textbf{Posture :} l'élève pense,\\
\hspace{0.6cm}argumente, prend position\\
\textbf{Question centrale :} « qu'en\\
\hspace{0.6cm}penser et pourquoi ? »};
\node[draw,thick,poporange,fill=poporange!15,rounded corners,minimum width=6.4cm,minimum height=4.6cm,align=left,anchor=north west] at (7,0)
{\textbf{L'EXPOSÉ}\\[2pt]
\textbf{Objectif :} informer, présenter\\
\textbf{Structure :} catalogue de\\
\hspace{0.6cm}connaissances juxtaposées\\
\textbf{Posture :} l'élève récite le cours\\
\hspace{0.6cm}et les auteurs\\
\textbf{Question centrale :} « que sait-on\\
\hspace{0.6cm}sur ce thème ? »};
\draw[<->,very thick,popdark] (6.55,-2.3) -- (6.95,-2.3);
\end{tikzpicture}
\end{center}
\legende{Tableau comparatif : dissertation et exposé. La dissertation problématise et engage la pensée personnelle de l'élève ; l'exposé se contente d'informer en juxtaposant des connaissances.}
\end{popfigure}

\begin{exoresolu}[Reconnaître une dissertation et un exposé]
\textbf{Énoncé.} Sujet : « L'État est-il nécessaire à la liberté ? » Voici deux débuts de copie. Lequel relève de la dissertation, lequel de l'exposé ? Justifiez.

\emph{Copie A :} « L'État est une institution politique. Hobbes a écrit le \emph{Léviathan} en 1651. Rousseau a écrit le \emph{Contrat social} en 1762. Il existe plusieurs formes d'État : la monarchie, la république\ldots »

\emph{Copie B :} « Obéir à l'État, n'est-ce pas renoncer à sa liberté ? Pourtant, sans lois communes, la liberté de chacun serait menacée par la force du plus puissant. Faut-il alors voir dans l'État l'ennemi de la liberté ou sa condition ? »

\tcblower
\textbf{Solution.} La \textbf{copie A} est un exposé : elle aligne des informations (dates, titres, classifications) sans poser de problème ni répondre à la question posée. La \textbf{copie B} est un début de dissertation : elle transforme le sujet en tension (obéissance vs liberté), montre l'enjeu (la sécurité de tous) et formule une problématique. C'est elle qui sera valorisée au Baccalauréat technique.
\end{exoresolu}

\courssub{La gestion du temps à l'examen}

L'épreuve de philosophie au Baccalauréat technique dure quatre heures. Une mauvaise répartition du temps ruine les meilleures connaissances : l'élève qui rédige trop tôt fonce vers le hors-sujet ; celui qui analyse trop longtemps n'a plus le temps d'écrire. La répartition recommandée est la suivante :

\begin{center}
\begin{tabularx}{\linewidth}{|l|X|X|}
\hline
\rowcolor{popgreenL} \textbf{Phase} & \textbf{Durée} & \textbf{Travail à accomplir} \\
\hline
\textbf{Analyse du sujet et plan} & 30 min & Définir les termes, dégager la tension, formuler la problématique, construire le plan détaillé au brouillon. \\
\hline
\textbf{Rédaction} & 2 h 30 & Rédiger l'introduction, le développement (thèse, antithèse, synthèse) et la conclusion, au propre, avec transitions. \\
\hline
\textbf{Relecture} & 30 min & Corriger l'orthographe, vérifier la cohérence, contrôler que chaque partie répond au sujet. \\
\hline
\end{tabularx}
\end{center}

\begin{attention}[Les pièges de la gestion du temps]
\begin{itemize}
\item Ne \textbf{jamais} commencer à rédiger au propre avant d'avoir un plan complet au brouillon.
\item Ne pas sacrifier la relecture : une copie pleine de fautes perd des points sur la clarté, même si le fond est bon.
\item Surveiller la montre : si la synthèse n'est pas rédigée à 15 minutes de la fin, écrire au moins la conclusion, qui est obligatoire.
\end{itemize}
\end{attention}

\begin{savaistu}[La notation de la copie au Baccalauréat technique du Cameroun]
Au Baccalauréat technique du Cameroun, la copie de philosophie est notée sur 20 selon des critères explicites : la \textbf{pertinence} (la copie répond-elle au sujet ?), la \textbf{cohérence} (le raisonnement est-il logique et structuré ?), la \textbf{richesse des références} (auteurs, exemples bien choisis et bien utilisés) et la \textbf{qualité de la langue} (orthographe, expression, présentation). Comprendre ces critères, c'est déjà préparer sa réussite : chaque exigence méthodologique étudiée dans cette leçon correspond à des points concrets sur la copie.
\end{savaistu}

\begin{exercice}[Entraînement à l'analyse de sujet]
Pour chacun des sujets suivants, souligne les termes clés, définis-les, dégage la tension et formule une problématique :
\begin{enumerate}
\item « Le travail est-il une aliénation ? »
\item « Peut-on dire que la technique libère l'homme ? Discutez. »
\item « Le bonheur dépend-il de nous ou des circonstances ? »
\end{enumerate}
Précise pour chaque sujet s'il s'agit d'une question directe, d'une affirmation à discuter ou d'une comparaison de deux thèses.
\end{exercice}

\begin{corrige}[Exercice : Entraînement à l'analyse de sujet]
\textbf{Sujet 1 (question directe).} Termes clés : \emph{travail} (activité humaine visant à transformer la nature et à produire) et \emph{aliénation} (perte de soi, dépossession du produit ou du sens de son activité). Tension : le travail réalise l'homme (Hegel : il humanise) mais peut aussi l'asservir (Marx : l'ouvrier est dépossédé). Problématique : « Le travail est-il le chemin de l'accomplissement humain ou la source de sa dépossession ? »

\textbf{Sujet 2 (affirmation à discuter).} Termes clés : \emph{technique} (ensemble des outils et procédés inventés par l'homme) et \emph{libérer} (rendre autonome, affranchir d'une contrainte). Tension : la technique soulage des contraintes naturelles (machines, médecine) mais crée de nouvelles dépendances (écrans, obsolescence). Problématique : « La technique affranchit-elle l'homme de la nature ou l'enchaîne-t-elle à de nouvelles servitudes ? »

\textbf{Sujet 3 (comparaison de deux thèses).} Termes clés : \emph{bonheur} (état de satisfaction durable) et \emph{circonstances} (conditions extérieures indépendantes de notre volonté). Tension : le stoïcisme soutient que le bonheur dépend de notre jugement (Épictète), tandis que le sens commun l'attribue à la fortune. Problématique : « Sommes-nous les artisans de notre bonheur ou les jouets des événements ? »
\end{corrige}

En résumé, cette première étape de consolidation a rappelé l'essentiel : la dissertation philosophique est une réflexion personnelle, structurée et argumentée sur un problème ; elle obéit aux trois exigences d'exactitude, de cohérence et de clarté ; elle commence par une analyse rigoureuse du sujet en quatre étapes ; elle se distingue radicalement de l'exposé ; et elle exige une gestion méthodique du temps. Fort de ces principes, nous pouvons maintenant passer aux travaux pratiques, en commençant par la construction de l'introduction.

\coursec{TP 1 : L'introduction}

Après avoir rappelé les principes généraux qui gouvernent la dissertation philosophique — rigueur conceptuelle, articulation dialectique, exigence d'argumentation —, nous entrons maintenant dans l'atelier concret. L'introduction est la \cle{porte d'entrée} de votre copie : c'est le moment où vous prouvez à l'examinateur que vous avez compris le sujet, que vous en avez identifié la tension interne et que vous possédez une boussole pour ne pas vous égarer. Une introduction ratée, c'est souvent une dissertation qui dérive ; une introduction maîtrisée, c'est la promesse d'un développement cohérent. Nous allons disséquer ses quatre étapes, les illustrer par un exemple ancré dans notre réalité camerounaise, puis rédiger collectivement une introduction modèle.

\courssub{La fonction de l'introduction et ses quatre étapes}

L'introduction n'est pas un préambule décoratif ni un résumé de ce qui va suivre. Elle remplit trois fonctions indissociables : \textbf{présenter le sujet} en le situant dans un horizon de sens (l'accroche), \textbf{dégager le problème} philosophique qui s'y cache (l'analyse des termes et la problématique), et \textbf{annoncer la démarche} qui va permettre de le résoudre (l'annonce du plan). Ces quatre étapes forment un \cle{entonnoir} : on part du large (le sens commun, l'évidence) pour aboutir au précis (la question unique qui commande la copie).

\begin{popfigure}
\begin{tikzpicture}[
    node distance=1.2cm and 1.5cm,
    stepnode/.style={rectangle, rounded corners, draw=popblue, thick, fill=popblueL, text width=3.2cm, align=center, font=\small\bfseries, minimum height=1.8cm},
    arrownode/.style={-Stealth, thick, popdark},
    labelnode/.style={font=\footnotesize\textit, text width=3.5cm, align=center, popmuted}
]
% Nodes
\node[stepnode, align=center] (acc){1. L'ACCROCHE\\ (Amorce)};
\node[labelnode, below=0.1cm of acc, align=center] (acclab){Sens commun, actualité,\\ situation concrète};
\node[stepnode, right=of acc, align=center] (ana){2. L'ANALYSE\\ DES TERMES};
\node[labelnode, below=0.1cm of ana, align=center] (analab){Définition, distinctions,\\ oppositions, étymologie};
\node[stepnode, right=of ana] (prob) {3. LA PROBLÉMATIQUE};
\node[labelnode, below=0.1cm of prob, align=center] (problab){Question centrale,\\ tension, paradoxe};
\node[stepnode, right=of prob, align=center] (plan){4. L'ANNONCE\\ DU PLAN};
\node[labelnode, below=0.1cm of plan, align=center] (planlab){Trois moments logiques,\\ rédigés, liés};

% Arrows
\draw[arrownode] (acc.east) -- (ana.west) node[midway, above, font=\small\bfseries, poporange] {$\rightarrow$};
\draw[arrownode] (ana.east) -- (prob.west) node[midway, above, font=\small\bfseries, poporange] {$\rightarrow$};
\draw[arrownode] (prob.east) -- (plan.west) node[midway, above, font=\small\bfseries, poporange] {$\rightarrow$};

% Staircase visual effect (background)
\begin{scope}[on background layer]
\fill[popblue!10] (acc.south west) rectangle (plan.north east);
\draw[popblue, dashed, thick] (acc.south west) -- (plan.south east);
\draw[popblue, dashed, thick] (acc.north west) -- (plan.north east);
\end{scope}
\end{tikzpicture}
\legende{Schéma en escalier des quatre étapes de l'introduction : chaque marche prépare logiquement la suivante. L'accroche ouvre l'horizon, l'analyse en dégage les concepts, la problématique en extrait la tension, le plan annonce la résolution.}
\end{popfigure}

\begin{methode}[Les quatre marches de l'introduction — Structure type]
\textbf{Étape 1 — L'accroche (2--3 lignes).} Partez d'un fait observable, d'une évidence partagée, d'une actualité locale ou d'une citation courte. Montrez immédiatement que cette évidence cache une difficulté.
\textbf{Étape 2 — L'analyse des termes (3--4 lignes).} Définissez les concepts clés du sujet. Relevez les polysemies, les oppositions (ex. : coutume vs loi, obéissance vs respect), les nuances. C'est ici que se joue la pertinence de la suite.
\textbf{Étape 3 — La problématique (1--2 lignes).} Formulez \emph{une} question principale, claire, qui exprime le conflit entre des thèses plausibles. Elle doit être la fille directe de l'analyse.
\textbf{Étape 4 — L'annonce du plan (2--3 lignes).} Présentez la logique de votre démonstration en trois temps (thèse, antithèse, synthèse) \emph{sans} numéroter « I, II, III » ni écrire « Premièrement... ». Utilisez des connecteurs logiques : « Nous examinerons d'abord... avant de voir... pour enfin montrer... ».
\end{methode}

\courssub{L'accroche (ou amorce) : entrer dans le sujet par le concret}

L'erreur classique est l'accroche « catalogue » : « Depuis la nuit des temps, les hommes ont des coutumes... » Cela ne dit rien, ne concerne personne et ne prépare aucune analyse. Une bonne accroche est un \cle{piège à pensée} : elle présente une situation qui semble aller de soi, puis révèle une fissure.

\begin{methode}[Rédiger une accroche efficace — Démarche en 3 temps]
\emph{Temps 1 : L'ancrage.} Choisissez un point de départ concret : un fait de société au Cameroun (débat sur le « bride price », place des chefferies, réseaux sociaux vs traditions), une scène de vie (conseil de famille, héritage), une phrase d'auteur courte et percutante.
\emph{Temps 2 : L'évidence.} Énoncez ce que « tout le monde pense » ou ce que la situation montre de prime abord (ex. : « La coutume assure la cohésion du village »).
\emph{Temps 3 : Le basculement.} Introduisez le « mais » philosophique : « Pourtant, cette même coutume peut exclure les femmes, figer les générations, entrer en contradiction avec les droits humains. » C'est ce basculement qui appelle l'analyse des termes.
\end{methode}

\begin{attention}[Piège n°1 : L'accroche « Nuit des temps »]
\emph{Interdite :} « Depuis l'aube de l'humanité, la coutume règle la vie des sociétés. » \emph{Pourquoi ?} Trop vague, aucun ancrage spatio-temporel, ne distingue pas la coutume de la loi ou de l'habitude, ne crée pas de tension.
\emph{Préférez :} « À Bafoussam comme à Yaoundé, le règlement d'un litige foncier devant le chef de quartier selon la coutume ancestrale apaise les cœurs... jusqu'à ce qu'une fille réclame sa part d'héritage, exclue par la tradition. » (Ancré, concret, tension immédiate).
\end{attention}

\begin{exemplebox}[Accroches réussies pour « Faut-il obéir à la coutume ? »]
\textbf{Accroche A (Actualité/Société) :} « En 2023, le débat sur le code de la famille au Cameroun a ravivé la tension entre droit moderne et droit coutumier : l'âge du mariage, la polygamie, l'héritage des filles. La coutume, garante de l'ordre social, devient-elle un obstacle à la justice ? »
\textbf{Accroche B (Situation villageoise) :} « Dans le village de Nkolbisson, le conseil des notables décide de l'exclusion d'un jeune ayant refusé le mariage arrangé. L'obéissance à la coutume assure la paix du clan, mais au prix de la liberté d'un individu. »
\textbf{Accroche C (Citation brève + basculement) :} « "La coutume est le droit des peuples" (Montesquieu). Mais si ce droit fige les inégalités, l'obéissance n'est-elle pas complicité ? »
\end{exemplebox}

\courssub{L'analyse des termes du sujet : clarifier les concepts}

C'est l'étape la plus technique et la plus négligée. Elle ne consiste pas à donner des définitions de dictionnaire, mais à \cle{problématiser le vocabulaire}. Pour « Faut-il obéir à la coutume ? », il faut distinguer :
\begin{itemize}
    \item \textbf{Coutume} vs \textbf{Loi} (écrite, étatique, sanctionnée) vs \textbf{Habitude} (individuelle, non normative).
    \item \textbf{Obéissance} (soumission à une autorité, contrainte extérieure) vs \textbf{Respect} (adhésion intérieure, reconnaissance de valeur) vs \textbf{Conformité} (imitation sociale).
    \item \textbf{Faut-il} : obligation morale ? contrainte sociale ? condition de validité ? Le « il » est impersonnel : qui est le sujet de l'obligation ? L'individu ? Le citoyen ? Le membre de la communauté ?
\end{itemize}
L'analyse doit faire apparaître les \emph{oppositions internes} au sujet : la coutume est à la fois \emph{identité} (ce qui nous unit) et \emph{inertie} (ce qui nous fige) ; l'obéissance est à la fois \emph{devoir} (reconnaissance) et \emph{aliénation} (perte d'autonomie).

\courssub{La problématique : le cœur moteur de la dissertation}

\begin{definition}[Problématique]
La problématique est la \cle{question centrale}, unique et formulée, qui naît de l'analyse des termes et exprime la tension philosophique du sujet. Elle ne demande pas un savoir (type « Qu'est-ce que la coutume ? ») mais une \emph{prise de position raisonnée} (type « L'obéissance à la coutume est-elle compatible avec l'autonomie morale ? »). Toute la dissertation n'est que la réponse argumentée à cette question.
\end{definition}

Une bonne problématique possède trois caractères :
\begin{enumerate}
    \item \textbf{Elle est issue de l'analyse} : on y retrouve les concepts définis et leurs oppositions.
    \item \textbf{Elle est dialectique} : elle met en présence deux thèses défendables (pour/contre, ou plutôt : « dans quelle mesure / sous quelle condition »).
    \item \textbf{Elle est opératoire} : elle commande le plan. Si le plan ne répond pas \emph{précisément} à la question posée, la copie est hors sujet.
\end{enumerate}

\begin{exemplebox}[Problématique réussie — Analyse]
\textbf{Sujet :} « Faut-il obéir à la coutume ? »
\textbf{Problématique :} \emph{« Si la coutume fonde l'identité communautaire et assure l'ordre social, son obéissance aveugle ne menace-t-elle pas la liberté individuelle, la justice et le progrès social ? »}

\emph{Pourquoi ça marche ?}
\begin{itemize}
    \item \textbf{Thèse implicite (le "Si...") :} Reprend l'analyse : coutume = identité + ordre (valeur positive).
    \item \textbf{Tension (le "ne menace-t-elle pas...") :} Introduit l'antithèse : obéissance = aliénation + injustice + conservatisme (valeur négative).
    \item \textbf{Ouverture (le "progrès social") :} Annonce la synthèse : critère de validité de la coutume (justice, droits humains).
    \item \textbf{Commande du plan :} I. La coutume comme fondement nécessaire de l'ordre et de l'identité (Thèse). II. Les dangers de l'obéissance servile : injustice, blocage du progrès (Antithèse). III. Une obéissance critique : trier, réformer, soumettre la coutume à la raison universelle (Synthèse).
\end{itemize}
\end{exemplebox}

\courssub{L'annonce du plan : guider le lecteur sans numéroter}

L'annonce du plan n'est pas un sommaire administratif. C'est une \cle{phrase de transition rédigée} qui montre la \emph{logique} de votre pensée. Elle doit faire sentir la nécessité du passage d'une partie à l'autre.

\begin{attention}[Piège n°2 : L'annonce « liste de courses »]
\emph{Interdit :} « Premièrement, nous verrons que la coutume est nécessaire. Deuxièmement, nous verrons qu'elle est dangereuse. Troisièmement, nous verrons comment la réformer. »
\emph{Pourquoi ?} Cela coupe la pensée en tranches étanches. La dissertation n'est pas une addition de paragraphes, c'est un \emph{mouvement}.
\emph{Modèle correct :} « Nous analyserons d'abord en quoi la coutume, comme mémoire collective, fonde la cohésion du groupe et justifie une obéissance légitime ; nous interrogerons ensuite les dérives d'une soumission sans discernement qui peut consacrer l'injuste ; nous montrerons enfin que la véritable fidélité à la tradition exige de la soumettre à l'épreuve de la raison et du droit, pour ne conserver que ce qui élève l'humain. »
\end{attention}

\begin{methode}[Annonce du plan — Les connecteurs de la logique dialectique]
\begin{itemize}
    \item \textbf{Thèse (Moment 1) :} « Nous examinerons d'abord / Il convient d'analyser dans un premier temps / La première étape consiste à montrer... » $\rightarrow$ Valorisation du concept (fonction positive).
    \item \textbf{Antithèse (Moment 2) :} « Mais / Toutefois / Cependant, cette analyse révèle des limites / Il faut alors s'interroger sur les dangers de... / Une lecture critique met en lumière... » $\rightarrow$ Critique, renversement, complexification.
    \item \textbf{Synthèse (Moment 3) :} « C'est pourquoi / Au terme de cette dialectique / Il apparaît finalement que... / La solution réside dans... » $\rightarrow$ Dépassement, critère supérieur, réconciliation conditionnelle.
\end{itemize}
\end{methode}

\courssub{TP guidé : rédaction collective sur « Faut-il obéir à la coutume ? »}

Appliquons la méthode. Contexte : un village de la région du Centre, Nkolmekok. Le chef traditionnel, appuyé par le conseil des anciens, impose à une jeune fille, Amina, 17 ans, bachelière, un mariage avec un notable polygame de 60 ans, au nom de la « coutume des alliances ». Le père d'Amina, fonctionnaire à Yaoundé, s'y oppose au nom de la loi et du consentement. Le village se divise.

\begin{exoresolu}[Rédaction guidée de l'introduction]
\textbf{Énoncé :} Rédigez l'introduction complète (10--15 lignes) en signalant les 4 étapes.
\tcblower
\textbf{Solution détaillée :}

\begin{center}
\begin{tabularx}{\linewidth}{|>{\bfseries\color{popblue}}p{2.2cm}|X|}
\hline
\textbf{[ACCROCHE]} & Dans le village de Nkolmekok, le tambour appelle au palais : le chef et les notables scellent l'alliance d'Amina, 17 ans, bachelière, à un notable quadragénaire. La coutume des alliances, garante de la paix entre clans, s'impose à tous. Pourtant, le père d'Amina, instituteur à la retraite, brandit le Code civil et la Convention des droits de l'enfant : le consentement libre manque, l'âge légal n'est pas atteint. \\ \hline
\textbf{[ANALYSE]} & Le sujet oppose deux normes : la \emph{coutume}, ensemble de règles non écrites, transmises oralement, sanctionnées par la pression sociale et l'autorité des ancêtres, et la \emph{loi}, texte étatique, universel, coercitif. \emph{Obéir} implique une soumission à une autorité reconnue ; mais ici, l'autorité est double : celle des ancêtres (coutume) et celle de la République (loi). Le « faut-il » interroge la \emph{légitimité} de l'obéissance : est-elle morale (devoir), sociale (conformité) ou juridique (obligation) ? \\ \hline
\textbf{[PROBLÉMATIQUE]} & \emph{Si la coutume tisse le lien social et préserve l'identité collective, l'obéissance qu'elle exige peut-elle s'affranchir de l'examen critique de la raison et du droit, au risque de sacrifier la liberté et la dignité de l'individu ?} \\ \hline
\textbf{[ANNONCE DU PLAN]} & Nous montrerons d'abord que la coutume, comme fondement de l'ordre communautaire, légitime une obéissance qui n'est pas servilité mais reconnaissance. Nous verrons ensuite que l'obéissance aveugle devient aliénation lorsqu'elle fige des injustices et nie l'autonomie morale. Nous conclurons que la véritable fidélité à la tradition commande de la soumettre à l'épreuve de l'universel, pour ne lui obéir que lorsqu'elle élève l'humain. \\ \hline
\end{tabularx}
\end{center}

\textbf{Commentaire de la correction :} L'accroche est ancrée (lieu, noms, âge, conflit concret). L'analyse distingue coutume/loi, obéissance/contrainte, et ouvre le « faut-il ». La problématique reprend la tension « lien social vs liberté individuelle » et annonce le critère de la synthèse (« épreuve de l'universel »). L'annonce du plan est rédigée, fluide, utilise « d'abord / ensuite / conclurons que » et exprime la logique dialectique (fondement $\rightarrow$ limite $\rightarrow$ critère de dépassement). Longueur : environ 13 lignes. \textbf{Objectif atteint.}
\end{exoresolu}

\begin{popfigure}
\begin{tikzpicture}[
    node distance=0.8cm,
    box/.style={rectangle, draw=popdark, thick, fill=popcream, rounded corners, inner sep=6pt, text width=\linewidth-12pt},
    tag/.style={rectangle, fill=poporange, text=white, font=\tiny\bfseries, rounded corners=2pt, inner sep=2pt, anchor=west},
    tagblue/.style={tag, fill=popblue},
    taggreen/.style={tag, fill=popgreen},
    tagpurple/.style={tag, fill=poppurple},
]
\node[box] (intro) {
    \begin{minipage}{\linewidth}
    \textbf{[ACCROCHE]} Dans le village de Nkolmekok, le tambour appelle au palais : le chef et les notables scellent l'alliance d'Amina, 17 ans, bachelière, à un notable quadragénaire. La coutume des alliances, garante de la paix entre clans, s'impose à tous. Pourtant, le père d'Amina, instituteur à la retraite, brandit le Code civil et la Convention des droits de l'enfant : le consentement libre manque, l'âge légal n'est pas atteint.\par\vspace{2pt}
    \textbf{[ANALYSE]} Le sujet oppose deux normes : la \emph{coutume}, ensemble de règles non écrites, transmises oralement, sanctionnées par la pression sociale et l'autorité des ancêtres, et la \emph{loi}, texte étatique, universel, coercitif. \emph{Obéir} implique une soumission à une autorité reconnue ; mais ici, l'autorité est double : celle des ancêtres (coutume) et celle de la République (loi). Le « faut-il » interroge la \emph{légitimité} de l'obéissance : est-elle morale (devoir), sociale (conformité) ou juridique (obligation) ?\par\vspace{2pt}
    \textbf{[PROBLÉMATIQUE]} \emph{Si la coutume tisse le lien social et préserve l'identité collective, l'obéissance qu'elle exige peut-elle s'affranchir de l'examen critique de la raison et du droit, au risque de sacrifier la liberté et la dignité de l'individu ?}\par\vspace{2pt}
    \textbf{[ANNONCE DU PLAN]} Nous montrerons d'abord que la coutume, comme fondement de l'ordre communautaire, légitime une obéissance qui n'est pas servilité mais reconnaissance. Nous verrons ensuite que l'obéissance aveugle devient aliénation lorsqu'elle fige des injustices et nie l'autonomie morale. Nous conclurons que la véritable fidélité à la tradition commande de la soumettre à l'épreuve de l'universel, pour ne lui obéir que lorsqu'elle élève l'humain.
    \end{minipage}
};
% Marginal tags (simulated via overlay nodes on left)
\node[tagblue, left=0.2cm of intro.north west, yshift=-0.8cm] {ACCROCHE};
\node[taggreen, left=0.2cm of intro.west, yshift=-3.2cm] {ANALYSE};
\node[tagpurple, left=0.2cm of intro.west, yshift=-5.8cm] {PROBLÉMATIQUE};
\node[tag, left=0.2cm of intro.south west, yshift=0.8cm] {PLAN};
\end{tikzpicture}
\legende{Introduction type sur « Faut-il obéir à la coutume ? » avec repérage visuel des quatre étapes. Notez la progressivité : du fait villageois concret (accroche) à la question universelle (problématique) en passant par la clarification conceptuelle (analyse).}
\end{popfigure}

\courssub{Critères d'évaluation d'une introduction réussie}

Au Probatoire comme au Baccalauréat technique, l'introduction est notée sur 4 à 5 points (souvent 1 pt par étape + 1 pt cohérence). Voici la grille interne de l'examinateur.

\begin{aretenir}[Grille d'auto-évaluation — Introduction (Cible : 10--15 lignes)]
\begin{enumerate}
    \item \textbf{Accroche (1 pt) :} Ancrage concret (Cameroun, actualité, situation vécue) + Basculement vers le problème. \emph{Pas de "Nuit des temps".}
    \item \textbf{Analyse des termes (1 pt) :} Définitions pertinentes + Distinctions opératoires (coutume/loi, obéissance/respect) + Mise en évidence des oppositions internes au sujet.
    \item \textbf{Problématique (1 pt) :} \emph{Une} question unique, dialectique, issue de l'analyse, qui commande le plan. \emph{Pas de "Qu'est-ce que... ?" ni de liste de questions.}
    \item \textbf{Annonce du plan (1 pt) :} Trois moments logiques (Thèse/Antithèse/Synthèse) annoncés de manière \emph{rédigée}, avec connecteurs dialectiques (d'abord / mais / enfin). \emph{Pas de "I, II, III" ni "Premièrement".}
    \item \textbf{Cohérence globale (1 pt) :} Enchaînement fluide, pas de répétition, ton philosophique, longueur respectée (10--15 lignes manuscrites $\approx$ 120--180 mots).
\end{enumerate}
\end{aretenir}

\begin{savaistu}[Le saviez-vous ? — L'introduction au Bac Technique]
Au Probatoire/Bac Technique, le correcteur lit souvent \emph{d'abord} l'introduction et la conclusion pour se faire une idée du niveau. Une introduction qui « tient la route » (problématique nette, plan dialectique annoncé) place la copie d'emblée dans le « bon paquet » (note $\ge$ 12/20). À l'inverse, une introduction floue ou hors-sujet crée un préjugé négatif difficile à rattraper au développement. Soignez-la comme votre carte de visite.
\end{savaistu}

\courssub{Exercice d'entraînement : À vous de jouer}

\begin{exercice}[Rédaction d'introductions — Sujets types Bac Tech]
Rédigez une introduction complète (10--15 lignes, 4 étapes signalées) pour \textbf{deux} des sujets suivants. Utilisez une accroche ancrée au Cameroun.
\begin{enumerate}
    \item « La technique libère-t-elle l'homme ? » (Accroche possible : pénibilité agricole vs tracteur / téléphone portable vs isolement).
    \item « Faut-il tout dire ? » (Accroche possible : secret de famille / whistleblower / réseaux sociaux).
    \item « Le travail est-il une contrainte ? » (Accroche possible : moto-taxi à Douala / fonctionnaire vs entrepreneur / stage non rémunéré).
    \item « La science peut-elle nous dispenser de philosopher ? » (Accroche possible : diagnostic médical vs choix éthique / IA et décision de justice).
\end{enumerate}
\textit{Conseil :} Chronométrez-vous (15 min max par introduction). Relisez avec la grille d'auto-évaluation ci-dessus. Échangez avec un camarade pour correction croisée.
\end{exercice}

\begin{corrige}[Exercice n°1 — Pistes pour « La technique libère-t-elle l'homme ? »]
\textbf{Accroche :} À Mbalmayo, le cultivateur Pierre a troqué sa daba contre un motoculteur : la même surface se laboure en deux heures au lieu de deux jours. Son corps souffre moins, son rendement double. Pourtant, le bruit du moteur chasse le chant des oiseaux, la dette d'achat l'enchaîne à la banque, et son fils, parti à la ville, ne reviendra pas hériter du savoir-faire ancestral.
\textbf{Analyse :} \emph{Technique} : ensemble de moyens artificiels pour transformer la nature (outils, machines, algorithmes). \emph{Libérer} : affranchir de la contrainte (naturelle, physique, sociale), accroître l'autonomie. \emph{Homme} : être de besoins et de fins. Le sujet oppose \emph{libération par la technique} (maîtrise de la nature, gain de temps) et \emph{aliénation par la technique} (dépendance, déshumanisation, asservissement aux machines).
\textbf{Problématique :} \emph{Si la technique libère l'homme de la pénibilité naturelle et du manque, ne risque-t-elle pas, par sa logique propre d'efficacité et de rentabilité, de le soumettre à de nouvelles servitudes — techniques, économiques, existentielles ?}
\textbf{Annonce du plan :} Nous analyserons d'abord comment la technique, en augmentant nos puissances, nous affranchit de l'hostilité du milieu et du travail forcé. Nous interrogerons ensuite le paradoxe d'une émancipation qui engendre dépendance, standardisation et perte de sens. Nous montrerons enfin que la libération véritable suppose une \emph{maîtrise politique et éthique} de la technique, pour qu'elle reste un moyen au service de fins humaines choisies.
\end{corrige}

\coursec{TP 2 : Le développement — thèse, antithèse, synthèse}

Dans le TP 1, nous avons appris à construire une \cle{introduction} complète : accroche, analyse du sujet, problématique, annonce du plan. Mais une belle introduction ne suffit pas : elle ouvre une porte, encore faut-il marcher. C'est dans le \cle{développement} que se joue l'essentiel de la note au Probatoire et au Baccalauréat technique, car c'est là que le correcteur vérifie si l'élève sait réellement \emph{penser}, c'est-à-dire confronter des arguments, et non simplement réciter des opinions. Ce TP 2 est donc le cœur de notre consolidation méthodologique.

\courssub{La structure dialectique du développement}

\textbf{L'intuition d'abord.} Imaginez une palabre villageoise, sous le fromager, à Bafoussam ou à Maroua. Un conflit oppose deux familles. Le premier ancien prend la parole et défend une position. Puis un second ancien se lève et la conteste avec vigueur. Enfin, le chef, après avoir écouté les deux, prononce une parole qui dépasse le désaccord et réconcilie les points de vue. La dissertation philosophique fonctionne exactement comme cette palabre : elle met en scène un dialogue entre des positions contraires, puis elle conclut par une parole de sagesse qui dépasse l'opposition.

\begin{definition}[Le plan dialectique]
Le \cle{plan dialectique} (ou plan thèse--antithèse--synthèse) est une structure de développement en trois parties : la \cle{thèse} expose une première réponse au problème (souvent celle du sens commun ou la réponse la plus immédiate) ; l'\cle{antithèse} critique cette réponse et défend la position opposée ; la \cle{synthèse} dépasse l'opposition en apportant une réponse nuancée et justifiée à la problématique.
\end{definition}

\begin{definition}[La synthèse]
La \cle{synthèse} est le moment de la dissertation où la réflexion \emph{dépasse} l'opposition entre la thèse et l'antithèse. Elle ne consiste pas à dire « les deux ont raison », mais à \textbf{distinguer}, \textbf{nuancer} ou \textbf{reformuler le problème} de manière à conserver ce que chaque position avait de vrai, tout en écartant leurs erreurs respectives. C'est la réponse personnelle et argumentée de l'élève à la problématique.
\end{definition}

\begin{popfigure}
\begin{center}
\begin{tikzpicture}[scale=1.0]
\coordinate (T) at (-4,0);
\coordinate (A) at (4,0);
\coordinate (S) at (0,4.2);
\draw[line width=1.2pt, popmuted, dashed] (T) -- (A);
\draw[-{Stealth[length=3mm]}, line width=1.4pt, popblue] (T) -- (S);
\draw[-{Stealth[length=3mm]}, line width=1.4pt, poporange] (A) -- (S);
\node[draw=popblue, fill=popblueL, rounded corners=4pt, align=center, line width=1.2pt, minimum width=3.4cm] at (T) {\textbf{THÈSE}\\ \small Première réponse\\ \small (sens commun)};
\node[draw=poporange, fill=poporangeL, rounded corners=4pt, align=center, line width=1.2pt, minimum width=3.4cm] at (A) {\textbf{ANTITHÈSE}\\ \small Critique et\\ \small position opposée};
\node[draw=poppurple, fill=poppurpleL, rounded corners=4pt, align=center, line width=1.2pt, minimum width=3.6cm] at (S) {\textbf{SYNTHÈSE}\\ \small Dépassement nuancé\\ \small et justifié};
\node[popmuted, align=center, font=\small] at (0,-0.6) {opposition à examiner};
\node[popblue, font=\small\itshape, rotate=45] at (-2.6,2.3) {critique};
\node[poporange, font=\small\itshape, rotate=-45] at (2.6,2.3) {critique};
\end{tikzpicture}
\end{center}
\legende{Le triangle dialectique : la thèse et l'antithèse s'opposent à la base ; la synthèse, au sommet, dépasse leur conflit en intégrant la part de vérité de chacune.}
\end{popfigure}

\begin{attention}[La synthèse n'est pas un compromis mou]
L'erreur la plus fréquente consiste à écrire en synthèse : « la vérité est entre les deux » ou « tout dépend des points de vue ». C'est ce qu'on appelle le \emph{relativisme}, et le correcteur le sanctionne. Une vraie synthèse \textbf{tranche} : elle établit une distinction précise (par exemple : distinguer les coutumes vivantes des coutumes aliénantes) et elle \textbf{justifie} cette distinction par un argument. Elle répond à la problématique, elle ne l'esquive pas.
\end{attention}

\courssub{La règle d'or : un paragraphe = une idée}

Le développement compte en général \textbf{trois parties} (thèse, antithèse, synthèse), et chaque partie contient \textbf{deux à trois paragraphes}. Mais qu'est-ce qu'un bon paragraphe philosophique ?

\begin{methode}[Construire un paragraphe philosophique]
Un paragraphe de dissertation obéit à une structure rigoureuse en cinq temps :
\begin{enumerate}
\item \textbf{Annoncer l'idée} : une seule idée par paragraphe, formulée clairement dès la première phrase (la phrase d'idée).
\item \textbf{Développer l'argument} : expliquer \emph{pourquoi} cette idée est vraie, par un raisonnement logique et non par une simple affirmation.
\item \textbf{Illustrer par un exemple} : donner un cas concret, de préférence tiré de la réalité camerounaise ou de l'expérience commune.
\item \textbf{Appuyer par une référence} : citer brièvement un auteur, une œuvre ou un proverbe qui conforte l'analyse (voir TP 4 sur les citations).
\item \textbf{Conclure provisoirement et transiter} : ramener l'idée à la problématique et préparer le paragraphe suivant.
\end{enumerate}
Règle absolue : \textbf{une seule idée par paragraphe}. Si vous avez deux idées, faites deux paragraphes.
\end{methode}

\begin{popfigure}
\begin{center}
\begin{tikzpicture}[node distance=0.35cm,
box/.style={draw=popdark, rounded corners=3pt, align=center, minimum width=2.5cm, minimum height=1cm, line width=1pt, font=\small}]
\node[box, fill=popblueL, draw=popblue, align=center] (i){\textbf{1. Idée}\\ annoncée};
\node[box, fill=popgreenL, draw=popgreen, right=of i, align=center] (a){\textbf{2. Argument}\\ développé};
\node[box, fill=poporangeL, draw=poporange, right=of a, align=center] (e){\textbf{3. Exemple}\\ concret};
\node[box, fill=poppurpleL, draw=poppurple, right=of e, align=center] (r){\textbf{4. Référence}\\ auteur};
\node[box, fill=popgoldL, draw=popgold, right=of r, align=center] (t){\textbf{5. Transition}\\ vers la suite};
\foreach \x/\y in {i/a, a/e, e/r, r/t}
  \draw[-{Stealth[length=2.5mm]}, line width=1.2pt, popdark] (\x) -- (\y);
\end{tikzpicture}
\end{center}
\legende{La chaîne d'un paragraphe type : chaque maillon est obligatoire. Un paragraphe sans argument n'est qu'une opinion ; un paragraphe sans transition laisse le lecteur en panne.}
\end{popfigure}

\courssub{La thèse : exposer la première réponse avec force}

La \cle{thèse} présente la réponse la plus immédiate au problème, souvent celle du \emph{sens commun} ou de la tradition. Deux exigences :

\begin{itemize}
\item Il faut la défendre \textbf{sincèrement et fortement}, avec ses meilleurs arguments, même si vous n'êtes pas d'accord : le correcteur juge votre capacité à comprendre une position, pas vos opinions personnelles.
\item Il faut l'exposer \textbf{pour ensuite l'examiner}, non pour y rester : la thèse n'est qu'une étape de la réflexion.
\end{itemize}

\begin{exemplebox}[Exemple de thèse : « La coutume est-elle un obstacle au progrès ? »]
Premier paragraphe de thèse : la coutume assure la \textbf{cohésion sociale}. \emph{Idée} : en transmettant des règles communes, la coutume fait vivre le groupe. \emph{Argument} : sans normes partagées, aucune vie collective n'est possible ; la coutume est la mémoire qui unit les générations. \emph{Exemple} : dans les sociétés traditionnelles camerounaises, les rites d'initiation, les funérailles ou les règles du mariage coutumier créent des liens de solidarité entre les familles et les villages. \emph{Référence} : on peut rappeler que pour les sociologues, les coutumes sont des « institutions » qui structurent la vie sociale. \emph{Transition} : mais ce qui unit peut-il aussi enfermer ?
\end{exemplebox}

\courssub{L'antithèse : critiquer la thèse et défendre la position inverse}

L'\cle{antithèse} est le moment de la \textbf{critique}. Elle comporte deux gestes : d'abord montrer les \emph{limites} de la thèse (ce qu'elle ne voit pas, ce qu'elle prétend à tort), ensuite défendre la position \emph{opposée} avec la même rigueur.

\begin{exemplebox}[Exemple d'antithèse : la critique de la coutume]
\emph{Idée} : la coutume peut figer la société et opprimer l'individu. \emph{Argument} : une règle n'est pas juste parce qu'elle est ancienne ; certaines traditions consacrent l'inégalité et interdisent toute remise en cause. \emph{Exemple} : certaines pratiques coutumières (mariages forcés, lévirat imposé, exclusion des filles de l'héritage foncier dans certaines régions) sont aujourd'hui contestées au Cameroun car contraires aux droits humains et à la Constitution. \emph{Référence} : les défenseurs des droits de la personne rappellent que la dignité de l'individu prime sur la tradition. \emph{Transition} : faut-il alors rejeter toute coutume ?
\end{exemplebox}

\begin{attention}[Erreur fréquente : juxtaposer sans discuter]
Beaucoup de copies se contentent d'\emph{exposer} les opinions : « certains disent que... d'autres disent que... ». C'est insuffisant et c'est sanctionné. En philosophie, chaque thèse doit être \textbf{examinée critiquement} : il faut montrer ses arguments, mais aussi ses présupposés et ses limites. La dissertation n'est pas un catalogue d'opinions, c'est un \emph{procès} équitable où chaque position est interrogée. Posez-vous toujours la question : « cette thèse a-t-elle raison ? jusqu'à quel point ? que cache-t-elle ? »
\end{attention}

\courssub{La synthèse : dépasser l'opposition}

La synthèse est le moment le plus difficile et le plus valorisé. Trois stratégies possibles :

\begin{enumerate}
\item \textbf{Distinguer} : montrer que thèse et antithèse parlent de choses différentes. Exemple : distinguer les \emph{coutumes vivantes} (celles qui transmettent des valeurs de solidarité et peuvent évoluer) des \emph{coutumes aliénantes} (celles qui violent la dignité humaine). Les unes doivent être préservées, les autres réformées.
\item \textbf{Nuancer} : montrer que chaque position est vraie \emph{sous une condition} ou \emph{dans un domaine} précis.
\item \textbf{Reformuler le problème} : montrer que la vraie question était ailleurs. Exemple : le problème n'est pas « coutume ou modernité », mais « quel critère permet de juger une coutume ? » — critère qui peut être la dignité de la personne humaine.
\end{enumerate}

\courssub{Les transitions : faire progresser la réflexion}

Les \cle{transitions} sont les phrases de liaison entre les parties. Elles ont une double fonction : \textbf{bilan} de ce qui vient d'être établi, et \textbf{ouverture} d'une difficulté nouvelle qui justifie la partie suivante. Elles montrent au correcteur que la réflexion \emph{progresse} au lieu de tourner en rond.

\begin{exemplebox}[Une transition réussie]
« Mais si la coutume unit la communauté, ne risque-t-elle pas d'étouffer l'individu ? C'est ce que contestent les partisans de la modernité. »
\par\smallskip
Analysons : la première proposition fait le \emph{bilan} de la thèse (« la coutume unit ») ; la question introduit la \emph{difficulté} (« étouffer l'individu ») ; la dernière phrase annonce l'\emph{antithèse}. Trois mouvements en deux phrases : voilà le modèle à imiter.
\end{exemplebox}

\begin{savaistu}[Le plan dialectique n'est pas le seul accepté]
Contrairement à une idée reçue, le plan thèse--antithèse--synthèse n'est pas obligatoire au Baccalauréat technique. Un \textbf{plan thématique} (étude successive de plusieurs aspects du problème) ou un \textbf{plan analytique} (constat, causes, conséquences, solution) peut obtenir d'excellentes notes, à condition de rester \emph{problématisé} : chaque partie doit répondre à une question et faire progresser la réflexion. Le plan dialectique reste cependant le plus sûr pour les sujets qui opposent naturellement deux réponses (« faut-il... ? », « est-ce... ou... ? »).
\end{savaistu}

\courssub{TP guidé : « Le travail est-il une aliénation ? »}

Construisons collectivement le plan détaillé du développement de ce sujet classique.

\begin{exoresolu}[Plan détaillé : « Le travail est-il une aliénation ? »]
Énoncé : après avoir rédigé l'introduction (problématique : « le travail est-il une servitude qui dépossède l'homme de lui-même, ou au contraire la voie de son accomplissement ? »), construire le développement dialectique complet, avec pour chaque partie les idées, arguments, exemples et références.
\tcblower
\textbf{Solution détaillée.}
\par\smallskip
\textbf{I. Thèse : le travail est une aliénation.}
\par
\emph{Paragraphe 1.} Idée : le travailleur est dépossédé du produit de son travail. Argument : dans le salariat, l'ouvrier produit un objet qui ne lui appartient pas et qui se retourne contre lui comme une puissance étrangère ; plus il produit, plus il s'appauvrit. Référence : \textbf{Marx} (Manuscrits de 1844) parle du travailleur « aliéné » dans le produit, dans l'acte de production et dans son essence humaine. Exemple : l'ouvrier d'une usine de Douala qui assemble des produits qu'il ne pourra jamais acheter.
\par
\emph{Paragraphe 2.} Idée : le travail subi détruit le corps et l'esprit. Argument : la répétition mécanique des gestes, les horaires imposés, la fatigue font du travail une souffrance plutôt qu'un épanouissement. Exemple : les travailleurs des chantiers ou des plantations, épuisés par des tâches pénibles pour un salaire de survie. Transition : « Pourtant, si le travail n'était qu'une malédiction, comment expliquer que l'homme s'y réalise et que le chômage soit vécu comme une humiliation ? »
\par\smallskip
\textbf{II. Antithèse : le travail est libération et humanisation.}
\par
\emph{Paragraphe 1.} Idée : par le travail, l'homme transforme la nature et se transforme lui-même. Argument : en façonnant le monde, l'homme se reconnaît dans son œuvre ; le travail est éducateur, il forge la patience, la discipline et la dignité. Référence : le proverbe camerounais « la main qui travaille ne manque pas de manger » exprime la valeur du travail comme source d'autonomie. Exemple : l'artisan de Foumban qui sculpte et tire fierté de son savoir-faire transmis de père en fils.
\par
\emph{Paragraphe 2.} Idée : le travail est ce qui distingue l'homme de l'animal. Référence : \textbf{Hannah Arendt} (Condition de l'homme moderne) distingue le \emph{labeur} (asservi aux besoins vitaux, répétitif) de l'\emph{œuvre} (qui fabrique un monde durable) : tout travail n'est pas aliénant. Transition : « Le travail serait donc tantôt servitude, tantôt liberté ? Ne faut-il pas, plutôt que de condamner ou de glorifier le travail en général, distinguer ses formes ? »
\par\smallskip
\textbf{III. Synthèse : distinguer travail aliénant et travail épanouissant.}
\par
\emph{Paragraphe 1.} Idée : l'aliénation ne tient pas au travail lui-même, mais à ses conditions. Argument : le même geste peut être servitude (travail imposé, répétitif, dépossédé) ou accomplissement (travail choisi, créatif, reconnu) ; c'est le rapport social au travail qui est en cause, comme Marx lui-même le suggérait en rêvant d'un travail libéré.
\par
\emph{Paragraphe 2.} Idée : le vrai problème est donc : « quelles conditions rendent le travail humain ? » Réponse : un travail humain est un travail où l'homme garde l'initiative, comprend le sens de sa tâche et est reconnu. Exemple : les coopératives agricoles camerounaises où les producteurs sont associés aux décisions montrent qu'on peut concilier travail et dignité. Ce dépassement prépare la conclusion, qui répondra : le travail n'est pas par nature une aliénation ; il le devient quand ses conditions nient l'homme.
\end{exoresolu}

\begin{exercice}[À vous de construire]
Sujet : « La télévision nous rend-elle plus libres ? » Rédigez le plan détaillé du développement dialectique : pour chaque partie (thèse, antithèse, synthèse), formulez deux idées avec leur argument, un exemple camerounais et une référence possible, puis rédigez les deux transitions entre les parties.
\end{exercice}

\begin{corrige}[Exercice : « La télévision nous rend-elle plus libres ? »]
\textbf{Thèse} : la télévision libère. Idée 1 : elle informe et ouvre l'esprit (argument : l'accès à l'information est une condition de la liberté de juger ; exemple : les journaux télévisés de la CRTV ou les documentaires éducatifs). Idée 2 : elle divertit et détend (exemple : les émissions culturelles et sportives qui rassemblent les familles). Transition : « Mais un instrument qui informe peut aussi manipuler : la télévision ne nous rend-elle pas passifs ? »
\par
\textbf{Antithèse} : la télévision aliène. Idée 1 : elle rend spectateur passif au lieu de citoyen actif (argument : l'image s'impose sans laisser le temps de réfléchir ; exemple : les jeunes qui passent des heures devant les séries au détriment des études). Idée 2 : elle uniformise les désirs par la publicité (exemple : les spots qui créent des besoins artificiels). Transition : « Faut-il alors condamner la télévision, ou n'est-ce pas l'usage que nous en faisons qui est en cause ? »
\par
\textbf{Synthèse} : distinguer l'outil et l'usage. La télévision est un moyen : elle libère celui qui la regarde avec un esprit critique et la choisit ; elle aliène celui qui la subit sans distance. La vraie question devient : comment former des téléspectateurs critiques ? Ce qui renvoie au rôle de l'éducation.
\end{corrige}

\begin{aretenir}[L'essentiel du TP 2]
\begin{itemize}
\item Le développement dialectique suit trois moments : \cle{thèse} (première réponse, défendue fortement), \cle{antithèse} (critique et position opposée), \cle{synthèse} (dépassement par distinction, nuance ou reformulation).
\item Chaque partie compte 2 à 3 paragraphes ; chaque paragraphe suit la chaîne : \textbf{idée → argument → exemple → référence → transition}.
\item Ne jamais juxtaposer les opinions : chaque position doit être \emph{examinée critiquement}.
\item Les transitions font le bilan d'une partie et ouvrent la difficulté suivante : elles montrent le progrès de la réflexion.
\item La synthèse tranche et justifie ; elle n'est jamais un « les deux ont raison ».
\end{itemize}
\end{aretenir}

\coursec{TP 3 : La conclusion}

Après avoir construit le développement — la thèse, l'antithèse et la synthèse —, il reste à l'élève une dernière étape, souvent négligée et pourtant décisive : la \cle{conclusion}. Beaucoup de copies perdent des points précieux non pas parce que le développement est mauvais, mais parce que la conclusion est absente, bâclée ou hors sujet. Or la conclusion est la dernière impression laissée au correcteur du Baccalauréat technique : c'est elle qui montre que l'élève a réellement répondu à la question posée. Ce troisième travail pratique a pour objectif de consolider les acquis méthodologiques relatifs à la rédaction de la conclusion.

\courssub{La fonction de la conclusion}

\begin{definition}[La conclusion]
La \cle{conclusion} est la dernière partie de la dissertation philosophique. Elle a pour fonction de \textbf{répondre explicitement à la problématique} posée dans l'introduction, en tirant les conséquences du chemin parcouru dans le développement. Elle comporte trois moments : le \cle{bilan} de la démarche, la \cle{réponse} argumentée à la problématique et l'\cle{ouverture}.
\end{definition}

Commençons par une intuition simple. Imaginez un élève de Terminale technique qui part de Yaoundé pour se rendre à Douala. Le voyage comporte des étapes : Mbankomo, Boumnyebel, Edéa. Arrivé à Douala, personne ne répète tout l'itinéraire village par village ; on dit simplement : « nous sommes partis de Yaoundé, nous avons traversé la Sanaga, et nous voici arrivés ». La conclusion fonctionne de la même manière : elle \textbf{rappelle brièvement le trajet} (le bilan), elle \textbf{annonce l'arrivée} (la réponse à la question), et elle \textbf{indique un nouveau départ possible} (l'ouverture).

La conclusion n'est donc ni un résumé détaillé du développement, ni une simple formule de politesse. Elle est le \emph{verdict} du procès intellectuel que constitue la dissertation : après avoir entendu la thèse (l'accusation) et l'antithèse (la défense), le juge — c'est-à-dire l'élève — doit \textbf{rendre sa sentence} (la synthèse assumée comme position personnelle).

\begin{attention}[Une conclusion n'est jamais neutre]
L'erreur la plus fréquente au Baccalauréat technique est de conclure par des formules du type : « chacun a son opinion », « tout dépend des points de vue », « la question reste ouverte ». Ces formules sont \textbf{interdites} en dissertation philosophique : elles signifient que deux heures de réflexion n'ont abouti à rien. La dissertation exige une \textbf{position personnelle justifiée}. L'élève doit trancher, en s'appuyant sur la synthèse construite dans le développement. Relativiser toutes les opinions, c'est précisément refuser de philosopher.
\end{attention}

\courssub{Les trois moments de la conclusion}

\textbf{Premier moment : le bilan de la démarche.}\par
Le \cle{bilan} consiste à reformuler \emph{brièvement} les positions examinées dans le développement, sans répéter les arguments ni les exemples. Il s'agit de rappeler le chemin parcouru : quelle thèse avons-nous d'abord envisagée ? Quelle objection lui avons-nous opposée ? Quelle position dépassée avons-nous dégagée ? Le bilan se rédige en deux ou trois phrases, souvent introduites par des formules comme « Nous avons vu que... », « Notre réflexion nous a d'abord conduit à... ».

\begin{attention}[Ne pas répéter le développement]
Le bilan n'est pas un second développement en miniature. Si l'élève ressort tous ses exemples (les ouvriers de Bonabéri, les citations de Marx, etc.), il allonge inutilement sa copie et montre qu'il n'a pas compris la différence entre \emph{exposer} et \emph{conclure}. Le bilan ne retient que les \textbf{étapes} du raisonnement, pas son contenu détaillé.
\end{attention}

\textbf{Deuxième moment : la réponse à la problématique.}\par
C'est le cœur de la conclusion. L'élève doit ici \textbf{prendre clairement position} sur la question posée dans l'introduction, en s'appuyant sur la synthèse dégagée dans la troisième partie du développement. La réponse doit être formulée de manière directe et assumée : « Ainsi, nous pouvons affirmer que... », « Il apparaît donc que... ». Une réponse nuancée n'est pas une réponse vague : on peut répondre « oui, mais à certaines conditions » ou « non, sauf dans tel cas », à condition que la nuance soit \emph{argumentée} et non une fuite.

\textbf{Troisième moment : l'ouverture.}\par
L'\cle{ouverture} est une question nouvelle, liée au sujet, qui prolonge la réflexion sans la clore définitivement. Elle montre que la pensée continue au-delà de la copie. Par exemple, après un sujet sur la coutume, on peut ouvrir vers la question de l'identité culturelle dans la mondialisation : « Si la coutume fonde nos sociétés, comment les cultures camerounaises peuvent-elles demeurer elles-mêmes face à la mondialisation ? » L'ouverture doit rester \textbf{liée au sujet traité} : une ouverture sans rapport avec la problématique est une faute méthodologique.

\begin{popfigure}
\begin{center}
\begin{tikzpicture}[font=\small]
  % Entonnoir inversé : bilan (étroit) -> réponse -> ouverture (large)
  \fill[popblueL] (0,0) -- (4,0) -- (3.6,-1.6) -- (0.4,-1.6) -- cycle;
  \fill[poporangeL] (0.4,-1.8) -- (3.6,-1.8) -- (4.4,-3.4) -- (-0.4,-3.4) -- cycle;
  \fill[popgreenL] (-0.4,-3.6) -- (4.4,-3.6) -- (5.2,-5.2) -- (-1.2,-5.2) -- cycle;
  \node at (2,-0.8) {\textbf{1. Bilan} : « Nous avons vu que... »};
  \node at (2,-2.6) {\textbf{2. Réponse} : « Ainsi, nous affirmons que... »};
  \node at (2,-4.4) {\textbf{3. Ouverture} : « Reste à se demander... »};
  \draw[->, thick, popdark] (5.6,-0.8) -- (5.6,-4.4) node[midway, right, align=left] {élargissement\\ progressif};
\end{tikzpicture}
\legende{Schéma en entonnoir inversé de la conclusion : du bilan (resserrement du propos) vers la réponse, puis vers l'ouverture (élargissement de la réflexion).}
\end{center}
\end{popfigure}

\begin{methode}[Rédiger la conclusion en trois temps]
\begin{enumerate}
  \item \textbf{Le bilan} : commencer par « Nous avons vu que... » et rappeler en deux phrases les positions examinées (thèse, antithèse, synthèse), sans répéter les arguments.
  \item \textbf{La réponse} : poursuivre par « Ainsi... » ou « Nous pouvons donc affirmer que... » et répondre directement à la problématique de l'introduction, en assumant la synthèse comme position personnelle.
  \item \textbf{L'ouverture} : terminer par « Reste à se demander... » ou « Cette réflexion invite à s'interroger sur... » et poser une question nouvelle, liée au sujet, qui prolonge la réflexion.
\end{enumerate}
Vérifier enfin que la conclusion compte entre 8 et 12 lignes et qu'elle ne contient \textbf{aucun argument ni auteur nouveau}.
\end{methode}

\begin{attention}[Aucun élément nouveau dans la conclusion]
Il est \textbf{formellement interdit} d'introduire dans la conclusion un argument, un exemple, un auteur ou une thèse qui n'ont pas été traités dans le développement. La conclusion ne fait que \emph{tirer les conséquences} de ce qui a été démontré. Introduire Nietzsche ou un nouvel exemple à la dernière ligne, c'est avouer que le développement était incomplet — et le correcteur le sanctionne. Tout ce qui se trouve dans la conclusion doit pouvoir être rattaché à une étape du développement.
\end{attention}

\courssub{Les erreurs de conclusion et leur correction}

L'expérience des jurys du Baccalauréat technique permet de dresser un tableau des erreurs les plus fréquentes. Observer ces erreurs, c'est déjà apprendre à les éviter.

\begin{popfigure}
\begin{center}
\begin{tikzpicture}[font=\small]
  \node[draw=popink, fill=poppinkL, rounded corners, text width=4.6cm, align=center, minimum height=1.5cm] (e1) at (0,0) {\textbf{Conclusion absente}\\ la copie s'arrête après la synthèse};
  \node[draw=popink, fill=poppinkL, rounded corners, text width=4.6cm, align=center, minimum height=1.5cm] (e2) at (0,-2.4) {\textbf{Conclusion hors sujet}\\ réponse à une autre question que celle posée};
  \node[draw=popink, fill=poppinkL, rounded corners, text width=4.6cm, align=center, minimum height=1.5cm] (e3) at (0,-4.8) {\textbf{Nouvelle thèse}\\ un argument ou un auteur apparaît à la fin};
  \node[draw=popgreen, fill=popgreenL, rounded corners, text width=5.4cm, align=center, minimum height=1.5cm] (c1) at (8,0) {\textbf{Correction}\\ rédiger les trois moments : bilan, réponse, ouverture};
  \node[draw=popgreen, fill=popgreenL, rounded corners, text width=5.4cm, align=center, minimum height=1.5cm] (c2) at (8,-2.4) {\textbf{Correction}\\ relire la problématique de l'introduction et y répondre mot à mot};
  \node[draw=popgreen, fill=popgreenL, rounded corners, text width=5.4cm, align=center, minimum height=1.5cm] (c3) at (8,-4.8) {\textbf{Correction}\\ ne tirer que les conséquences des arguments déjà développés};
  \draw[->, very thick, popgreen] (e1) -- (c1);
  \draw[->, very thick, popgreen] (e2) -- (c2);
  \draw[->, very thick, popgreen] (e3) -- (c3);
\end{tikzpicture}
\legende{Tableau des erreurs fréquentes de conclusion et de leur correction : 1. conclusion absente ; 2. conclusion hors sujet ; 3. nouvelle thèse introduite à la fin.}
\end{center}
\end{popfigure}

\begin{aretenir}[L'essentiel sur la conclusion]
\begin{itemize}
  \item La conclusion répond \textbf{explicitement} à la problématique posée dans l'introduction.
  \item Elle comporte trois moments : le \cle{bilan} (« Nous avons vu que... »), la \cle{réponse} (« Ainsi... ») et l'\cle{ouverture} (« Reste à se demander... »).
  \item Longueur conseillée : \textbf{8 à 12 lignes}.
  \item Interdictions : conclure par « chacun a son opinion » ; introduire un argument ou un auteur nouveau ; répéter le développement ; s'écarter du sujet.
\end{itemize}
\end{aretenir}

\courssub{TP guidé : rédaction d'une conclusion}

\textbf{Consigne de travail.}\par
Sujet : \emph{« Le travail est-il une aliénation ? »} Le développement a suivi le plan suivant : (1) thèse — le travail est une aliénation (Marx : le travailleur est dépossédé du produit de son travail) ; (2) antithèse — le travail est une libération (Hegel : par le travail, l'homme se transforme et se reconnaît dans son œuvre) ; (3) synthèse — le travail aliène ou libère selon ses conditions (travail subi ou travail choisi et reconnu).

Rédigez \textbf{individuellement}, en 15 minutes, une conclusion de 8 à 12 lignes comportant les trois moments étudiés. La correction collective suivra la grille : bilan présent ? réponse claire à la problématique ? ouverture liée au sujet ? aucun élément nouveau ?

\begin{exoresolu}[Rédaction d'une conclusion — « Le travail est-il une aliénation ? »]
Énoncé : à partir du développement rappelé ci-dessus, rédiger la conclusion complète du devoir.
\tcblower
\textbf{Solution détaillée (corrigé type).}

\emph{Bilan} : « Nous avons vu que le travail peut d'abord apparaître comme une aliénation : dans certaines conditions, le travailleur est dépossédé du produit de son effort et se sent étranger à sa propre activité, comme le montre Marx. Mais notre réflexion nous a ensuite conduit à reconnaître, avec Hegel, que le travail est aussi le moyen par lequel l'homme transforme la nature, se transforme lui-même et accède à la conscience de sa valeur. »

\emph{Réponse} : « Ainsi, nous pouvons affirmer que le travail n'est pas par nature une aliénation : il aliène lorsqu'il est subi, répétitif et méconnu, mais il libère lorsqu'il est choisi, créateur et socialement reconnu. Tout dépend donc des conditions dans lesquelles il s'exerce. »

\emph{Ouverture} : « Reste à se demander si la société camerounaise offre aujourd'hui à tous ces conditions de libération : si le travail peut libérer, qu'en est-il du chômage des jeunes diplômés au Cameroun, qui prive cette libération de ses conditions mêmes ? »

\emph{Vérification méthodologique} : la conclusion compte environ 10 lignes ; elle répond directement à la question posée (« le travail n'est pas par nature une aliénation ») ; elle ne cite que Marx et Hegel, déjà mobilisés dans le développement ; l'ouverture prolonge le sujet sans le quitter.
\end{exoresolu}

\begin{exemplebox}[Une ouverture pertinente]
Une bonne ouverture relie le sujet à une réalité vécue. Pour « Le travail est-il une aliénation ? », l'ouverture suivante est excellente : « Si le travail peut libérer, qu'en est-il du chômage des jeunes diplômés au Cameroun, qui prive cette libération de ses conditions ? » Elle est pertinente parce qu'elle (1) part de la réponse donnée, (2) pose une question \emph{nouvelle} (le chômage, non traité dans le devoir), (3) reste \emph{liée} au sujet (le travail et la liberté). À l'inverse, ouvrir sur « la pollution de l'environnement » serait hors sujet.
\end{exemplebox}

\begin{exercice}[Entraînement personnel]
Sujet : « La coutume est-elle un obstacle au progrès ? » Rédigez une conclusion de 8 à 12 lignes comportant les trois moments, puis proposez une ouverture vers la question de l'identité culturelle dans la mondialisation. Faites corriger votre production par un camarade à l'aide de la grille du TP guidé.
\end{exercice}

\begin{corrige}[Exercice — « La coutume est-elle un obstacle au progrès ? »]
\emph{Éléments de correction.} Bilan attendu : rappel que la coutume peut freiner l'innovation (poids des traditions, refus du changement), mais qu'elle assure aussi la cohésion sociale et transmet des valeurs (solidarité, respect des aînés). Réponse attendue : position personnelle nuancée et assumée, par exemple : « la coutume n'est un obstacle au progrès que lorsqu'elle devient immobile ; intégrée à une réflexion critique, elle peut au contraire orienter le progrès ». Ouverture attendue : « Reste à se demander comment les cultures camerounaises peuvent préserver leur identité face à la mondialisation sans se fermer au progrès. » Barème indicatif : bilan 2 pts, réponse claire 3 pts, ouverture liée au sujet 2 pts, respect des interdictions (aucun élément nouveau, pas de « chacun a son opinion ») 3 pts.
\end{corrige}

La conclusion maîtrisée, il reste à consolider un dernier savoir-faire qui donne à toute la dissertation sa force persuasive : le choix et l'utilisation des citations, qui feront l'objet du TP 4.

\coursec{TP 4 : Choix et utilisation des citations}

Dans le TP 3, nous avons appris à conclure une dissertation : répondre clairement au problème, résumer le chemin parcouru, éventuellement ouvrir. Mais une dissertation ne vit pas seulement de sa structure. Entre l'introduction et la conclusion, le développement a besoin de \cle{références} : des textes d'auteurs, des formules célèbres, des proverbes. Encore faut-il savoir les choisir et les utiliser. Une citation mal employée peut ruiner une copie ; une citation bien intégrée la fait passer dans la bonne moyenne. Ce TP 4 vous donne les règles du jeu.

\courssub{Le rôle de la citation : appuyer, jamais remplacer}

Commençons par une intuition simple. Imaginez un élève de Terminale technique qui, au cours d'un débat, affirme : « La liberté, c'est faire ce qu'on veut. » Un camarade lui répond : « Ce n'est pas toi qui le dis, c'est tout le monde. » La remarque est juste : une opinion personnelle, sans appui, reste fragile. Maintenant, si l'élève précise : « Rousseau montre que la vraie liberté consiste à obéir à la loi qu'on s'est prescrite », sa position change de statut : elle s'adosse à une réflexion reconnue, elle devient défendable.

\begin{definition}[La citation en dissertation philosophique]
Une \cle{citation} est le fait de rapporter, entre guillemets, les paroles exactes d'un auteur (philosophe, écrivain, savant) ou d'une tradition (proverbe, texte sacré) pour \textbf{appuyer un argument} de sa propre réflexion. La citation est un \emph{outil au service} de la pensée personnelle du candidat : elle ne pense jamais à sa place.
\end{definition}

\begin{aretenir}[La règle d'or]
La citation \textbf{appuie} l'argument, elle ne le \textbf{remplace} pas. Une dissertation philosophique est d'abord un exercice de réflexion personnelle : le correcteur évalue \emph{votre} capacité à problématiser, à argumenter, à articuler. Les citations sont des renforts, des preuves d'autorité, des tremplins — jamais le contenu principal de la copie.
\end{aretenir}

Autrement dit, la citation joue trois rôles possibles :
\begin{itemize}
\item \textbf{Rôle d'autorité} : montrer qu'un grand penseur a défendu l'idée que vous avancez (« Comme le montre Kant... »).
\item \textbf{Rôle de formulation} : un auteur a parfois dit mieux que nous, en une phrase frappante, ce que nous cherchons à exprimer.
\item \textbf{Rôle de repoussoir} : citer une thèse pour la discuter ou la critiquer ensuite (« Descartes affirme que... ; mais cette position se heurte à une difficulté... »).
\end{itemize}

\courssub{Les règles de choix d'une citation}

Toutes les citations ne se valent pas pour un sujet donné. Trois critères permettent de faire le bon choix.

\begin{propriete}[Les trois critères de choix d'une citation]
\begin{enumerate}
\item \textbf{La pertinence} : la citation doit avoir un rapport direct et vérifiable avec le sujet et avec l'argument précis qu'elle vient soutenir. Une citation « à peu près dans le thème » affaiblit la copie.
\item \textbf{L'exactitude} : l'auteur doit être correctement identifié et la formulation fidèle. Attribuer à Platon une phrase de Nietzsche est une faute grave aux yeux du correcteur.
\item \textbf{La brièveté} : une citation efficace tient en une ou deux lignes. Les longs extraits recopiés donnent l'impression de « remplissage » et masquent l'absence de réflexion personnelle.
\end{enumerate}
\end{propriete}

\begin{attention}[Piège : la citation décorative]
Beaucoup de candidats placent une citation célèbre en tête de copie ou au milieu d'un paragraphe simplement parce qu'elle « fait sérieux ». Si le lien avec l'argument n'est pas établi, la citation est \emph{décorative} : le correcteur s'en aperçoit immédiatement et sanctionne. Avant d'écrire une citation, posez-vous la question : « Quel argument précis vient-elle renforcer ? » Si vous ne trouvez pas de réponse, ne la mettez pas.
\end{attention}

\courssub{Les règles d'utilisation : introduire, citer, commenter}

Une citation ne se jette pas dans la copie comme une pierre dans l'eau. Elle suit un circuit en trois temps, qu'il faut rendre visible dans la rédaction.

\begin{popfigure}
\begin{center}
\begin{tikzpicture}[
  box/.style={draw=popblue, thick, rounded corners=4pt, fill=popblueL, text width=3.4cm, minimum height=2.6cm, align=center, font=\small},
  arr/.style={-{Stealth[length=3mm]}, very thick, poporange}
]
\node[box, align=center] (a) at (0,0){\textbf{1. Introduire}\\[2pt] Qui parle ? Dans quelle œuvre, quel contexte ? Pourquoi cet auteur ici ?};
\node[box, align=center] (b) at (5.2,0){\textbf{2. Citer}\\[2pt] Rapporter les paroles exactes, entre guillemets, brièvement et fidèlement.};
\node[box, align=center] (c) at (10.4,0){\textbf{3. Commenter}\\[2pt] Reformuler en ses propres termes et montrer l'apport à l'argumentation.};
\draw[arr] (a) -- (b);
\draw[arr] (b) -- (c);
\draw[arr, popgreen, dashed] (c.south) to[bend right=25] node[below, font=\small, text=popgreen]{la réflexion personnelle continue} (a.south);
\end{tikzpicture}
\end{center}
\legende{Le circuit d'intégration d'une citation : introduire (1), citer (2), commenter (3). La citation s'insère dans le flux de la réflexion personnelle, elle ne l'interrompt pas.}
\end{popfigure}

\begin{methode}[La formule d'intégration d'une citation]
Pour intégrer correctement une citation, suivez ce schéma en trois étapes :
\begin{enumerate}
\item \textbf{Introduire} : « Comme le souligne [auteur] dans [œuvre ou contexte]... »
\item \textbf{Citer} : « ... : “[paroles exactes, entre guillemets]”. »
\item \textbf{Commenter} : « — autrement dit, [reformulation en vos propres termes], ce qui montre que [lien explicite avec votre argument]. »
\end{enumerate}
Modèle complet : \emph{« Comme le souligne [auteur], “...” — autrement dit, [reformulation], ce qui montre que [lien avec l'argument]. »}
\end{methode}

\begin{exemplebox}[Une citation correctement intégrée]
Kant affirme dans \emph{Réponse à la question : qu'est-ce que les Lumières ?} que « l'illumination est la sortie de l'homme de sa minorité » : autrement dit, penser par soi-même suppose le courage de se passer des tutelles qui nous gouvernent, ce qui rejoint notre critique de l'obéissance aveugle à la coutume.
\end{exemplebox}

Analysons cet exemple : l'auteur est nommé (Kant), le contexte est donné (l'opuscule sur les Lumières), la citation est brève et exacte, la reformulation (« autrement dit ») prouve que le candidat a compris, et le lien avec l'argument du devoir est explicite (« ce qui rejoint notre critique... »). Tout y est.

\courssub{Le commentaire de citation : reformuler et exploiter}

Le commentaire est l'étape la plus souvent négligée, et pourtant la plus importante. Commenter une citation, c'est faire deux choses :

\begin{itemize}
\item \textbf{Reformuler} : traduire la pensée de l'auteur en vos propres termes. Cela prouve au correcteur que vous n'avez pas appris la citation par cœur sans la comprendre. Les formules utiles sont : « autrement dit », « c'est-à-dire », « en d'autres termes », « ce que l'auteur signifie ici, c'est que... ».
\item \textbf{Exploiter} : tirer de la citation une conséquence pour votre argumentation. La citation doit \emph{faire avancer} le raisonnement : elle confirme une étape, elle autorise une transition, elle ouvre une objection à discuter.
\end{itemize}

\begin{attention}[Erreur fréquente : le catalogue de citations]
Enchaîner les citations sans commentaire — « Platon dit que... Aristote dit que... Kant dit que... » — transforme la dissertation en catalogue. Le correcteur n'y voit pas une réflexion, mais une liste de course apprise par cœur. Règle pratique : \textbf{une citation = au moins deux phrases de commentaire}. Mieux vaut deux citations bien exploitées que six citations plaquées.
\end{attention}

\begin{exoresolu}[Transformer une citation plaquée en citation exploitée]
\textbf{Énoncé.} Voici un extrait de copie sur le sujet « Le travail libère-t-il l'homme ? » : \emph{« Hegel dit que “le travail est le désir réfréné”. Marx dit que le travail aliène l'ouvrier. Donc le travail est ambivalent. »} Corrigez cet extrait.
\tcblower
\textbf{Solution détaillée.} L'extrait juxtapose deux citations sans les introduire ni les commenter : la conclusion « donc le travail est ambivalent » tombe sans être démontrée. Version corrigée : « Hegel, dans la \emph{Phénoménologie de l'esprit} (§190), définit le travail comme “le désir réfréné” : autrement dit, en travaillant, l'homme renonce à détruire immédiatement l'objet de son désir et le transforme durablement, ce qui montre que le travail éduque et humanise. Cependant, Marx observe dans les \emph{Manuscrits de 1844} que le travailleur salarié est séparé du produit de son travail : le travail peut donc aussi aliéner, c'est-à-dire déposséder l'homme de lui-même. Ainsi, le travail n'est pas libérateur par nature : tout dépend des conditions dans lesquelles il s'exerce. » Ici, chaque citation est introduite (auteur, œuvre), citée brièvement, reformulée, et mise au service d'une conclusion nuancée.
\end{exoresolu}

\courssub{Citation bien utilisée, citation mal utilisée : tableau comparatif}

\begin{popfigure}
\begin{center}
\begin{tikzpicture}
\node[draw=popgreen, thick, rounded corners=4pt, fill=popgreenL, text width=6.6cm, align=left, font=\small, anchor=north west] (good) at (0,0)
{\textbf{Bien utilisée}\\[3pt]
« Comme le souligne Sénèque, “ce n'est pas que nous ayons peu de temps, c'est que nous en perdons beaucoup” : autrement dit, le temps ne manque pas, c'est notre usage qui le gaspille, ce qui invite à philosopher sur l'emploi de notre existence. »\\[3pt]
\textcolor{popgreen}{Introduite, citée exactement, reformulée, reliée à l'argument.}};
\node[draw=poppink, thick, rounded corners=4pt, fill=poppinkL, text width=6.6cm, align=left, font=\small, anchor=north east] (bad) at (14.6,0)
{\textbf{Mal utilisée}\\[3pt]
« Sénèque a dit que le temps passe vite. Platon a dit que le temps est l'image mobile de l'éternité. Einstein a dit que le temps est relatif. Le temps est donc un mystère. »\\[3pt]
\textcolor{poppink}{Catalogue sans commentaire, formulations approximatives, aucun lien argumentatif.}};
\end{tikzpicture}
\end{center}
\legende{Deux usages contrastés de la citation : à gauche, une citation intégrée selon le circuit introduire--citer--commenter ; à droite, un catalogue de citations plaquées qui ne fait pas avancer la réflexion.}
\end{popfigure}

\courssub{Quand on ne se souvient plus exactement : la paraphrase honnête}

Au baccalauréat, la mémoire fait parfois défaut : vous savez que tel auteur a défendu telle idée, mais la formulation exacte vous échappe. Que faire ?

\begin{attention}[Ne jamais inventer une citation]
Il est \textbf{formellement interdit} d'inventer une citation ou de l'attribuer au mauvais auteur. Une fausse citation entre guillemets est une faute d'honnêteté intellectuelle que le correcteur repère souvent — et elle discrédite toute la copie. La solution : \textbf{paraphraser}. Si la formulation exacte vous échappe, nommez l'auteur et exposez son idée sans guillemets : « Selon Nietzsche, l'homme est quelque chose qui doit être dépassé » devient, en paraphrase honnête : « Nietzsche soutient que l'humanité n'est pas un aboutissement mais un stade à surmonter. » Une paraphrase exacte vaut infiniment mieux qu'une fausse citation.
\end{attention}

\begin{aretenir}[Guillemets = exactitude]
Les guillemets sont un contrat : ils garantissent que les mots rapportés sont \emph{exactement} ceux de l'auteur. Si vous n'êtes pas sûr de la formulation, retirez les guillemets et paraphrasez en nommant l'auteur. Vous gardez ainsi le bénéfice de la référence sans risquer la faute.
\end{aretenir}

\courssub{TP guidé : choisir la bonne citation pour un sujet}

\textbf{Sujet : « Peut-on se libérer de ses passions ? »}

Voici une liste de références. À vous de choisir celle qui convient le mieux et de l'intégrer selon la méthode.

\begin{center}
\begin{tabularx}{\linewidth}{|l|X|l|}
\hline
\rowcolor{popblueL} \textbf{Auteur / source} & \textbf{Référence} & \textbf{Pertinence} \\
\hline
Platon & Le mythe de l'attelage ailé : l'âme est un cocher conduisant un bon et un mauvais cheval (les passions) & Forte (thèse : les passions peuvent être dirigées) \\
\hline
Kant & « L'illumination est la sortie de l'homme de sa minorité » & Faible (sujet : la raison, non les passions) \\
\hline
Nietzsche & « Deviens ce que tu es » (\emph{Ecce Homo} / \emph{Ainsi parlait Zarathoustra}) ; les instincts doivent être affirmés, non réprimés & Forte (antithèse : non se libérer, mais assumer) \\
\hline
Sénèque & « Ce n'est pas que nous ayons peu de temps, c'est que nous en perdons beaucoup » & Nulle (sujet : le temps) \\
\hline
Proverbe africain & « Un seul doigt ne ramasse pas la farine » & Moyenne (utilisable pour la dimension sociale de la maîtrise de soi) \\
\hline
\end{tabularx}
\end{center}

\textbf{Corrigé du choix.} Les citations de Kant et de Sénèque sont hors sujet : les utiliser serait décoratif. Platon convient pour la première partie (la raison peut dominer les passions, comme le cocher dirige l'attelage) ; Nietzsche convient pour l'antithèse (vouloir supprimer les passions, c'est mutiler la vie : il faut les canaliser et les affirmer) ; le proverbe peut servir en synthèse (la maîtrise de soi s'apprend aussi dans la communauté). Exemple d'intégration de Nietzsche : « À l'inverse, Nietzsche invite à “devenir ce que l'on est” : autrement dit, les passions ne sont pas des ennemies à éliminer mais des forces à orienter, ce qui suggère que la liberté n'est pas l'absence de passions mais leur mise en ordre. »

\begin{savaistu}[La sagesse orale africaine, une référence légitime]
Les proverbes et la tradition orale africaine — par exemple « Un seul doigt ne ramasse pas la farine » (solidarité), « Quand un vieillard meurt, c'est une bibliothèque qui brûle » (transmission, formule attribuée à Amadou Hampâté Bâ) — peuvent servir de références \textbf{légitimes} dans une dissertation, à condition d'être \emph{analysés philosophiquement} : il faut dégager l'idée qu'ils expriment, l'interroger, la confronter aux thèses des philosophes. Un proverbe cité sans analyse reste du folklore ; un proverbe problématisé devient de la philosophie.
\end{savaistu}

\courssub{Constituer son répertoire personnel de citations}

Pour être prêt le jour du baccalauréat, il est vivement conseillé de constituer, tout au long de l'année, un \cle{répertoire personnel} de 15 à 20 citations, classées par thème du programme de Terminale technique.

\begin{methode}[Construire son répertoire de citations]
\begin{enumerate}
\item Créez une fiche par grand thème du programme : la conscience, l'inconscient, le temps, la liberté, le devoir, le bonheur, la vérité, l'État, la justice, le travail, la technique, l'art, la religion, la mort, l'autrui.
\item Pour chaque thème, retenez \textbf{une ou deux citations courtes} (une ligne maximum), vérifiées dans le manuel ou le cours, avec l'auteur et si possible l'œuvre.
\item Pour chaque citation, rédigez \textbf{une phrase de reformulation} : si vous savez la reformuler, vous la comprenez, et vous saurez l'exploiter.
\item Apprenez le répertoire progressivement et réutilisez-le dans chaque dissertation d'entraînement : une citation n'est vraiment acquise que lorsqu'elle a été employée.
\end{enumerate}
\end{methode}

\begin{exercice}[À faire chez soi]
Constituez les trois premières fiches de votre répertoire (thèmes : la liberté, le devoir, la vérité). Pour chaque thème : une citation vérifiée, son auteur, une reformulation en vos propres termes, et un sujet de baccalauréat où elle pourrait servir.
\end{exercice}

\begin{corrige}[Exercice : répertoire de citations]
Exemple de fiche attendue — Thème : la liberté. Citation : « La liberté consiste à obéir à la loi qu'on s'est prescrite » (Rousseau, \emph{Du contrat social}, Livre I, ch. 8). Reformulation : être libre, ce n'est pas faire ce qu'on veut, mais se donner à soi-même sa règle d'action ; l'obéissance à soi-même est la vraie autonomie. Sujet possible : « La liberté est-elle l'absence de contraintes ? » — la citation servirait dans l'antithèse ou la synthèse pour montrer que la loi n'est pas l'ennemie de la liberté mais sa condition. Les fiches « devoir » et « vérité » se construisent sur le même modèle (par exemple Kant pour le devoir : « Agis de telle sorte que la maxime de ta volonté puisse valoir comme principe d'une législation universelle », \emph{Fondements de la métaphysique des mœurs}).
\end{corrige}

\begin{aretenir}[Bilan du TP 4]
\begin{itemize}
\item La citation \textbf{appuie} l'argument, elle ne le remplace jamais.
\item Choix : \textbf{pertinence}, \textbf{exactitude}, \textbf{brièveté}.
\item Utilisation : \textbf{introduire} (qui ? où ?), \textbf{citer} (entre guillemets, fidèlement), \textbf{commenter} (reformuler + exploiter).
\item Une citation = au moins deux phrases de commentaire ; jamais de catalogue.
\item En cas de doute sur la formulation : \textbf{paraphraser} en nommant l'auteur, sans guillemets.
\item Les proverbes africains sont des références légitimes s'ils sont analysés philosophiquement.
\item Constituez un répertoire de 15 à 20 citations classées par thème du programme.
\end{itemize}
\end{aretenir}

Armés de ces règles, nous pouvons maintenant passer à la dernière étape de notre consolidation : l'application de l'ensemble de la méthode — introduction, développement, conclusion, citations — à des situations concrètes et à des sujets d'examen complets.

\coursec{Application à la vie courante et entraînement à l'examen}

Dans les sections précédentes, nous avons travaillé chaque pièce du moteur séparément : l'introduction, le développement en trois moments, la conclusion, et enfin le choix et l'utilisation des citations. Il est temps maintenant d'assembler toutes ces pièces et de faire tourner le moteur en conditions réelles. Car une méthode ne vaut que si elle sait s'appliquer à des sujets nouveaux, et en particulier à ces situations de la vie quotidienne camerounaise que le Baccalauréat technique aime transformer en sujets de dissertation. Cette dernière section est donc un véritable entraînement au combat : nous apprendrons à convertir une situation de vie en problématique, à gérer notre temps, puis nous réaliserons une simulation complète d'épreuve.

\courssub{De la situation de vie à la problématique philosophique}

\textbf{Pourquoi partir de la vie courante ?}\par

Commençons par une intuition simple. Quand deux familles se disputent à propos d'une dot exigée à Ngaoundéré, quand un proviseur confisque les téléphones portables dans un lycée technique de Douala, quand un jeune reproche à son père de « ne rien comprendre à la vie moderne », quand un usager glisse un billet à un gendarme sur l'axe Yaoundé--Bafoussam : ce ne sont pas seulement des faits divers. Ce sont des \cle{conflits de valeurs}. Or la philosophie est précisément la discipline qui pense les valeurs : le juste, le bien, la liberté, la tradition, le devoir. Toute situation de vie où deux valeurs s'affrontent est donc un sujet de dissertation en puissance.

L'examinateur du Baccalauréat technique le sait bien : il prend une tension vécue par les élèves et la reformule en question abstraite. Notre travail consiste à faire le chemin inverse : partir de la question abstraite, mais l'alimenter sans cesse du concret que nous connaissons.

\begin{methode}[Problématiser une situation de vie]
Pour transformer une situation concrète en problématique philosophique, suivre quatre étapes :
\begin{enumerate}
\item \textbf{Décrire la situation} en une phrase neutre, sans prendre parti (exemple : « Au Cameroun, la famille du mari verse une dot à la famille de l'épouse »).
\item \textbf{Identifier les valeurs en tension} : chercher les couples d'opposés que la situation met en jeu — \cle{tradition/modernité}, \cle{liberté/autorité}, \cle{individu/communauté}, \cle{devoir/intérêt}, \cle{égalité/hiérarchie}.
\item \textbf{Formuler la question} que ces valeurs soulèvent, sous forme d'alternative ou de question ouverte, en utilisant les notions du programme (devoir, liberté, culture, État, justice, technique).
\item \textbf{Vérifier la problématique} : elle doit être discutable (on peut y répondre oui \emph{et} non), générale (elle dépasse le cas particulier) et philosophique (elle porte sur une valeur, pas sur un fait).
\end{enumerate}
\end{methode}

\begin{exemplebox}[La dot : de la situation au sujet]
Appliquons la méthode à la dot.
\begin{itemize}
\item \textbf{Situation} : dans de nombreuses régions du Cameroun, le mariage coutumier exige que la famille du prétendant remette des biens et de l'argent à la famille de la jeune fille.
\item \textbf{Valeurs en tension} : d'un côté, la \emph{tradition} et la \emph{communauté} (la dot scelle l'alliance entre deux familles, honore la femme, transmet la culture) ; de l'autre, la \emph{liberté} individuelle et l'\emph{égalité} (la dot peut enfermer la femme dans un statut d'objet acheté, empêcher le divorce, exclure les jeunes sans ressources du mariage).
\item \textbf{Problématique} : « La dot est-elle une tradition à préserver ou une marchandisation de la femme ? »
\item \textbf{Vérification} : la question est discutable (les deux réponses se défendent), générale (elle vaut pour toutes les coutumes) et philosophique (elle porte sur la valeur de la culture et sur la dignité de la personne — notions du programme).
\end{itemize}
On remarque que cette problématique ouvre naturellement sur une thèse (la dot est un héritage culturel précieux), une antithèse (la dot réduit la femme à une marchandise) et une synthèse (distinguer l'esprit de la coutume de ses dérives matérielles).
\end{exemplebox}

\begin{exemplebox}[Trois autres situations à problématiser]
Entraînons-nous rapidement :
\begin{itemize}
\item \textbf{Le téléphone mobile à l'école} : tension entre la \emph{liberté} de communication et l'\emph{autorité} éducative, entre l'outil technique et la distraction. Problématique possible : « L'usage du téléphone portable en classe est-il un droit de l'élève ou un obstacle à l'éducation ? »
\item \textbf{Les conflits entre générations} : tension entre l'expérience des anciens et l'autonomie des jeunes, entre fidélité et innovation. Problématique : « Les jeunes doivent-ils obéir aux anciens ou penser par eux-mêmes ? »
\item \textbf{La corruption et le civisme} : tension entre l'intérêt individuel immédiat et le devoir envers la communauté, entre la survie et la justice. Problématique : « Le civisme est-il un devoir ou une affaire de circonstances ? »
\end{itemize}
\end{exemplebox}

\courssub{La démarche complète : vue d'ensemble}

Avant la simulation d'examen, fixons dans une frise l'enchaînement des opérations que nous avons apprises dans toute la leçon. Cette frise doit devenir un réflexe : chaque flèche correspond à une étape obligée, aucune ne peut être sautée.

\begin{popfigure}
\begin{center}
\begin{tikzpicture}[
  etape/.style={rectangle, rounded corners=3pt, draw=popblue, fill=popblueL, text width=2.6cm, align=center, minimum height=1.1cm, font=\small},
  fleche/.style={-{Stealth[length=2.5mm]}, thick, popdark}
]
\node[etape, align=center] (a) at (0,0){\textbf{1. Analyse du sujet}\\ mots-clés, problématique};
\node[etape, align=center] (b) at (3.6,0){\textbf{2. Introduction}\\ accroche, problème, plan};
\node[etape, align=center] (c) at (7.2,0){\textbf{3. Thèse}\\ première réponse argumentée};
\node[etape, align=center] (d) at (10.8,0){\textbf{4. Antithèse}\\ objection argumentée};
\node[etape, align=center] (e) at (3.6,-2.2){\textbf{5. Synthèse}\\ dépassement, réponse nuancée};
\node[etape, align=center] (f) at (7.2,-2.2){\textbf{6. Conclusion}\\ bilan + ouverture};
\draw[fleche] (a) -- (b);
\draw[fleche] (b) -- (c);
\draw[fleche] (c) -- (d);
\draw[fleche] (d.south) |- (e.east);
\draw[fleche] (e) -- (f);
\end{tikzpicture}
\end{center}
\legende{Frise récapitulative de la démarche complète de la dissertation philosophique : six étapes enchaînées, de l'analyse du sujet à la conclusion. Les étapes 3, 4 et 5 constituent le développement.}
\end{popfigure}

\begin{exoresolu}[Rédiger et justifier une démarche complète sur la dot]
Énoncé : à partir de la problématique « La dot est-elle une tradition à préserver ou une marchandisation de la femme ? », rédiger le squelette complet d'une dissertation (introduction rédigée, plan détaillé du développement, conclusion rédigée) et justifier chaque choix.
\tcblower
\textbf{Introduction rédigée} : « Au Cameroun, aucun mariage coutumier ne se célèbre sans la remise de la dot, ce don que la famille du mari offre à celle de l'épouse. Pratique ancienne et vivante, elle divise pourtant les esprits : les uns y voient le ciment de nos cultures, les autres une survivance qui avilit la femme. Le problème est donc de savoir si la dot est une tradition à préserver ou une marchandisation de la femme. Nous verrons d'abord que la dot fonde l'alliance et honore la femme ; nous montrerons ensuite qu'elle peut la réduire au rang de marchandise ; nous soutiendrons enfin qu'il faut distinguer l'esprit de la coutume de ses dérives. »
\emph{Justification} : accroche factuelle et locale, définition implicite du terme clé, problématique reprise mot pour mot, annonce de plan en trois temps.

\textbf{Développement (plan détaillé)} :
\begin{itemize}
\item \emph{Thèse} : la dot est une tradition à préserver. Argument 1 : elle scelle l'alliance entre deux familles et non entre deux individus seuls (exemple : les négociations familiales chez les Bamiléké ou les Ewondo). Argument 2 : elle honore la femme en reconnaissant sa valeur (référence : la coutume comme identité, notion de culture au programme).
\item \emph{Antithèse} : la dot est une marchandisation. Argument 1 : quand elle devient exorbitante, elle fait de la femme un bien acheté et du mari un propriétaire. Argument 2 : elle bloque le divorce (la femme restée « payée ») et exclut les jeunes pauvres du mariage (exemple : jeunes salariés de Douala reportant le mariage faute de moyens).
\item \emph{Synthèse} : le problème n'est pas la dot elle-même mais sa dérive mercantile. Il faut préserver le geste symbolique (alliance, reconnaissance) en régulant les montants, comme le tentent certaines autorités coutumières et municipales qui plafonnent les listes de dot.
\end{itemize}

\textbf{Conclusion rédigée} : « La dot n'est donc ni à abolir ni à sacraliser : tradition à préserver dans sa signification d'alliance, elle devient condamnable dès qu'elle se transforme en transaction. Ouverte sur l'avenir, la question demeure : toute tradition peut-elle être réformée sans être trahie ? »
\emph{Justification} : la conclusion répond explicitement à la problématique, résume le dépassement obtenu, et l'ouverture prolonge la réflexion sans introduire un nouveau sujet.
\end{exoresolu}

\courssub{Simulation d'examen : traiter un sujet en conditions réelles}

\textbf{La gestion du temps}\par

L'épreuve de philosophie au Baccalauréat technique dure trois heures. L'erreur la plus coûteuse n'est pas de mal philosopher, c'est de mal répartir son temps : qui rédige sans plan produit un devoir incohérent ; qui perfectionne son plan pendant deux heures n'a plus le temps d'écrire.

\begin{exemplebox}[Répartition chiffrée du temps à l'examen (3 heures)]
Pour une épreuve de 3 h, la répartition recommandée est la suivante :
\begin{itemize}
\item \textbf{30 minutes} : analyse du sujet (mots-clés, problématique), recherche d'idées au brouillon, construction du plan détaillé ;
\item \textbf{2 heures} : rédaction au propre (introduction, trois parties du développement, conclusion), soit environ 30 minutes par grande partie ;
\item \textbf{30 minutes} : relecture complète (orthographe, grammaire, cohérence, transitions, citations).
\end{itemize}
Règle d'or : à la fin de la première demi-heure, le plan doit être arrêté, même imparfait. Un plan moyen bien rédigé vaut mieux qu'un plan parfait jamais écrit.
\end{exemplebox}

\textbf{La grille de correction}\par

Pour s'entraîner efficacement, il faut connaître les critères du correcteur. La grille ci-dessous reproduit l'esprit des grilles harmonisées utilisées au Baccalauréat technique camerounais (sur 20 points).

\begin{popfigure}
\begin{center}
\begin{tikzpicture}
\node[anchor=north west, align=center] at (0,0){
\begin{tabular}{|p{4.6cm}|p{6.2cm}|c|}
\hline
\rowcolor{popblueL} \textbf{Critère} & \textbf{Ce qui est évalué} & \textbf{Points} \\
\hline
Pertinence et compréhension du sujet & Sujet compris, problématique dégagée, pas de hors-sujet & 4 \\
\hline
Structure et organisation & Introduction, thèse, antithèse, synthèse, conclusion ; transitions & 4 \\
\hline
Argumentation & Arguments développés, exemples concrets, objections traitées & 5 \\
\hline
Références philosophiques & Auteurs et doctrines cités à bon escient et exactement & 4 \\
\hline
Langue et présentation & Orthographe, grammaire, clarté, lisibilité de la copie & 3 \\
\hline
\textbf{Total} & & \textbf{20} \\
\hline
\end{tabular}
};
\end{tikzpicture}
\end{center}
\legende{Grille d'évaluation type Baccalauréat technique : cinq critères totalisant 20 points. L'argumentation est le critère le plus lourd (5 points), mais la langue pèse 3 points : la relecture est rentable.}
\end{popfigure}

\begin{attention}[Ne jamais négliger la relecture]
L'erreur fréquente consiste à rédiger jusqu'à la dernière seconde et à rendre la copie sans la relire. Or les fautes de langue (orthographe, accords, conjugaison), les incohérences (une conclusion qui contredit la synthèse) et les oublis (une citation attribuée au mauvais auteur) coûtent des points réels : le critère « langue et présentation » vaut à lui seul environ 3 points, et une incohérence interne dégrade aussi la note d'argumentation. Trente minutes de relecture peuvent faire gagner 2 à 3 points : c'est l'investissement le plus rentable de l'épreuve. Réservez-les \emph{avant} de commencer à écrire, non pas « s'il reste du temps ».
\end{attention}

\begin{savaistu}[Des sujets proches de la vie sociale]
Les sujets du Baccalauréat technique camerounais portent régulièrement sur des notions directement liées à la vie sociale : le devoir, la liberté, l'État, la justice, la culture, la technique. Des sujets comme « Faut-il obéir aux lois ? », « La tradition est-elle un obstacle au progrès ? » ou « La technique libère-t-elle l'homme ? » sont typiques. S'entraîner sur des situations locales — la dot, le téléphone à l'école, les relations entre générations, le civisme face à la corruption — est donc une stratégie payante : elle fournit des exemples authentiques que le correcteur valorise, et elle prépare précisément aux notions les plus fréquentes de l'épreuve.
\end{savaistu}

\courssub{Auto-évaluation et correction entre pairs}

La correction entre pairs est un exercice formateur à double titre : celui qui corrige apprend à lire une copie avec l'œil du jury, et celui qui est corrigé reçoit un retour immédiat et sans enjeu. Voici la démarche à suivre en classe.

\begin{methode}[Corriger la copie d'un camarade avec la grille]
\begin{enumerate}
\item Lire la copie une première fois sans rien noter, pour saisir l'ensemble.
\item Relire critère par critère : la problématique est-elle formulée ? Les trois parties sont-elles présentes et équilibrées ? Chaque argument est-il développé et illustré ? Les références sont-elles exactes ? La langue est-elle correcte ?
\item Attribuer une note par critère en s'appuyant sur la grille, et justifier chaque note par une remarque écrite précise (jamais « c'est bien », mais « ta deuxième partie manque d'exemple »).
\item Formuler un conseil unique et prioritaire : l'\emph{unique} amélioration qui ferait gagner le plus de points.
\item Rendre la copie et discuter : l'auteur peut défendre ses choix, le correcteur doit justifier les siens.
\end{enumerate}
\end{methode}

\begin{exercice}[Simulation complète d'examen]
En conditions réelles (3 heures, sans documents, copie double), traiter l'un des sujets suivants :
\begin{enumerate}
\item « La dot est-elle une tradition à préserver ou une marchandisation de la femme ? »
\item « L'usage du téléphone portable en classe est-il un droit de l'élève ou un obstacle à l'éducation ? »
\item « Le civisme est-il un devoir ou une affaire de circonstances ? »
\end{enumerate}
À l'issue de l'épreuve, échanger les copies par binôme et les corriger à l'aide de la grille d'évaluation, puis comparer l'auto-évaluation de chacun avec la note attribuée par le pair.
\end{exercice}

\begin{corrige}[Exercice — indications de correction]
Il n'existe pas de corrigé unique en philosophie, mais des exigences communes. Pour le sujet 1, la copie doit définir la dot, opposer la valeur d'alliance (thèse) à la réduction marchande (antithèse), et aboutir à une distinction entre esprit et dérive de la coutume. Pour le sujet 2, on attend la tension entre liberté individuelle et finalité de l'école, avec une synthèse possible sur l'usage raisonné et encadré de la technique. Pour le sujet 3, la thèse défend le civisme comme devoir inconditionnel (référence possible : Kant et le devoir), l'antithèse invoque les circonstances (pauvreté, mauvais exemple des élites), et la synthèse montre que l'excuse des circonstances détruit la vie commune. Dans tous les cas : problématique explicite (4 pts), trois parties équilibrées avec transitions (4 pts), arguments développés et exemples camerounais (5 pts), au moins deux références exactes (4 pts), langue soignée (3 pts). L'écart entre auto-évaluation et note du pair est un excellent indicateur des points à retravailler.
\end{corrige}

\courssub{Bilan de la leçon : ma méthode en 10 points}

Pour clore cette leçon de consolidation, chaque élève rédige sa fiche de synthèse personnelle. En voici le modèle, à recopier et à garder dans son classeur jusqu'au jour de l'examen.

\begin{aretenir}[Ma méthode de dissertation en 10 points]
\begin{enumerate}
\item Je lis le sujet trois fois et je souligne les mots-clés avant toute écriture.
\item Je définis chaque terme important et je formule la problématique par écrit.
\item Je consacre 30 minutes au brouillon : idées, exemples, références, plan en trois parties.
\item Je rédige une introduction en quatre temps : accroche, définitions, problématique, annonce du plan.
\item Je développe une thèse argumentée, avec au moins deux arguments et un exemple concret.
\item Je développe une antithèse aussi forte que la thèse, sans la caricaturer.
\item Je construis une synthèse qui dépasse l'opposition au lieu de la répéter.
\item Je conclus en répondant clairement à la problématique, avec une ouverture mesurée.
\item Je n'utilise que des citations exactes, brèves, attribuées et expliquées.
\item Je réserve 30 minutes de relecture : orthographe, cohérence, transitions.
\end{enumerate}
\end{aretenir}

Ainsi s'achève notre parcours de consolidation. La dissertation philosophique n'est pas un don réservé à quelques-uns : c'est une technique, au sens noble du terme, c'est-à-dire un savoir-faire qui s'acquiert par la méthode et l'entraînement. En partant des situations que vous vivez chaque jour au Cameroun, en les élevant jusqu'aux notions du programme, puis en redescendant vers l'exemple concret, vous disposez désormais de tous les outils pour affronter l'épreuve du Baccalauréat technique avec confiance.

\newpage

\coursec{Activité d'intégration}

\coursec{Leçon 14 : Consolidation des acquis de la méthodologie de la dissertation philosophique}

\courssub{Objectifs de l'activité}

À l'issue de cette activité d'intégration, l'élève doit être capable de :
\begin{itemize}
\item analyser un sujet de dissertation philosophique (dégager les termes clés, identifier la tension, formuler une problématique) ;
\item rédiger une introduction complète et conforme aux exigences du Baccalauréat technique (accroche, analyse des termes, problématique, annonce du plan) ;
\item construire un développement dialectique en trois moments (thèse, antithèse, synthèse) avec arguments, exemples et références ;
\item choisir, introduire et commenter correctement au moins deux citations philosophiques ;
\item rédiger une conclusion qui répond à la problématique et ouvre sur une question nouvelle ;
\item évaluer une copie à l'aide de la grille de correction du Baccalauréat technique.
\end{itemize}

\courssub{Situation}

Dans le cadre de l'émission \emph{Cameroon Feeling} diffusée sur la CRTV, un reportage consacré aux mariages traditionnels dans la région de l'Ouest et dans la ville de Douala a relancé un vif débat national. Le reportage présente les données suivantes, recueillies auprès des familles et des notables :

\begin{center}
\begin{tabularx}{\linewidth}{|l|X|X|}\hline
\rowcolor{popblueL} \textbf{Localité} & \textbf{Montant moyen demandé (dot + liste)} & \textbf{Évolution sur 10 ans} \\ \hline
Bafoussam (Ouest) & \num{800000} FCFA & $+60\,\%$ \\ \hline
Douala (Littoral) & \num{1200000} FCFA & $+85\,\%$ \\ \hline
Bertoua (Est) & \num{350000} FCFA & $+25\,\%$ \\ \hline
\end{tabularx}
\end{center}

Le reportage donne la parole à deux témoins aux positions contradictoires :

\textbf{Témoignage 1 — Sa Majesté N., chef traditionnel :} « La dot n'est pas un prix d'achat. C'est un pacte entre deux familles, le signe que l'homme est capable d'assumer une épouse. Supprimer la dot, c'est détruire le lien social qui unit nos communautés depuis nos ancêtres. »

\textbf{Témoignage 2 — Aïcha, 23 ans, étudiante à l'Université de Douala :} « Quand on exige plus d'un million de francs, on transforme la femme en marchandise. Beaucoup de jeunes renoncent au mariage ou s'endettent. La femme devient un bien que l'on acquiert, et certains maris le lui rappellent ensuite. »

À la suite de ce reportage, ton professeur de philosophie te propose le sujet de dissertation suivant :

\begin{center}
\textbf{\emph{« La dot est-elle une tradition à préserver ou une marchandisation de la femme ? »}}
\end{center}

Tu disposes de 4 heures, en conditions d'examen, pour rédiger une dissertation complète.

\courssub{Consignes}

\textbf{Travail en groupes (étapes 1 et 2) :}
\begin{enumerate}
\item Souligne les termes clés du sujet et définis chacun d'eux (\emph{dot}, \emph{tradition}, \emph{préserver}, \emph{marchandisation}, \emph{femme}).
\item Dégage la tension contenue dans le sujet : quelles sont les deux positions qui s'affrontent ? À quels témoins du reportage correspondent-elles ?
\item Formule la problématique sous forme d'une question unique et précise.
\item Rédige en groupe l'introduction complète (accroche, analyse des termes, problématique, annonce du plan en trois parties).
\item Élabore le plan détaillé : pour chaque partie (thèse, antithèse, synthèse), indique l'idée directrice, deux arguments, un exemple camerounais tiré de la situation et une référence philosophique.
\end{enumerate}

\textbf{Travail individuel (étape 3) :}
\begin{enumerate}
\setcounter{enumi}{5}
\item Rédige l'une des trois parties du développement, en intégrant au moins une citation correctement introduite (auteur annoncé, citation entre guillemets, commentaire qui la relie à l'argument).
\item Rédige la conclusion : réponse claire à la problématique, bilan des trois parties, ouverture.
\end{enumerate}

\textbf{Correction croisée (étape 4) :}
\begin{enumerate}
\setcounter{enumi}{7}
\item Échange ta copie avec un autre groupe et évalue-la selon la grille : pertinence de la problématique (2 pts), structure en trois parties (4 pts), qualité des arguments et exemples (6 pts), citations (4 pts), langue et présentation (4 pts).
\item Participe à la mise en commun de la copie modèle au tableau.
\end{enumerate}

\courssub{Ressources à mobiliser}

\begin{aretenir}[Structure canonique de la dissertation dialectique]
\begin{itemize}
\item \textbf{Introduction} : accroche (fait, situation, citation) $\rightarrow$ analyse des termes du sujet $\rightarrow$ problématique (une question) $\rightarrow$ annonce du plan.
\item \textbf{Développement} : I. Thèse (la dot est une tradition à préserver) ; II. Antithèse (la dot est une marchandisation de la femme) ; III. Synthèse (dépassement : une tradition à réformer, non à abolir).
\item \textbf{Conclusion} : réponse à la problématique, bilan, ouverture.
\item \textbf{Citations} : toujours introduites (nom de l'auteur, contexte), placées entre guillemets, puis commentées et reliées à l'argument.
\end{itemize}
Références utiles : Kant (la personne comme fin, jamais comme simple moyen) ; Hegel (la dialectique thèse--antithèse--synthèse) ; Lévi-Strauss (le mariage comme échange et alliance) ; Simone de Beauvoir (la condition de la femme).
\end{aretenir}

\begin{methode}[Les 4 étapes de la rédaction en examen]
\begin{enumerate}
\item \textbf{Analyse et brouillon (1h) :} Dégager les termes, la tension, la problématique ; construire le plan détaillé (idée directrice, arguments, exemples, citations) pour les 3 parties.
\item \textbf{Introduction (30 min) :} La rédiger au net soigneusement ; elle conditionne la note.
\item \textbf{Développement (2h) :} Rédiger les trois parties en respectant la structure : idée directrice $\rightarrow$ argument 1 (explication + exemple + citation) $\rightarrow$ argument 2 (explication + exemple) $\rightarrow$ transition.
\item \textbf{Conclusion et relecture (30 min) :} Répondre à la problématique, faire le bilan, ouvrir. Relire pour corriger la langue et vérifier les citations.
\end{enumerate}
\end{methode}

\begin{savaistu}[La dot au Cameroun : diversité culturelle et enjeux contemporains]
Le Cameroun compte plus de 250 groupes ethniques aux coutumes matrimoniales variées. Chez les Bamiléké (Ouest), la dot (\emph{tsi} ou \emph{ndop}) est souvent élevée, symbolisant la valeur de la femme et l'alliance entre clans puissants. Chez les Bassa (Littoral/Centre) ou les Gbaya (Est), elle inclut traditionnellement des biens en nature (tissus, colas, bétail, vivres) avant la monétarisation récente. L'inflation des montants (jusqu'à \num{1200000} FCFA à Douala selon le reportage) pose le problème de l'accessibilité au mariage pour la jeunesse et de la dette conjugale. Le Code civil camerounais (Loi n°81-02 du 29 juin 1981) reconnaît le mariage coutumier mais ne fixe pas le montant de la dot, laissant place aux négociations familiales. La philosophie permet d'interroger la norme : la coutume doit-elle s'adapter aux droits humains (dignité, non-discrimination) ou la modernité doit-elle respecter l'intangible culturel ?
\end{savaistu}

\courssub{Démarche et corrigé commenté}

\begin{exoresolu}[De la dot à la dissertation : copie modèle commentée]
Rédiger une dissertation complète sur le sujet : \emph{« La dot est-elle une tradition à préserver ou une marchandisation de la femme ? »}
\tcblower
\textbf{Étape 1 — Analyse du sujet.} Les termes clés sont : \emph{dot} (biens remis par la famille du marié à celle de la mariée lors du mariage coutumier), \emph{tradition} (usage transmis de génération en génération), \emph{préserver} (conserver, maintenir), \emph{marchandisation} (transformation d'une personne ou d'une valeur en bien échangeable contre de l'argent), \emph{femme} (personne humaine de sexe féminin, sujet de droit). La tension oppose la \emph{valeur symbolique et sociale} de la dot à son \emph{dévoiement économique} possible. D'où la problématique : \textbf{la dot, en tant qu'institution coutumière, honore-t-elle la femme en l'intégrant dans une alliance entre familles, ou la réduit-elle au statut de marchandise ?}

\textbf{Étape 2 — Introduction (copie modèle).} Au Cameroun, le reportage de la CRTV sur les mariages coutumiers a montré que le montant de la dot atteint parfois \num{1200000} FCFA à Douala, suscitant un débat entre les défenseurs de la coutume, comme le chef traditionnel N., et ses critiques, comme l'étudiante Aïcha. La dot désigne les biens remis par le prétendant à la famille de sa future épouse ; la tradition est un héritage culturel transmis ; marchandiser, c'est traiter comme un objet de commerce. La question est alors de savoir si la dot est une tradition à préserver ou une marchandisation de la femme. Nous verrons d'abord que la dot est une institution symbolique fondatrice du lien social, ensuite que ses dérives financières en font une véritable marchandisation, enfin qu'il convient de la réformer pour concilier culture et dignité.

\textbf{Étape 3 — Développement (extraits commentés).}

\emph{I. Thèse.} La dot est d'abord un pacte symbolique entre deux familles. Comme l'écrit Lévi-Strauss dans \emph{Les Structures élémentaires de la parenté}, « le mariage est un échange » qui fonde l'alliance entre les groupes humains : cette citation montre que la dot n'a pas pour fonction d'acheter la femme, mais de sceller une relation sociale durable. Le témoignage du chef N. l'illustre : la dot atteste la capacité de l'homme à assumer ses responsabilités. À Bertoua, où le montant reste modéré (\num{350000} FCFA), la dot conserve cette fonction symbolique.

\emph{II. Antithèse.} Cependant, l'inflation des montants ($+85\,\%$ en dix ans à Douala) transforme la dot en transaction commerciale. Kant rappelle dans la \emph{Fondation de la métaphysique des mœurs} qu'il faut traiter l'humanité « toujours en même temps comme une fin, et jamais simplement comme un moyen » : or, exiger plus d'un million de francs, c'est faire de la femme un moyen d'enrichissement, ce qui viole sa dignité. Simone de Beauvoir montre en outre que la femme est souvent ramenée au statut d'objet dans les institutions patriarcales. Le témoignage d'Aïcha confirme cette dérive : endettement des jeunes mariés, sentiment de propriété chez certains époux.

\emph{III. Synthèse.} Il faut dépasser l'alternative abolition/conservation. La dot possède une valeur culturelle réelle, mais sa forme actuelle doit être régulée : plafonnement symbolique des montants, retour aux dons non monétaires (colas, tissus, vivres), sensibilisation des notables. La dialectique hégélienne enseigne que la vérité naît du dépassement des contraires : une tradition vivante est une tradition capable de se réformer.

\textbf{Étape 4 — Conclusion (copie modèle).} La dot n'est donc ni purement une tradition à préserver telle quelle, ni essentiellement une marchandisation de la femme : elle est une institution symbolique dont les dérives économiques contemporaines menacent la dignité de la personne. La préserver suppose de la réformer. Reste alors une question plus large : toute tradition doit-elle être jugée à l'aune de la dignité universelle de la personne humaine ?

\textbf{Justification méthodologique :} chaque exigence du programme est respectée — termes définis, problématique unique, trois parties dialectiques équilibrées, deux citations introduites et commentées (Lévi-Strauss, Kant), exemples chiffrés tirés de la situation, conclusion répondant à la problématique avec ouverture.
\end{exoresolu}

\begin{attention}[Erreurs fréquentes à éviter]
\begin{itemize}
\item Citer un auteur sans introduire ni commenter la citation (citation « parachutée »).
\item Traiter le sujet comme un débat de société sans définitions philosophiques des termes.
\item Déséquilibrer les parties : la thèse et l'antithèse doivent avoir un poids comparable.
\item Conclure sans répondre explicitement à la problématique posée.
\item Recopier les témoignages du document au lieu de les exploiter comme exemples au service d'un argument.
\end{itemize}
\end{attention}

\courssub{Bilan de l'activité}

Cette activité a permis de mobiliser, dans une situation d'examen réaliste, l'ensemble des savoir-faire de la leçon : analyse des termes et formulation d'une problématique à partir d'un document authentique, rédaction d'une introduction complète, construction d'un développement dialectique équilibré, usage rigoureux des citations et rédaction d'une conclusion. Elle a en outre montré que la philosophie n'est pas coupée de la vie courante : un débat social camerounais actuel, celui de la dot, peut être traité avec les outils de la réflexion philosophique (dignité kantienne, échange lévi-straussien, dialectique hégélienne). Enfin, la correction croisée à l'aide de la grille du Baccalauréat a entraîné les élèves à s'auto-évaluer, compétence décisive pour le jour de l'examen : celui qui sait corriger une copie sait aussi éviter les erreurs dans la sienne.

\newpage

\coursec{Méthodes et exercices résolus}

\coursec{Consolidation des acquis de la méthodologie de la dissertation philosophique}

\courssub{Méthode 1 : Rédiger une introduction complète}

\begin{methode}[Les 5 étapes de l'introduction philosophique]
Une introduction réussie au Probatoire / Baccalauréat technique obéit à un plan rigoureux en cinq moments :
\begin{enumerate}
\item \textbf{L'accroche (ou amorce)} : une phrase générale qui situe le sujet dans l'expérience humaine ou dans l'histoire de la pensée. Exemple : « De tout temps, l'homme s'est interrogé sur... »
\item \textbf{La présentation du sujet} : reformuler le sujet avec ses propres mots, en définissant les \cle{termes clés} (notions principales).
\item \textbf{L'analyse du sujet} : dégager la tension ou le paradoxe contenu dans le sujet (ce qui fait \emph{problème}).
\item \textbf{La problématique} : formuler une \cle{question principale} claire, qui commande tout le devoir. Elle commence souvent par « Dans quelle mesure... », « Faut-il... », « Peut-on... ».
\item \textbf{L'annonce du plan} : présenter brièvement les trois moments du développement (thèse, antithèse, synthèse), sans les détailler.
\end{enumerate}
\textbf{Réflexes à retenir} : l'introduction fait environ 10 à 15 lignes ; elle ne contient \emph{aucune} réponse personnelle ; la problématique est une question, jamais une affirmation.
\end{methode}

\begin{exoresolu}[Rédiger une introduction]
\textbf{Énoncé (type Bac)} : Sujet : « La technique est-elle un danger pour l'homme ? » Rédigez une introduction complète à ce sujet de dissertation.
\tcblower
\textbf{Solution détaillée.}

\textbf{Étape 1 — L'accroche.} On part d'un constat général et indiscutable : « Depuis la révolution industrielle, et plus encore à l'ère du numérique, la technique occupe une place centrale dans l'existence humaine : elle transforme le travail, la communication, la santé et jusqu'à notre rapport au monde. »

\textbf{Étape 2 — Présentation du sujet.} On reformule et on définit les termes clés : « Le sujet nous invite à interroger le rapport entre la \emph{technique}, c'est-à-dire l'ensemble des procédés et des outils que l'homme fabrique pour agir sur la nature, et le \emph{danger}, c'est-à-dire ce qui menace l'intégrité ou la liberté de l'homme. »

\textbf{Étape 3 — Analyse du sujet (le paradoxe).} « Or la technique apparaît d'abord comme un bienfait : elle libère l'homme de la fatigue, soigne, rapproche les hommes — au Cameroun, le téléphone mobile permet au commerçant de Douala de gérer ses transactions à distance. Mais cette même technique peut asservir : dépendance aux écrans, destruction des emplois, menaces sur l'environnement. La technique serait donc à la fois libération et menace. »

\textbf{Étape 4 — La problématique.} « Dès lors, une question s'impose : la technique est-elle un danger pour l'homme, ou le danger vient-il de l'usage que l'homme en fait ? »

\textbf{Étape 5 — Annonce du plan.} « Pour répondre à cette question, nous montrerons d'abord que la technique peut constituer un danger réel pour l'homme ; nous verrons ensuite qu'elle est d'abord un instrument de progrès et de liberté ; enfin, nous soutiendrons que le danger ne réside pas dans la technique elle-même, mais dans l'usage que l'homme en fait. »

\textbf{Conclusion de l'exercice} : l'introduction ainsi construite respecte les cinq étapes exigées (accroche, présentation, analyse, problématique, annonce du plan) et annonce clairement un plan dialectique en trois parties.
\end{exoresolu}

\courssub{Méthode 2 : Construire la thèse (première partie du développement)}

\begin{methode}[Rédiger la thèse en 4 étapes]
La \cle{thèse} est la première réponse apportée à la problématique : c'est la position la plus immédiate, souvent celle du sens commun.
\begin{enumerate}
\item \textbf{Formuler la position} en une phrase claire (le « point de vue » défendu dans cette partie).
\item \textbf{Définir les notions} mobilisées par cette position.
\item \textbf{Argumenter} : développer au moins deux arguments logiques, chacun illustré par un \cle{exemple concret} (vie quotidienne, réalité camerounaise, histoire).
\item \textbf{Appuyer sur un auteur} : citer brièvement un philosophe dont la doctrine soutient cette position (référence exacte, sans inventer).
\end{enumerate}
\textbf{Réflexes} : annoncer la partie par une phrase de transition (« Dans un premier temps, nous montrerons que... ») ; un argument = une idée + sa justification + un exemple ; ne jamais accumuler les exemples sans analyse.
\end{methode}

\begin{exoresolu}[Rédiger la thèse]
\textbf{Énoncé (type Bac)} : Sujet : « La technique est-elle un danger pour l'homme ? » Rédigez la première partie (thèse) du développement.
\tcblower
\textbf{Solution détaillée.}

\textbf{Étape 1 — Formulation de la position.} « Dans un premier temps, il semble que la technique constitue un danger réel pour l'homme, car elle risque de le déposséder de sa liberté et de détruire les conditions mêmes de son existence. »

\textbf{Étape 2 — Définition.} « Par danger, on entend ici tout ce qui menace l'homme dans son corps, sa liberté ou son humanité. »

\textbf{Étape 3 — Argumentation.}
\emph{Argument 1 : la technique asservit celui qu'elle devait servir.} « L'homme devient dépendant de ses propres outils : incapable de se passer de son téléphone, de l'électricité, des réseaux sociaux, il perd son autonomie. Ainsi, à Yaoundé comme dans les villages, nombre de jeunes passent plusieurs heures par jour sur les écrans, au détriment des études et des relations réelles : l'outil commande au lieu d'obéir. »
\emph{Argument 2 : la technique menace la nature et la survie humaine.} « La déforestation industrielle, la pollution des fleuves par les rejets des usines, le réchauffement climatique montrent que la technique, livrée à elle-même, détruit le milieu naturel dont dépend la vie humaine. »

\textbf{Étape 4 — Appui sur un auteur.} « C'est ce que soulignait le philosophe Martin Heidegger : la technique moderne n'est pas un simple outil, elle est une manière de « dévoiler » la nature comme un simple fonds exploitable, ce qui fait oublier à l'homme sa véritable essence. »

\textbf{Conclusion de l'exercice} : la thèse est construite selon le schéma exigé : position claire, définition, deux arguments illustrés, un appui philosophique. Elle répond à la problématique par l'affirmative, ce qui appelle logiquement une antithèse.
\end{exoresolu}

\courssub{Méthode 3 : Construire l'antithèse (deuxième partie)}

\begin{methode}[Rédiger l'antithèse en 4 étapes]
L'\cle{antithèse} oppose à la thèse une position contraire, tout aussi bien argumentée. Elle évite le devoir « à sens unique ».
\begin{enumerate}
\item \textbf{Transition critique} : montrer la limite ou l'insuffisance de la thèse (« Cependant, cette position est-elle suffisante ? »).
\item \textbf{Formuler la position contraire} en une phrase claire.
\item \textbf{Argumenter} avec au moins deux arguments, chacun illustré par un exemple concret.
\item \textbf{Appuyer sur un auteur} dont la doctrine défend cette position opposée.
\end{enumerate}
\textbf{Réflexes} : l'antithèse n'est pas une simple négation ; elle apporte des arguments \emph{nouveaux} ; elle doit être aussi développée que la thèse ; la transition est obligatoire pour montrer la logique du devoir.
\end{methode}

\begin{exoresolu}[Rédiger l'antithèse]
\textbf{Énoncé (type Bac)} : Sujet : « La technique est-elle un danger pour l'homme ? » Rédigez la deuxième partie (antithèse) du développement.
\tcblower
\textbf{Solution détaillée.}

\textbf{Étape 1 — Transition critique.} « Cependant, accuser la technique de tous les maux, n'est-ce pas oublier qu'elle est d'abord une création de l'intelligence humaine et la condition du progrès ? »

\textbf{Étape 2 — Formulation de la position contraire.} « Dans un second temps, nous montrerons que la technique est avant tout un instrument de libération et de progrès pour l'homme. »

\textbf{Étape 3 — Argumentation.}
\emph{Argument 1 : la technique libère l'homme de la nécessité naturelle.} « Grâce aux machines, l'homme échappe aux travaux épuisants : le tracteur remplace la houe éreintante, et le paysan camerounais peut cultiver de plus grandes surfaces avec moins de fatigue. Le temps gagné peut être consacré à l'éducation et à la vie sociale. »
\emph{Argument 2 : la technique sauve des vies et rapproche les hommes.} « Les vaccins, les médicaments, les ambulances reculent la mort ; le téléphone mobile et l'Internet permettent à un étudiant de Bafoussam de suivre des cours à distance et à une famille de rester en contact avec ses proches de la diaspora. Un tel outil peut-il être appelé un danger ? »

\textbf{Étape 4 — Appui sur un auteur.} « Comme le rappelle Bergson, l'homme est \emph{homo faber} : un fabricant d'outils ; l'intelligence humaine est par nature technique, et la technique prolonge l'élan créateur de la vie. Renoncer à la technique serait donc renoncer à l'humanité elle-même. »

\textbf{Conclusion de l'exercice} : l'antithèse oppose à la thèse une position contraire argumentée (deux arguments nouveaux, exemples camerounais, appui sur Bergson), précédée d'une transition critique. Le devoir présente désormais deux positions opposées qu'il faudra dépasser dans la synthèse.
\end{exoresolu}

\courssub{Méthode 4 : Construire la synthèse (troisième partie)}

\begin{methode}[Rédiger la synthèse en 4 étapes]
La \cle{synthèse} dépasse l'opposition thèse/antithèse : elle n'est ni un compromis mou, ni une simple répétition.
\begin{enumerate}
\item \textbf{Bilan critique} : rappeler en une phrase ce que la thèse et l'antithèse ont chacune de vrai et de limité.
\item \textbf{Dégager le vrai problème} : montrer que l'opposition reposait sur une confusion ou un point de vue partiel (c'est le \cle{dépassement}).
\item \textbf{Formuler la position personnelle argumentée} : la réponse nuancée et supérieure, soutenue par un argument décisif et éventuellement un auteur.
\item \textbf{Ouvrir} éventuellement sur les conditions pratiques de cette position (éducation, loi, responsabilité).
\end{enumerate}
\textbf{Réflexes} : la synthèse doit apporter quelque chose de \emph{nouveau} ; éviter la formule paresseuse « la vérité est entre les deux » ; c'est la partie la plus importante du devoir car elle manifeste la pensée personnelle du candidat.
\end{methode}

\begin{exoresolu}[Rédiger la synthèse]
\textbf{Énoncé (type Bac)} : Sujet : « La technique est-elle un danger pour l'homme ? » Rédigez la troisième partie (synthèse) du développement.
\tcblower
\textbf{Solution détaillée.}

\textbf{Étape 1 — Bilan critique.} « La thèse a montré que la technique peut asservir et détruire ; l'antithèse a établi qu'elle libère et fait progresser. Mais ces deux positions opposées parlent en réalité de la même chose : un instrument entre les mains de l'homme. »

\textbf{Étape 2 — Le vrai problème dégagé.} « L'opposition entre les deux thèses repose sur une confusion : elle attribue à la technique une intention qu'elle ne possède pas. La technique est en elle-même \emph{neutre} : c'est un moyen, non une fin. Le couteau peut couper le pain ou blesser ; le réseau social peut informer ou manipuler. Le danger ne vient donc pas de la technique, mais de l'\emph{usage} que l'homme en fait. »

\textbf{Étape 3 — Position personnelle argumentée.} « Nous soutiendrons donc que la technique n'est dangereuse que pour un homme dépourvu de sagesse et de responsabilité. Comme l'écrivait déjà Platon, le savoir-faire sans la connaissance du bien est une puissance aveugle. Ce qui manque à l'homme technique, ce n'est pas la puissance, mais la \emph{maîtrise morale} de cette puissance. »

\textbf{Étape 4 — Ouverture pratique.} « Dès lors, la réponse au danger technique n'est pas le refus de la technique, mais l'éducation : former des utilisateurs responsables, doter les États de lois protégeant l'environnement et les personnes. C'est à ce prix que la technique restera ce qu'elle doit être : une servante, et non un maître. »

\textbf{Conclusion de l'exercice} : la synthèse dépasse réellement l'opposition en déplaçant le problème (de la technique vers l'usage), formule une position personnelle appuyée sur Platon, et ouvre sur des conditions pratiques. Le plan dialectique est complet.
\end{exoresolu}

\courssub{Méthode 5 : Rédiger une conclusion efficace}

\begin{methode}[Les 3 étapes de la conclusion]
La \cle{conclusion} clôt le devoir ; elle ne doit rien introduire de nouveau.
\begin{enumerate}
\item \textbf{Rappeler la problématique} en la reformulant brièvement (« Nous nous demandions... »).
\item \textbf{Bilan du parcours} : résumer en deux ou trois phrases les trois moments du développement (thèse, antithèse, synthèse).
\item \textbf{Réponse claire à la problématique} : affirmer nettement la position défendue, puis \textbf{ouverture} : une question nouvelle, en lien avec le sujet, qui prolonge la réflexion.
\end{enumerate}
\textbf{Réflexes} : la conclusion fait environ 6 à 10 lignes ; interdiction d'ajouter un argument nouveau ; l'ouverture est une \emph{question}, jamais une affirmation gratuite ; éviter les formules vides (« en conclusion, je dirai que tout dépend... »).
\end{methode}

\begin{exoresolu}[Rédiger une conclusion]
\textbf{Énoncé (type Bac)} : Sujet : « La technique est-elle un danger pour l'homme ? » Rédigez la conclusion du devoir.
\tcblower
\textbf{Solution détaillée.}

\textbf{Étape 1 — Rappel de la problématique.} « Au terme de cette réflexion, revenons à notre question initiale : la technique est-elle un danger pour l'homme ? »

\textbf{Étape 2 — Bilan du parcours.} « Nous avons d'abord reconnu que la technique peut asservir l'homme et détruire la nature, ce qui la fait apparaître comme un danger réel. Nous avons ensuite montré qu'elle est d'abord un instrument de libération et de progrès, prolongeant l'intelligence humaine. Enfin, nous avons établi que la technique, étant un moyen neutre, ne saurait être dangereuse par elle-même. »

\textbf{Étape 3 — Réponse et ouverture.} « Nous répondrons donc que la technique n'est un danger que pour l'homme irresponsable : le véritable danger réside dans l'usage, non dans l'outil. Seule l'éducation morale et la responsabilité collective peuvent garantir que la puissance technique reste au service de l'homme. Mais une question demeure, qui prolonge notre réflexion : l'homme d'aujourd'hui est-il capable de la sagesse que réclame sa puissance ? »

\textbf{Conclusion de l'exercice} : cette conclusion respecte les trois étapes exigées (rappel de la problématique, bilan, réponse argumentée) et s'achève sur une ouverture pertinente, sans introduire aucun argument nouveau.
\end{exoresolu}

\courssub{Méthode 6 : Choisir et utiliser les citations}

\begin{methode}[Règles d'usage des citations en 5 points]
La \cle{citation} est un appui d'autorité : elle renforce un argument, elle ne le remplace jamais.
\begin{enumerate}
\item \textbf{Pertinence} : ne citer que des auteurs dont la doctrine correspond réellement au point défendu ; une citation hors sujet est sanctionnée.
\item \textbf{Exactitude} : ne jamais inventer une citation ni l'attribuer au mauvais auteur ; si la formulation exacte est incertaine, paraphraser (« selon Kant... ») plutôt que mettre entre guillemets.
\item \textbf{Brièveté} : une citation courte (une phrase au maximum), introduite par une formule (« Comme le dit Platon... », « C'est ce que souligne Descartes lorsqu'il écrit... »).
\item \textbf{Commentaire} : toute citation doit être \emph{expliquée} en une ou deux phrases : dire ce qu'elle signifie et en quoi elle appuie l'argument.
\item \textbf{Sobriété} : une à deux références par partie suffisent ; un devoir n'est pas un catalogue de citations.
\end{enumerate}
\textbf{Réflexes} : guillemets pour les mots exacts de l'auteur ; nom de l'auteur correctement orthographié ; la citation s'intègre grammaticalement à la phrase.
\end{methode}

\begin{exoresolu}[Bien utiliser une citation]
\textbf{Énoncé (type Bac)} : Dans un devoir sur le sujet « La liberté consiste-t-elle à faire ce que l'on veut ? », un élève écrit : « Comme le dit Rousseau, la liberté c'est faire ce qu'on veut, donc l'homme libre obéit à ses désirs. » Cette utilisation de la citation est-elle correcte ? Corrigez-la en rédigeant le passage attendu.
\tcblower
\textbf{Solution détaillée.}

\textbf{Étape 1 — Diagnostic de l'erreur.} La citation est \emph{fausse et inversée} : Rousseau soutient précisément le contraire. Dans le \emph{Contrat social}, il écrit que « l'obéissance à la loi qu'on s'est prescrite est liberté ». La liberté, pour Rousseau, n'est pas la satisfaction des désirs, mais l'obéissance à une loi que l'on se donne à soi-même (l'autonomie). Attribuer à Rousseau la thèse adverse détruit la crédibilité du devoir : c'est une erreur d'exactitude (règle 2) et de pertinence (règle 1).

\textbf{Étape 2 — Passage corrigé et rédigé.} « Faire ce que l'on veut, c'est-à-dire obéir à chacun de ses désirs, n'est pas la liberté mais l'esclavage des passions. C'est ce que montre Rousseau lorsqu'il écrit dans le \emph{Contrat social} que « l'obéissance à la loi qu'on s'est prescrite est liberté ». Par cette formule, Rousseau signifie que l'homme n'est véritablement libre que lorsqu'il suit une règle qu'il s'est lui-même donnée, par sa raison, et non lorsqu'il subit l'emprise de ses impulsions. Ainsi, le jeune qui sait refuser une distraction pour préparer son examen est plus libre que celui qui cède à tous ses désirs : le premier se gouverne, le second est gouverné. »

\textbf{Étape 3 — Vérification des règles.} La citation est exacte et correctement attribuée (règle 2), pertinente car elle soutient l'argument (règle 1), brève et introduite (règle 3), suivie d'un commentaire explicatif et d'un exemple (règle 4), et unique dans ce passage (règle 5).

\textbf{Conclusion de l'exercice} : une citation n'a de valeur que si elle est exacte, pertinente et expliquée ; mieux vaut paraphraser un auteur que l'on connaît bien que d'inventer une formule qui sera sanctionnée par le correcteur.
\end{exoresolu}

\courssub{Méthode 7 : Appliquer la dissertation à une situation concrète de la vie courante}

\begin{methode}[De la vie quotidienne à la dissertation en 4 étapes]
Au Probatoire / Baccalauréat technique, le sujet part souvent d'une situation concrète (conflit, rumeur, usage du téléphone, examen truqué...).
\begin{enumerate}
\item \textbf{Décrire la situation} et repérer le \cle{problème} qu'elle pose (un conflit de valeurs, un doute, une injustice).
\item \textbf{Dégager la notion philosophique} concernée (liberté, devoir, vérité, justice, technique, État...) et formuler la problématique.
\item \textbf{Traiter philosophiquement} : appliquer le plan dialectique (thèse, antithèse, synthèse) avec arguments et auteurs, en utilisant la situation comme \emph{exemple}, non comme sujet.
\item \textbf{Revenir à la situation} dans la synthèse ou la conclusion : tirer une leçon pratique, une conduite à tenir, une décision justifiée.
\end{enumerate}
\textbf{Réflexes} : la situation concrète est le point de départ et le point d'arrivée, mais le cœur du devoir reste l'argumentation philosophique ; toujours justifier sa réponse finale.
\end{methode}

\begin{exoresolu}[Traiter une situation concrète]
\textbf{Énoncé (type Bac)} : « Un élève de Terminale technique découvre, la veille du Baccalauréat, qu'un sujet d'examen circule sur un groupe WhatsApp. Ses camarades hésitent : faut-il l'utiliser ? » Dégager la problématique philosophique de cette situation et indiquer, en les justifiant, les grandes lignes du devoir à rédiger.
\tcblower
\textbf{Solution détaillée.}

\textbf{Étape 1 — Description et problème.} La situation met en conflit deux valeurs : l'\emph{intérêt personnel} (réussir facilement) et l'\emph{honnêteté} (respecter les règles communes de l'examen). Le problème est donc moral : faut-il faire son devoir même lorsque la tricherie est facile et impunie ?

\textbf{Étape 2 — Notion et problématique.} La notion philosophique mobilisée est le \cle{devoir} (la morale). Problématique : « Faut-il accomplir son devoir par intérêt ou par respect de la morale ? Autrement dit, l'honnêteté a-t-elle besoin d'une récompense pour être une valeur ? »

\textbf{Étape 3 — Grandes lignes du devoir (plan dialectique).}
\emph{Thèse} : on peut soutenir que l'élève a intérêt à utiliser le sujet : la réussite est le but, et « la fin justifie les moyens » (position machiavélienne) ; argument : dans une société où la réussite sociale prime, tricher semble rationnel.
\emph{Antithèse} : mais tricher, c'est se dérober au devoir et traiter autrui injustement : Kant montre que la moralité consiste à agir par devoir, selon une maxime universalisable — or si tous trichaient, l'examen ne vaudrait plus rien ; la tricherie détruit donc la condition même du diplôme convoité.
\emph{Synthèse} : la vraie réussite est celle qui a une valeur ; l'honnêteté n'est pas une contrainte extérieure mais la condition de la dignité et de la confiance sociale. Le devoir accompli par intérêt cesse d'être moral ; il faut agir par respect de la loi morale.

\textbf{Étape 4 — Retour à la situation.} Conclusion pratique justifiée : l'élève doit refuser le sujet divulgué et, si possible, signaler la fuite, car son diplôme ne certifiera sa compétence que s'il est obtenu honnêtement ; la réussite volée est une réussite sans valeur, pour lui-même comme pour la société camerounaise qui a besoin de techniciens réellement compétents.

\textbf{Conclusion de l'exercice} : la situation concrète a été convertie en problématique philosophique (le devoir), traitée selon le plan dialectique avec appui sur Kant, puis résolue par une conduite justifiée : c'est exactement la démarche attendue au Probatoire / Baccalauréat technique.
\end{exoresolu}

\begin{attention}[Les 5 erreurs qui coûtent des points au Bac]
\begin{enumerate}
\item \textbf{Le hors-sujet} : réciter un cours sur la notion sans répondre à la question précise posée. Réflexe : souligner les termes du sujet, les définir, et vérifier à chaque partie que l'on répond à la problématique.
\item \textbf{L'introduction sans problématique} (ou avec une problématique qui n'est pas une question) : sans question directrice, le devoir est sans fil conducteur et lourdement sanctionné. La problématique doit être une phrase interrogative.
\item \textbf{La synthèse-compromis} : écrire « la vérité est entre les deux » ou « tout dépend des points de vue ». La synthèse doit \emph{dépasser} l'opposition en apportant une position personnelle argumentée et nouvelle.
\item \textbf{Les citations inventées ou mal attribuées} : faire dire à Platon ce qu'a dit Kant, ou fabriquer une formule entre guillemets, décrédibilise tout le devoir. En cas de doute, paraphraser (« selon Kant... ») sans guillemets.
\item \textbf{Les exemples sans analyse et les arguments sans exemple} : aligner des faits (« au Cameroun... ») sans montrer ce qu'ils prouvent, ou affirmer sans illustrer. Règle d'or : un argument = une idée + une justification + un exemple + une phrase d'analyse qui relie l'exemple à l'idée.
\end{enumerate}
\end{attention}

\newpage

\coursec{Exercices}

\coursec{Consolidation des acquis de la méthodologie de la dissertation philosophique}

\courssub{Je vérifie mes connaissances}

\begin{exercice}[QCM : les exigences de la dissertation]
Pour chaque question, une seule réponse est correcte. Entoure la lettre correspondante et justifie brièvement ton choix.
\begin{enumerate}
\item La problématique d'une introduction philosophique se présente sous la forme :
\begin{enumerate}
\item d'une définition du dictionnaire ;
\item d'une question (ou d'un ensemble de questions) qui fait problème ;
\item d'une citation d'auteur ;
\item d'une opinion personnelle affirmée.
\end{enumerate}
\item Dans un développement dialectique, l'antithèse sert à :
\begin{enumerate}
\item répéter la thèse avec d'autres mots ;
\item contredire gratuitement l'auteur cité ;
\item montrer les limites de la thèse et faire apparaître la position opposée ;
\item conclure le devoir.
\end{enumerate}
\item Une citation est bien utilisée lorsqu'elle est :
\begin{enumerate}
\item copiée sans nom d'auteur ;
\item exacte, attribuée à son auteur et commentée ;
\item placée dans la conclusion ;
\item la plus longue possible.
\end{enumerate}
\item La conclusion d'une dissertation philosophique doit :
\begin{enumerate}
\item introduire un exemple nouveau ;
\item répondre clairement à la problématique et rappeler le chemin parcouru ;
\item poser une nouvelle problématique ;
\item rester volontairement vague.
\end{enumerate}
\end{enumerate}
\end{exercice}

\begin{exercice}[Vrai ou faux justifié]
Indique si chaque affirmation est vraie ou fausse, puis justifie ta réponse en deux ou trois phrases.
\begin{enumerate}
\item « L'accroche de l'introduction peut être une anecdote, un fait d'actualité ou un constat de la vie courante. »
\item « Il est interdit, dans une dissertation, de donner son avis personnel : il faut seulement réciter les auteurs. »
\item « Une citation approximative, dont on ne connaît pas l'auteur, peut quand même être utilisée entre guillemets. »
\item « La synthèse consiste toujours à dire que les deux positions précédentes ont raison. »
\item « Le hors-sujet est une erreur de contenu, jamais une erreur de méthode. »
\end{enumerate}
\end{exercice}

\begin{exercice}[Définitions express]
Définis en deux ou trois phrases chacun des termes suivants, tels qu'ils sont employés dans la méthodologie de la dissertation philosophique :
\begin{enumerate}
\item la \cle{problématique} ;
\item la \cle{thèse} ;
\item la \cle{transition} ;
\item le \cle{commentaire de citation}.
\end{enumerate}
\end{exercice}

\begin{exercice}[Remettre l'introduction en ordre]
Voici, dans le désordre, les cinq étapes d'une introduction rédigée sur le sujet « La technique libère-t-elle l'homme ? ». Remets-les dans l'ordre correct et nomme chaque étape (accroche, présentation du sujet, problématique, annonce du plan, définition des termes).
\begin{enumerate}
\item « Nous montrerons d'abord que la technique libère l'homme de la nature, puis qu'elle peut aussi l'asservir, avant de dégager les conditions d'un usage libre de la technique. »
\item « Au marché de Mokolo comme dans les ateliers de Douala, le téléphone portable est devenu un outil de travail quotidien : la technique est partout. »
\item « La technique désigne l'ensemble des procédés et des outils que l'homme fabrique pour transformer la nature ; la liberté est le pouvoir d'agir selon sa propre volonté. »
\item « Mais l'outil qui nous rend plus puissants ne risque-t-il pas de nous rendre dépendants ? La technique libère-t-elle l'homme ou l'asservit-elle ? »
\item « Le sujet porte donc sur le rapport entre l'homme et ses outils : ce rapport est-il un rapport de maîtrise ou de soumission ? »
\end{enumerate}
\end{exercice}

\courssub{Je m'entraîne}

\begin{exercice}[Rédiger une problématique]
Pour chacun des trois sujets suivants, définis les termes essentiels en une phrase, puis formule une problématique claire (une question principale, éventuellement suivie d'une question secondaire).
\begin{enumerate}
\item « Faut-il toujours dire la vérité ? »
\item « La culture nous éloigne-t-elle de la nature ? »
\item « Le bonheur est-il affaire de chance ? »
\end{enumerate}
\end{exercice}

\begin{exercice}[Construire une transition dialectique]
On te donne la fin d'une thèse et le début d'une antithèse sur le sujet « L'État peut-il imposer le civisme ? ».
\begin{itemize}
\item Fin de la thèse : « L'État, par ses lois et ses sanctions, rend donc possible la vie civique en obligeant chacun à respecter les règles communes. »
\item Début de l'antithèse : « Pourtant, un civisme obtenu par la seule contrainte n'est-il pas une simple obéissance, dépourvue de toute valeur morale ? »
\end{itemize}
\begin{enumerate}
\item Rédige la transition complète (trois à cinq phrases) qui fait passer de la thèse à l'antithèse en montrant la limite de la première position.
\item Indique quel mot ou quelle expression de ta transition signale le retournement dialectique.
\end{enumerate}
\end{exercice}

\begin{exercice}[Classer et commenter des citations]
Voici six citations. Classe-les dans le tableau selon la notion du programme à laquelle elles se rapportent (liberté, devoir, État, culture, bonheur, vérité), puis choisis-en une et rédige son commentaire en quatre ou cinq phrases (sens, enjeu, intérêt pour un sujet de dissertation).
\begin{enumerate}
\item « Agis de telle sorte que la maxime de ta volonté puisse valoir comme principe d'une législation universelle » (Kant).
\item « L'homme est né libre, et partout il est dans les fers » (Rousseau).
\item « Le bonheur est le but de la vie humaine » (Aristote).
\item « La culture est ce qui reste lorsqu'on a tout oublié » (Édouard Herriot).
\item « L'État, c'est moi » (attribué à Louis XIV).
\item « La vérité est fille du temps, non de l'autorité » (Francis Bacon).
\end{enumerate}
\begin{center}
\begin{tabularx}{\linewidth}{|l|X|}\hline
\rowcolor{popblueL} \textbf{Notion} & \textbf{Numéro(s) de la citation} \\ \hline
Liberté & \\ \hline
Devoir & \\ \hline
État & \\ \hline
Culture & \\ \hline
Bonheur & \\ \hline
Vérité & \\ \hline
\end{tabularx}
\end{center}
\end{exercice}

\begin{exercice}[Corriger une conclusion défaillante]
Un élève termine ainsi sa dissertation sur le sujet « Le travail est-il une malédiction ? » : « Voilà. On a vu plein de choses sur le travail. Chacun pense ce qu'il veut. Comme dit l'autre, tout travail mérite salaire. Et vous, qu'en pensez-vous ? »
\begin{enumerate}
\item Relève quatre erreurs méthodologiques commises dans cette conclusion.
\item Réécris une conclusion correcte (cinq à huit phrases) qui répond clairement à la question posée et rappelle les étapes du raisonnement.
\end{enumerate}
\end{exercice}

\courssub{Je me prépare au Bac}

\begin{exercice}[Sujet type Baccalauréat technique : introduction et plan détaillé]
Sujet : \emph{« La technique libère-t-elle l'homme ou l'asservit-elle ? »} (durée conseillée : 45 minutes).
\begin{enumerate}
\item Définis les deux termes essentiels du sujet (\cle{technique}, \cle{liberté}) en une phrase chacun.
\item Rédige l'introduction complète (accroche, présentation du sujet, définition des termes, problématique, annonce du plan), en douze à quinze lignes. Ton accroche devra partir d'une situation concrète de la vie camerounaise (par exemple l'usage du téléphone portable ou de la moto-taxi).
\item Élabore le plan détaillé du développement en trois parties (thèse, antithèse, synthèse) : pour chaque partie, indique l'idée directrice en une phrase, deux arguments et une référence (auteur ou exemple) que tu mobiliserais.
\item Rédige la phrase de transition qui conduit de la thèse à l'antithèse.
\end{enumerate}
\end{exercice}

\begin{exercice}[Repérage d'erreurs dans une copie d'élève]
Voici des extraits d'une copie d'élève sur le sujet « Faut-il toujours dire la vérité ? ». La copie contient \textbf{cinq fautes méthodologiques} : repère-les, nomme-les et propose pour chacune une correction rédigée.
\begin{enumerate}
\item « Depuis la nuit des temps, les hommes se demandent s'il faut dire la vérité. Je vais donc parler de la vérité. » (toute l'introduction)
\item « Comme le dit Platon dans \emph{Le Prince} : "Il faut toujours mentir pour plaire." »
\item « La vérité, c'est quand on dit ce qui est vrai. Le mensonge, c'est quand on ne dit pas la vérité. Donc il faut dire la vérité parce que c'est vrai. »
\item « Kant pense qu'il faut toujours dire la vérité. Moi aussi je pense comme Kant. D'ailleurs Kant a raison, donc j'ai raison. » (fin de la première partie, sans autre développement)
\item (La copie s'arrête après l'antithèse, sur la phrase :) « Voilà pourquoi certains préfèrent parfois le mensonge. »
\end{enumerate}
Questions :
\begin{enumerate}
\item Identifie et nomme précisément les cinq fautes (par exemple : citation inventée ou mal attribuée, hors-sujet, pétition de principe, absence de problématique, absence de conclusion, simple alignement d'opinions sans argumentation).
\item Rédige la problématique que cette copie aurait dû formuler.
\item Rédige la conclusion qui aurait dû terminer ce devoir, en répondant à la question posée et en intégrant correctement une référence à Kant.
\end{enumerate}
\end{exercice}

\begin{exercice}[Dissertation complète en conditions d'examen]
Sujet : \emph{« L'État peut-il imposer le civisme ? »} (durée : 3 heures, sans document).
\begin{enumerate}
\item Rédige la dissertation complète : introduction (avec problématique et annonce du plan), développement en trois parties (thèse, antithèse, synthèse) reliées par des transitions explicites, conclusion.
\item Ton développement devra mobiliser au moins : deux références philosophiques exactes et commentées (par exemple Hobbes, Rousseau ou Kant), un exemple tiré de la vie civique camerounaise (élections, propreté des cités, paiement des impôts, respect du code de la route) et un proverbe ou une sagesse africaine commenté.
\item Aucune citation ne devra être inventée : si tu n'es pas sûr de la formulation exacte, reformule la pensée de l'auteur sans guillemets en l'attribuant clairement.
\item Auto-évaluation finale : à l'aide de la grille ci-dessous, note ta propre copie sur 20 et rédige trois phrases de bilan (point fort, point faible, axe de progrès).
\end{enumerate}
\begin{center}
\begin{tabularx}{\linewidth}{|X|c|}\hline
\rowcolor{popblueL} \textbf{Critère} & \textbf{Barème} \\ \hline
Présentation générale, lisibilité, orthographe & /4 \\ \hline
Compréhension du sujet et pertinence (pas de hors-sujet) & /4 \\ \hline
Problématique et structure de l'introduction & /3 \\ \hline
Qualité de l'argumentation et des transitions & /4 \\ \hline
Exactitude et commentaire des références et citations & /3 \\ \hline
Conclusion : réponse claire à la problématique & /2 \\ \hline
\end{tabularx}
\end{center}
\end{exercice}

\newpage

\coursec{Corrigés des exercices}

\coursec{Rappel des principes et exigences de la dissertation philosophique}

\begin{definition}[La dissertation philosophique]
Exercice de réflexion argumentée qui répond à une question philosophique (le sujet) en suivant une démarche rigoureuse : analyser les termes, problématiser, construire un développement dialectique (thèse, antithèse, synthèse) et conclure en assumant une réponse nuancée. Elle mobilise des connaissances (auteurs, concepts) mais ne se réduit pas à une récitation de cours.
\end{definition}

\begin{methode}[Les quatre temps de la dissertation]
\begin{enumerate}
\item \textbf{Analyse du sujet (brouillon) :} identifier les termes clés, les oppositions, les présupposés ; formuler la problématique ; chercher arguments, contre-arguments, références, exemples concrets.
\item \textbf{L'introduction (5 étapes obligatoires) :} 
      \begin{enumerate}
      \item \textbf{Accroche (amorce) :} fait concret, actualité, anecdote, proverbe, situation camerounaise — brève et liée au sujet.
      \item \textbf{Présentation du sujet :} reformulation, mise en évidence de l'enjeu philosophique.
      \item \textbf{Définition des termes :} définitions philosophiques précises (pas du dictionnaire courant) des concepts clés.
      \item \textbf{Problématique :} question (ou ensemble de questions) qui exprime la tension entre les termes ; elle guide tout le devoir.
      \item \textbf{Annonce du plan :} les trois moments (thèse, antithèse, synthèse) annoncés avec marqueurs (« D'abord… Ensuite… Enfin… » ou « Nous montrerons que… Nous verrons que… Nous dégagerons… »).
      \end{enumerate}
\item \textbf{Le développement (structure dialectique) :}
      \begin{itemize}
      \item \textbf{Thèse :} première réponse argumentée à la problématique (idée directrice + arguments + références + exemples).
      \item \textbf{Transition thèse $\rightarrow$ antithèse :} bilan de la thèse, mise en évidence de sa limite, question qui lance l'antithèse (marqueur : « Cependant », « Pourtant », « Mais »).
      \item \textbf{Antithèse :} position opposée ou limite radicale de la thèse, argumentée de même façon.
      \item \textbf{Transition antithèse $\rightarrow$ synthèse :} bilan de l'opposition, besoin de dépassement.
      \item \textbf{Synthèse :} dépassement de l'opposition (pas un compromis mou) : montrer sous quelles conditions, à quel niveau, dans quelle mesure chaque position est valable ; apporter une réponse construite.
      \end{itemize}
\item \textbf{La conclusion (3 étapes) :}
      \begin{enumerate}
      \item \textbf{Bilan synthétique :} rappel en une ou deux phrases du chemin parcouru (thèse, antithèse, synthèse).
      \item \textbf{Réponse explicite à la problématique :} « Oui, mais… », « Non, sous condition… », « C'est… à condition que… ».
      \item \textbf{Ouverture (optionnelle, brève) :} prolongement vers un autre sujet, une question voisine, sans nouvel argument ni nouvelle citation.
      \end{enumerate}
\end{enumerate}
\end{methode}

\begin{propriete}[Règles d'or pour les citations]
Une citation n'a de valeur philosophique que si elle remplit \textbf{trois conditions} simultanées :
\begin{enumerate}
\item \textbf{Exactitude :} texte reproduit fidèlement (guillemets français « … »).
\item \textbf{Attribution :} auteur nommé avec précision (œuvre si possible).
\item \textbf{Commentaire :} explicitation immédiate du sens, de la portée et du lien avec l'argument en cours.
\end{enumerate}
\emph{Corollaire :} si l'on n'est pas sûr de l'exactitude ou de l'auteur, on reformule sans guillemets : « Selon Kant… », « On peut penser avec Rousseau que… ».
\end{propriete}

\begin{aretenir}[Vocabulaire méthodologique du Probatoire / Baccalauréat]
\begin{itemize}
\item \cle{Problématique} $\neq$ sujet $\neq$ plan. Elle naît de la \emph{tension} entre les termes.
\item \cle{Transition dialectique} = Bilan + Limite + Question lançant la partie suivante.
\item \cle{Synthèse} $\neq$ résumé $\neq$ compromis. Elle \emph{dépasse} l'opposition.
\item \cle{Hors-sujet} : erreur de méthode (mauvaise analyse, problématique absente, parties qui ne répondent pas à la question).
\item \cle{Pétition de principe} : supposer ce qu'il faut prouver (ex. : définir la vérité par « ce qui est vrai »).
\item \cle{Appel à l'autorité} : remplacer l'argument par « X a raison donc c'est vrai ».
\end{itemize}
\end{aretenir}

\begin{attention}[Pièges fréquents à l'examen]
\begin{itemize}
\item \textbf{Accroche clichée :} « Depuis la nuit des temps… », « De tout temps les hommes… » $\rightarrow$ sanctionnée.
\item \textbf{Définitions après la problématique :} l'ordre est immuable : Définitions $\rightarrow$ Problématique.
\item \textbf{Citation en conclusion :} interdite (aucune référence nouvelle).
\item \textbf{Relativisme de conclusion :} « Chacun pense ce qu'il veut », « Tout est relatif » $\rightarrow$ renoncement à philosopher.
\item \textbf{Copie « catalogue » :} alignement d'auteurs sans articulation dialectique ni prise de position personnelle.
\item \textbf{Oubli de la synthèse :} copie qui s'arrête à l'antithèse.
\end{itemize}
\end{attention}

\begin{exemplebox}[Exemple d'accroches concrètes camerounaises]
\begin{itemize}
\item Sujet « Technique / Liberté » : « À Douala, le conducteur de moto-taxi gagne sa vie grâce à sa machine ; le commerçant règle ses transactions par mobile money. »
\item Sujet « État / Civisme » : « À Yaoundé, les opérations "Ville propre" et les contraventions pour jet d'ordures montrent l'État au quotidien. »
\item Sujet « Travail / Malédiction » : « Dans les plantations de l'Ouest, l'agriculteur transforme la forêt en champ nourricier. »
\item Sujet « Culture / Nature » : « Les danses du Ngondo sur le Wouri illustrent comment la culture humanise l'espace naturel. »
\end{itemize}
\end{exemplebox}

\coursec{TP 1 : L'introduction}

\begin{corrige}[Exercice 1 — QCM : les exigences de la dissertation]
\begin{enumerate}
\item \textbf{Réponse b.} La problématique est une question (ou un ensemble de questions) qui fait problème : elle naît d'une tension entre les termes du sujet. Une définition de dictionnaire est plate et scolaire ; une citation ne peut pas tenir lieu de question ; une opinion affirmée court-circuite la recherche au lieu de l'ouvrir.
\item \textbf{Réponse c.} L'antithèse montre les limites de la thèse et fait apparaître la position opposée : c'est le moment du \emph{retournement dialectique}. Répéter la thèse (a) supprime toute progression ; contredire gratuitement (b) n'est pas argumenter ; conclure (d) est le rôle de la conclusion.
\item \textbf{Réponse b.} Une citation n'a de valeur que si elle est \cle{exacte}, \cle{attribuée} à son auteur et \cle{commentée} (expliquée et reliée à l'argumentation). Une citation sans auteur est invérifiable ; la conclusion n'est pas le lieu d'introduire des références nouvelles ; la longueur n'est pas un critère de qualité.
\item \textbf{Réponse b.} La conclusion répond clairement à la problématique et rappelle le chemin parcouru (bilan des trois parties). Elle ne doit introduire ni exemple nouveau (a), ni problématique nouvelle (c), ni rester vague (d) : elle tranche et assume une réponse nuancée.
\end{enumerate}
\textbf{Points de vigilance :} à l'examen, la justification demandée doit citer la règle méthodologique, pas seulement recopier la bonne réponse.
\end{corrige}

\begin{corrige}[Exercice 4 — Remettre l'introduction en ordre]
Ordre correct : \textbf{2 -- 5 -- 3 -- 4 -- 1}.
\begin{enumerate}
\item \textbf{Phrase 2 — Accroche} : constat concret de la vie courante (marché de Mokolo, ateliers de Douala) qui capte l'attention et amène le sujet.
\item \textbf{Phrase 5 — Présentation du sujet} : elle dégage l'enjeu général (le rapport de l'homme à ses outils) à partir de l'accroche.
\item \textbf{Phrase 3 — Définition des termes} : on définit \cle{technique} et \cle{liberté} avant de questionner, pour éviter tout malentendu.
\item \textbf{Phrase 4 — Problématique} : la tension entre maîtrise et dépendance est formulée en question(s).
\item \textbf{Phrase 1 — Annonce du plan} : les trois moments (libération, asservissement, conditions d'un usage libre) sont annoncés avec les marqueurs « d'abord... puis... avant de... ».
\end{enumerate}
\textbf{Points de vigilance :} la définition des termes précède toujours la problématique ; l'annonce du plan est toujours la dernière étape de l'introduction.
\end{corrige}

\begin{corrige}[Exercice 5 — Rédiger une problématique]
\begin{enumerate}
\item \textbf{« Faut-il toujours dire la vérité ? »} — Définitions : la \cle{vérité} est la conformité de ce qu'on dit à ce qui est ; « falloir » renvoie au \cle{devoir} moral. Problématique : le devoir de vérité admet-il des exceptions ? Dire la vérité est-il toujours un bien, ou peut-il parfois nuire à autrui ?
\item \textbf{« La culture nous éloigne-t-elle de la nature ? »} — Définitions : la \cle{culture} est l'ensemble des œuvres, savoirs et coutumes par lesquels l'homme transforme ce qui est donné ; la \cle{nature} est ce qui existe sans l'action humaine. Problématique : en transformant la nature, la culture nous en détache-t-elle, ou n'est-elle que l'expression de notre nature d'homme ?
\item \textbf{« Le bonheur est-il affaire de chance ? »} — Définitions : le \cle{bonheur} est un état durable de satisfaction ; la \cle{chance} est ce qui arrive sans que notre volonté y soit pour quelque chose. Problématique : le bonheur dépend-il des circonstances extérieures, ou dépend-il de nous, de notre sagesse et de nos efforts ?
\end{enumerate}
\textbf{Points de vigilance :} la problématique doit être une vraie question ouverte (pas une question de cours à réponse unique) et porter sur les termes du sujet, non sur un sujet voisin.
\end{corrige}

\coursec{TP 2 : Le développement — thèse, antithèse, synthèse}

\begin{corrige}[Exercice 2 — Vrai ou faux justifié]
\begin{enumerate}
\item \textbf{Vrai.} L'accroche (ou amorce) peut être une anecdote, un fait d'actualité, un constat de la vie courante ou une situation concrète. Elle doit cependant rester brève et conduire naturellement au sujet : une accroche sans lien avec le sujet est une erreur.
\item \textbf{Faux.} La dissertation exige une réflexion \emph{personnelle mais argumentée} : l'élève doit prendre position, notamment dans la synthèse et la conclusion. Les auteurs sont des appuis, non des substituts à la pensée : réciter sans discuter est la « dissertation-catalogue », faute classique.
\item \textbf{Faux.} Les guillemets engagent à l'exactitude : une citation approximative ou sans auteur connu ne doit pas être mise entre guillemets. La règle est alors de \emph{reformuler} la pensée de l'auteur sans guillemets, en l'attribuant clairement (« selon Kant... », « on peut penser avec Rousseau que... »).
\item \textbf{Faux.} La synthèse n'est pas un compromis mou (« tout le monde a raison ») : elle dépasse l'opposition en montrant sous quelles conditions, à quel niveau ou dans quelle mesure chaque position est valable. Elle apporte une réponse construite et nuancée à la problématique.
\item \textbf{Faux.} Le hors-sujet est d'abord une erreur de méthode : il naît d'une mauvaise analyse des termes du sujet, d'une problématique mal formulée ou de parties qui ne répondent pas à la question posée. Une copie peut être riche en contenu et totalement hors-sujet.
\end{enumerate}
\end{corrige}

\begin{corrige}[Exercice 3 — Définitions express]
\begin{enumerate}
\item \textbf{La problématique} : question (ou ensemble de questions) qui exprime le problème philosophique contenu dans le sujet. Elle naît de la tension entre les termes du sujet, oriente tout le devoir et doit recevoir une réponse explicite dans la conclusion.
\item \textbf{La thèse} : première position examinée dans le développement, qui répond à la problématique dans un sens déterminé. Elle doit être soutenue par des arguments rationnels et des références, et non simplement affirmée.
\item \textbf{La transition} : phrase ou court paragraphe qui relie deux parties du développement. Elle fait le bilan de la partie qui s'achève, en montre la limite, et annonce la partie suivante par une question : c'est le moteur de la progression dialectique.
\item \textbf{Le commentaire de citation} : explicitation qui suit obligatoirement une citation. Il en dégage le sens, montre l'enjeu de la pensée de l'auteur et relie la citation à l'argument en cours : une citation non commentée est une citation inutile.
\end{enumerate}
\end{corrige}

\begin{corrige}[Exercice 6 — Construire une transition dialectique]
\textbf{1. Transition rédigée (exemple) :} « Ainsi, sans la loi et la sanction, la vie en société serait exposée à la violence de chacun : l'État apparaît bien comme le garant du civisme. \emph{Pourtant}, cette obéissance extérieure suffit-elle à faire un citoyen ? Un conducteur de Yaoundé qui s'arrête au feu rouge uniquement par peur de l'amende agit-il moralement ? La contrainte produit des comportements réglés, mais non nécessairement des convictions : c'est pourquoi il faut examiner si le civisme véritable ne relève pas plutôt de l'éducation et de l'adhésion intérieure. »

\textbf{2. Marqueur du retournement :} le mot « \emph{Pourtant} » (on accepte aussi « mais », « cependant », « néanmoins ») suivi d'une question qui met en difficulté la thèse : c'est le signal du passage à l'antithèse.

\textbf{Points de vigilance :} une bonne transition comporte toujours trois moments : bilan de la thèse, limite montrée, question annonçant l'antithèse. Une simple juxtaposition (« Nous allons maintenant voir l'antithèse ») est une faute de méthode.
\end{corrige}

\begin{corrige}[Exercice 10 — Repérage d'erreurs dans une copie d'élève (focus développement)]
\textbf{1. Les cinq fautes :}
\begin{enumerate}
\item \textbf{Introduction vide et sans problématique :} « depuis la nuit des temps » est une accroche clichée ; « je vais parler de la vérité » annonce un exposé, pas un problème. Correction : formuler la tension du sujet (« le devoir de vérité peut-il entrer en conflit avec d'autres devoirs ? ») et annoncer un plan.
\item \textbf{Citation inventée et mal attribuée :} \emph{Le Prince} est de Machiavel, non de Platon, et la phrase « il faut toujours mentir pour plaire » est fabriquée. Correction : ne citer que du vérifié ; sinon reformuler sans guillemets (« on peut penser avec Platon que... »).
\item \textbf{Pétition de principe et définitions circulaires :} définir la vérité par « ce qui est vrai » ne définit rien ; conclure « il faut dire la vérité parce que c'est vrai » suppose déjà démontré ce qu'il fallait prouver. Correction : définir la vérité comme accord de la pensée avec la réalité, puis argumenter (valeur de la confiance, condition de la vie sociale).
\item \textbf{Alignement d'opinions sans argumentation :} « Kant a raison, donc j'ai raison » remplace la preuve par un appel à l'autorité. Correction : exposer l'argument de Kant (le mensonge détruit la confiance et ne peut être universalisé), le discuter, puis prendre position en justifiant.
\item \textbf{Absence de synthèse et de conclusion :} la copie s'arrête après l'antithèse, sans réponse à la question. Correction : rédiger une synthèse (la vérité comme règle, avec des cas limites discutés : le mensonge par humanité) et une conclusion qui tranche.
\end{enumerate}

\textbf{2. Problématique attendue :} « Dire la vérité est-il un devoir absolu, sans aucune exception, ou ce devoir peut-il céder devant d'autres exigences, comme la protection d'autrui ? »

\textbf{3. Conclusion rédigée (exemple) :} « Nous avons vu d'abord que la vérité est la condition de la confiance et de la vie sociale : Kant soutient à ce titre que le devoir de vérité ne souffre aucune exception, car le mensonge, universalisé, détruirait toute parole. Nous avons montré ensuite que ce rigorisme se heurte à des cas où la vérité tue : mentir pour protéger un innocent poursuivi semble moralement exigé. La synthèse consiste alors à distinguer la règle et l'usage qu'on en fait : la vérité demeure le principe, mais son application exige le discernement, c'est-à-dire la prudence morale. Faut-il donc toujours dire la vérité ? Oui, comme règle de la vie commune ; non, si l'on entend par là un devoir aveugle qui ignorerait la situation et la dignité des personnes. La vérité est un devoir, mais un devoir qui pense. »

\textbf{Barème indicatif :} repérage et nomination exacts des cinq fautes /10 ; problématique /3 ; conclusion /7.
\textbf{Points de vigilance :} nommer précisément chaque faute avec le vocabulaire méthodologique (pétition de principe, appel à l'autorité, citation inventée) ; la référence à Kant doit être exacte et commentée, pas plaquée.
\end{corrige}

\coursec{TP 3 : La conclusion}

\begin{corrige}[Exercice 8 — Corriger une conclusion défaillante]
\textbf{1. Quatre erreurs méthodologiques :}
\begin{itemize}
\item Registre de langue familier et oral (« Voilà », « plein de choses », « l'autre ») : une copie exige un français soutenu et précis.
\item Absence de réponse à la problématique : « chacun pense ce qu'il veut » est un relativisme qui renonce à trancher.
\item Aucun bilan du raisonnement : les étapes du développement ne sont pas rappelées.
\item Citation approximative et non attribuée (« comme dit l'autre ») introduite dans la conclusion, où aucune référence nouvelle ne doit apparaître ; de plus, la question finale au correcteur (« Et vous, qu'en pensez-vous ? ») est une faute : c'est l'élève qui doit répondre.
\end{itemize}

\textbf{2. Conclusion réécrite (exemple) :} « Nous avons d'abord vu que le travail peut apparaître comme une malédiction : pénibilité, contrainte, aliénation de l'ouvrier à la tâche. Nous avons ensuite montré que le travail est aussi ce par quoi l'homme transforme la nature, se réalise et conquiert sa dignité : l'agriculteur de l'Ouest qui cultive son champ n'y subit pas seulement une peine, il y construit son autonomie. La synthèse s'impose alors : le travail n'est pas une malédiction en soi ; ce sont ses conditions — exploitation, injustice, absence de reconnaissance — qui le rendent maudit. Répondre à notre problématique revient donc à distinguer le travail de ses formes dégradées : libéré de l'exploitation, le travail est un moyen d'émancipation. La question n'est plus de maudire le travail, mais de transformer les conditions dans lesquelles il s'exerce. »
\end{corrige}

\coursec{TP 4 : Choix et utilisation des citations}

\begin{corrige}[Exercice 7 — Classer et commenter des citations]
\textbf{1. Classement :}
\begin{center}
\begin{tabularx}{\linewidth}{|l|X|}\hline
\rowcolor{popblueL} \textbf{Notion} & \textbf{Numéro(s) de la citation} \\ \hline
Liberté & 2 (Rousseau) \\ \hline
Devoir & 1 (Kant) \\ \hline
État & 5 (attribué à Louis XIV) \\ \hline
Culture & 4 (Herriot) \\ \hline
Bonheur & 3 (Aristote) \\ \hline
Vérité & 6 (Bacon) \\ \hline
\end{tabularx}
\end{center}

\textbf{2. Commentaire (exemple avec la citation 2, Rousseau) :} Cette phrase célèbre de Rousseau oppose la liberté naturelle de l'homme (« né libre ») à sa situation sociale effective (« dans les fers »). Elle signifie que l'homme, libre par nature, se trouve partout enchaîné par les institutions, les conventions et les rapports de domination. L'enjeu est de distinguer la liberté de droit (l'essence humaine) et la liberté de fait (les conditions réelles d'existence). Dans une dissertation sur la liberté ou sur l'État, elle permet de poser le problème politique central : la société peut-elle légitimer les chaînes qu'elle impose, et comment rendre la liberté effective ? Elle sert aussi de problématique possible : l'État est-il ce qui enchaîne ou ce qui peut libérer ?

\textbf{Points de vigilance :} le commentaire doit comporter le sens, l'enjeu et l'usage possible dans un devoir ; pour la citation 5, noter la prudence méthodologique : « attribué à » signale une attribution incertaine, qu'il faut mentionner dans la copie.
\end{corrige}

\coursec{Application à la vie courante et entraînement à l'examen}

\begin{corrige}[Exercice 9 — Sujet type Bac : introduction et plan détaillé]
\textbf{1. Définitions :} la \cle{technique} est l'ensemble des procédés, outils et savoir-faire par lesquels l'homme transforme la nature et organise sa vie ; la \cle{liberté} est le pouvoir d'agir selon sa propre volonté, sans contrainte extérieure ni dépendance.

\textbf{2. Introduction (exemple) :} « À Douala comme à Bamenda, le conducteur de moto-taxi gagne sa vie grâce à sa machine, et le commerçant règle ses transactions par téléphone mobile : la technique est devenue la compagne quotidienne de nos existences. Le sujet interroge précisément ce rapport entre l'homme et ses outils. Par \emph{technique}, on entend l'ensemble des procédés et des instruments fabriqués pour transformer la nature ; par \emph{liberté}, le pouvoir d'agir selon sa volonté propre. Or l'outil qui accroît notre puissance ne crée-t-il pas en même temps de nouvelles dépendances ? La technique libère-t-elle l'homme, ou l'asservit-elle à ses propres créations ? Nous montrerons d'abord que la technique libère l'homme des contraintes naturelles ; nous verrons ensuite qu'elle peut devenir un instrument d'asservissement ; nous dégagerons enfin les conditions d'un usage libre et maîtrisé de la technique. »

\textbf{3. Plan détaillé :}
\begin{itemize}
\item \textbf{Thèse — la technique libère.} Idée directrice : en augmentant la puissance humaine, la technique affranchit l'homme de la nature et de la misère. Arguments : (a) elle réduit la peine et le temps de travail ; (b) elle vainc la distance et la maladie. Références : Descartes, qui fait de la technique un moyen de nous rendre « maîtres et possesseurs de la nature » (\emph{Discours de la méthode}) ; exemple : le forage qui donne l'eau potable dans un village du Nord.
\item \textbf{Antithèse — la technique asservit.} Idée directrice : l'outil retourne contre l'homme la dépendance qu'il devait vaincre. Arguments : (a) dépendance aux machines (panne de réseau, addiction au téléphone) ; (b) la technique peut servir la domination (surveillance, exploitation). Références : Rousseau, qui montre que le progrès technique corrompt et enchaîne (\emph{Discours sur l'origine de l'inégalité}) ; exemple : le jeune qui ne peut plus travailler sans son smartphone.
\item \textbf{Synthèse — un usage libre de la technique.} Idée directrice : la technique est un moyen ; tout dépend de la fin et de la maîtrise que l'homme en garde. Arguments : (a) la valeur de l'outil dépend de l'usage ; (b) l'éducation et la volonté politique peuvent orienter le progrès. Référence : la sagesse qui veut que « c'est la main qui commande à la machette, et non l'inverse » (proverbe commenté).
\end{itemize}

\textbf{4. Transition thèse $\rightarrow$ antithèse :} « La technique apparaît ainsi comme la grande libératrice de l'humanité. \emph{Mais} cette puissance nouvelle est-elle sans contrepartie ? Celui qui dépend de sa machine pour vivre, pour se déplacer, pour communiquer, est-il encore libre lorsque la machine lui échappe ? C'est la face obscure du progrès qu'il faut maintenant examiner. »

\textbf{Barème indicatif :} définitions /2 ; introduction complète en cinq étapes /6 ; plan détaillé cohérent (idées, arguments, références) /9 ; transition /3.
\textbf{Points de vigilance :} accroche concrète et courte ; problématique sous forme interrogative ; références exactes (Descartes, \emph{Discours de la méthode}) ; aucune partie ne doit oublier de répondre à la question « libère ou asservit ».
\end{corrige}

\begin{corrige}[Exercice 11 — Dissertation complète : « L'État peut-il imposer le civisme ? »]
\textbf{Corrigé-type (éléments attendus) :}

\textbf{Introduction.} Accroche concrète : à Yaoundé, les opérations « ville propre » et les contraventions pour jet d'ordures montrent l'État au quotidien dans la rue. Présentation du sujet : le civisme, ensemble des comportements de respect de la cité, relève-t-il de la contrainte publique ? Définitions : l'\cle{État} est l'institution politique détenant le pouvoir légitime de contrainte sur un territoire ; le \cle{civisme} est l'attitude du citoyen qui respecte les règles et le bien commun. Problématique : l'État peut-il produire le civisme par la loi et la sanction, ou le civisme véritable suppose-t-il une adhésion intérieure que nulle contrainte ne fabrique ? Annonce du plan : la contrainte comme condition du civisme (thèse) ; ses limites (antithèse) ; l'articulation de la loi et de l'éducation (synthèse).

\textbf{Thèse.} Sans loi, la vie commune est menacée : Hobbes montre que, hors de l'État, l'homme est exposé à la violence de tous, et que seul un pouvoir commun garantit la paix (\emph{Léviathan}). Argument 1 : la sanction rend les règles effectives (le code de la route sanctionné fait baisser les accidents). Argument 2 : l'impôt obligatoire finance les services dont tous profitent. Exemple camerounais : le paiement des impôts et taxes permet l'entretien des routes et des écoles.

\textbf{Transition.} « L'État apparaît ainsi comme le garant du civisme. \emph{Cependant}, l'obéissance par peur de la sanction est-elle déjà une vertu civique ? »

\textbf{Antithèse.} La contrainte produit des comportements, non des convictions. Argument 1 : un civisme de façade disparaît dès que le contrôle cesse (le chauffeur qui ralentit devant le gendarme et accélère après). Argument 2 : pour Rousseau, la légitimité vient de la volonté générale, c'est-à-dire de citoyens qui adhèrent à la loi comme à leur propre volonté (\emph{Du contrat social}) ; Kant ajoute qu'une action n'a de valeur morale que si elle est faite par devoir, non par crainte (\emph{Fondements de la métaphysique des mœurs}). La sagesse africaine le confirme : « on ne force pas un âne à boire » — commenté : la contrainte peut mener à l'eau, non donner la soif du bien commun.

\textbf{Synthèse.} L'État ne peut pas imposer le civisme comme on impose une taxe, mais il peut en créer les conditions : par la loi qui cadre, par l'exemple des dirigeants, et surtout par l'éducation civique qui transforme l'obéissance en adhésion. Le civisme véritable naît de l'articulation entre la contrainte extérieure et la conviction intérieure.

\textbf{Conclusion.} Réponse claire : non, l'État ne peut pas imposer le civisme par la seule force ; oui, il peut et doit le rendre possible par la loi juste, l'exemple et l'éducation. Bilan des trois parties rappelé en deux phrases.

\textbf{Auto-évaluation (exemple de bilan) :} « Point fort : ma problématique est claire et mes transitions explicites. Point faible : mes références sont peu développées, le commentaire de Hobbes est trop bref. Axe de progrès : approfondir le commentaire des citations et soigner la relecture orthographique. »

\textbf{Barème indicatif :} appliquer la grille de l'énoncé : présentation /4 ; pertinence /4 ; introduction /3 ; argumentation et transitions /4 ; références /3 ; conclusion /2.
\textbf{Points de vigilance :} (1) les deux références philosophiques doivent être exactes et commentées ; (2) l'exemple camerounais doit être précis et relié à l'argument ; (3) le proverbe doit être commenté, pas simplement cité ; (4) toute pensée d'auteur incertaine est reformulée sans guillemets ; (5) la conclusion répond par « oui/non, sous condition » et n'introduit rien de nouveau ; (6) l'auto-évaluation doit être honnête et chiffrée d'après la grille.
\end{corrige}

\newpage

\coursec{Fiche bilan}

\begin{popfigure}
\begin{tikzpicture}[mindmap, grow cyclic, every node/.style={concept, font=\small},
  root concept/.append style={concept color=popink, font=\bfseries},
  level 1/.append style={level distance=3.4cm, sibling angle=60},
  level 2/.append style={level distance=2.2cm, font=\scriptsize}]
\node[root concept]{Dissertation philosophique}
  child[concept color=popblue]{ node {Introduction}
    child { node {Accroche} }
    child { node {Problématique} }
    child { node {Plan annoncé} }
  }
  child[concept color=poporange]{ node {Thèse}
    child { node {Arguments} }
    child { node {Exemples} }
  }
  child[concept color=poppurple]{ node {Antithèse}
    child { node {Objection} }
    child { node {Critique} }
  }
  child[concept color=popgreen]{ node {Synthèse}
    child { node {Dépassement} }
    child { node {Réponse nuancée} }
  }
  child[concept color=poppink]{ node {Conclusion}
    child { node {Bilan} }
    child { node {Ouverture} }
  }
  child[concept color=popgold]{ node {Citations}
    child { node {Pertinentes} }
    child { node {Commentées} }
  };
\end{tikzpicture}
\legende{Carte mentale de la leçon : les six moments de la dissertation philosophique (introduction, thèse, antithèse, synthèse, conclusion, citations), avec leurs exigences essentielles.}
\end{popfigure}

\begin{bilan}
\textbf{Définitions clés}
\begin{itemize}
  \item \cle{Dissertation philosophique} : exercice écrit qui traite un problème de manière argumentée, ordonnée et personnelle, en mobilisant des notions, des auteurs et des exemples.
  \item \cle{Problématique} : question centrale, issue de l'analyse du sujet, qui oriente toute la réflexion.
  \item \cle{Thèse} : première réponse argumentée au problème posé.
  \item \cle{Antithèse} : objection ou position contraire qui met la thèse à l'épreuve.
  \item \cle{Synthèse} : dépassement du conflit thèse/antithèse ; réponse nuancée et personnelle.
  \item \cle{Citation} : parole exacte d'un auteur, rapportée entre guillemets, introduite et commentée.
\end{itemize}

\textbf{Structure à connaître par c\oe ur}
\begin{center}
\begin{tabularx}{\linewidth}{|l|X|}
\hline
\rowcolor{popblueL} \textbf{Partie} & \textbf{Contenu exigé} \\
\hline
Introduction & Accroche (amorce) ; définition des termes du sujet ; problématique ; annonce du plan. \\
\hline
Développement I & Thèse : 2 à 3 arguments, chacun étayé d'un exemple et, si possible, d'un auteur. \\
\hline
Développement II & Antithèse : objection sérieuse, arguments et exemples contraires. \\
\hline
Développement III & Synthèse : discussion critique, dépassement, position personnelle justifiée. \\
\hline
Conclusion & Bilan du parcours ; réponse claire à la problématique ; ouverture éventuelle. \\
\hline
\end{tabularx}
\end{center}

\textbf{Méthodes express}
\begin{itemize}
  \item Analyser le sujet : souligner les mots-clés, les définir, repérer la présupposition (ce que le sujet prend pour acquis).
  \item Formuler la problématique sous forme d'une vraie question (jamais une simple reprise du sujet).
  \item Un paragraphe = une idée = un argument + un exemple + une analyse.
  \item Citation : l'introduire (qui parle ?), la citer exactement, la commenter (que prouve-t-elle ?).
  \item Relire : transitions entre les parties, cohérence, orthographe.
\end{itemize}
\end{bilan}

\begin{aretenir}[Checklist avant le Bac]
Avant de rendre ma copie, je vérifie que :
\begin{itemize}
  \item[$\square$] j'ai bien analysé et défini les termes du sujet ;
  \item[$\square$] mon introduction contient une accroche, une problématique claire et un plan annoncé ;
  \item[$\square$] ma problématique est une véritable question philosophique ;
  \item[$\square$] mon développement suit le schéma thèse -- antithèse -- synthèse ;
  \item[$\square$] chaque partie contient au moins deux arguments étayés d'exemples concrets ;
  \item[$\square$] j'ai mobilisé au moins deux auteurs ou notions du programme ;
  \item[$\square$] mes citations sont exactes, entre guillemets, introduites et commentées ;
  \item[$\square$] mes exemples sont précis et reliés à l'argument (vie courante, contexte camerounais possible) ;
  \item[$\square$] il y a des transitions visibles entre les parties ;
  \item[$\square$] ma synthèse dépasse réellement l'opposition et exprime une position justifiée ;
  \item[$\square$] ma conclusion répond explicitement à la problématique ;
  \item[$\square$] j'ai relu l'orthographe, la grammaire et la présentation (paragraphes, alinéas).
\end{itemize}
\end{aretenir}

\findecours

\end{document}
